% BEGIN STANDALONE DEPENDENCIES
\begin{filecontents*}{iopart.cls}
%% 
%% This is file `iopart.cls'
%% 
%% This file is distributed in the hope that it will be useful, 
%% but WITHOUT ANY WARRANTY; without even the implied warranty of 
%% MERCHANTABILITY or FITNESS FOR A PARTICULAR PURPOSE. 
%% 
%% Licensed under the LPPL: http://www.latex-project.org/lppl.txt
%% Current Maintainer: IOP Publishing Ltd
%% 
%% \CharacterTable
%%  {Upper-case    \A\B\C\D\E\F\G\H\I\J\K\L\M\N\O\P\Q\R\S\T\U\V\W\X\Y\Z
%%   Lower-case    \a\b\c\d\e\f\g\h\i\j\k\l\m\n\o\p\q\r\s\t\u\v\w\x\y\z
%%   Digits        \0\1\2\3\4\5\6\7\8\9
%%   Exclamation   \!     Double quote  \"     Hash (number) \#
%%   Dollar        \$     Percent       \%     Ampersand     \&
%%   Acute accent  \'     Left paren    \(     Right paren   \)
%%   Asterisk      \*     Plus          \+     Comma         \,
%%   Minus         \-     Point         \.     Solidus       \/
%%   Colon         \:     Semicolon     \;     Less than     \<
%%   Equals        \=     Greater than  \>     Question mark \?
%%   Commercial at \@     Left bracket  \[     Backslash     \\
%%   Right bracket \]     Circumflex    \^     Underscore    \_
%%   Grave accent  \`     Left brace    \{     Vertical bar  \|
%%   Right brace   \}     Tilde         \~}
%
\NeedsTeXFormat{LaTeX2e}
\ProvidesClass{iopart}[1996/06/10 v0.0 IOP Journals LaTeX article class]
\newcommand\@ptsize{0}
\newif\if@restonecol
\newif\if@titlepage
\newif\ifiopams
\@titlepagefalse
\DeclareOption{a4paper}
   {\setlength\paperheight {297mm}%
    \setlength\paperwidth  {210mm}}
\DeclareOption{letterpaper}
   {\setlength\paperheight {11in}%
    \setlength\paperwidth  {8.5in}}
\DeclareOption{landscape}
   {\setlength\@tempdima   {\paperheight}%
    \setlength\paperheight {\paperwidth}%
    \setlength\paperwidth  {\@tempdima}}
\DeclareOption{10pt}{\renewcommand\@ptsize{0}}
\DeclareOption{11pt}{\renewcommand\@ptsize{2}} % No 11pt version
\DeclareOption{12pt}{\renewcommand\@ptsize{2}}
\DeclareOption{draft}{\setlength\overfullrule{5pt}}
\DeclareOption{final}{\setlength\overfullrule{0pt}}
\DeclareOption{titlepage}{\@titlepagetrue}
\DeclareOption{notitlepage}{\@titlepagefalse}
\ExecuteOptions{letterpaper,final}
\ProcessOptions
\DeclareMathAlphabet{\bi}{OML}{cmm}{b}{it}
\DeclareMathAlphabet{\bcal}{OMS}{cmsy}{b}{n}
\input{iopart1\@ptsize.clo}
\setlength\lineskip{1\p@}
\setlength\normallineskip{1\p@}
\renewcommand\baselinestretch{}
\setlength\parskip{0\p@ \@plus \p@}
\@lowpenalty   51
\@medpenalty  151
\@highpenalty 301
\setlength\parindent{2em}
\setcounter{topnumber}{8}
\renewcommand\topfraction{1}
\setcounter{bottomnumber}{3}
\renewcommand\bottomfraction{.99}
\setcounter{totalnumber}{8}
\renewcommand\textfraction{0.01}
\renewcommand\floatpagefraction{.8}
\setcounter{dbltopnumber}{6}
\renewcommand\dbltopfraction{1}
\renewcommand\dblfloatpagefraction{.8}
%
\pretolerance=5000
\tolerance=8000
%
% Headings for all pages apart from first
%
\def\ps@headings{\let\@oddfoot\@empty
      \let\@evenfoot\@empty
      \def\@evenhead{\thepage\hfil\itshape\rightmark}%
      \def\@oddhead{{\itshape\leftmark}\hfil\thepage}%
      \let\@mkboth\markboth
      \let\sectionmark\@gobble
      \let\subsectionmark\@gobble}
%
% Headings for first page
%
\def\ps@myheadings{\let\@oddfoot\@empty\let\@evenfoot\@empty
    \let\@oddhead\@empty\let\@evenhead\@empty
    \let\@mkboth\@gobbletwo
    \let\sectionmark\@gobble
    \let\subsectionmark\@gobble}
%
% \maketitle just ends page
%
\newcommand\maketitle{\newpage}
%
% Article titles
%
% Usage: \title[Short title]{Full title}
% [Short title] is optional; use where title is too long
% or contains footnotes, 50 characters maximum 
%
\renewcommand{\title}{\@ifnextchar[{\@stitle}{\@ftitle}}
\def\@stitle[#1]#2{\markboth{#1}{#1}%
    \thispagestyle{myheadings}%
    \vspace*{3pc}{\exhyphenpenalty=10000\hyphenpenalty=10000 
    \Large\raggedright\noindent
    \bf#2\par}}
\def\@ftitle#1{\markboth{#1}{#1}%
    \thispagestyle{myheadings}%
    \vspace*{3pc}{\exhyphenpenalty=10000\hyphenpenalty=10000 
    \Large\raggedright\noindent
    \bf#1\par}}
%
% Can use \paper instead of \title
%
\let\paper=\title
%
% Generic title command for articles other than papers
%
% Usage: \article[Short title]{Article Type}{Full title}
% [Short title] is optional; use where title is too long
% or contains footnotes, 50 characters maximum 
%
\newcommand{\article}{\@ifnextchar[{\@sarticle}{\@farticle}}
\def\@sarticle[#1]#2#3{\markboth{#1}{#1}%
    \thispagestyle{myheadings}%
     \vspace*{.5pc}%
    {\parindent=\mathindent \bf #2\par}%
     \vspace*{1.5pc}%
    {\exhyphenpenalty=10000\hyphenpenalty=10000
     \Large\raggedright\noindent
     \bf#3\par}}%
\def\@farticle#1#2{\markboth{#2}{#2}%
    \thispagestyle{myheadings}%
     \vspace*{.5pc}%
    {\parindent=\mathindent \bf #1\par}%
     \vspace*{1.5pc}%
    {\exhyphenpenalty=10000\hyphenpenalty=10000
     \Large\raggedright\noindent
     \bf#2\par}}%
%
% Letters to the Editor
%
% Usage \letter{Full title}
% No short title is required for Letters
%
\def\letter#1{\article[Letter to the Editor]{Letter to the Editor}{#1}}
%
% Fast Track Communications (added by sxb 9 March 2011)
%
% Usage \ftc{Full title} - there's no short title
\def\ftc#1{\article[Fast Track Communication]{Fast Track Communication}{#1}}
%
%
% Review articles
%
% Usage: \review[Short title]{Full title}
% [Short title] is optional; use where title is too long
% or contains footnotes, 50 characters maximum 
%
\def\review{\@ifnextchar[{\@sreview}{\@freview}}
\def\@sreview[#1]#2{\@sarticle[#1]{Review Article}{#2}}
\def\@freview#1{\@farticle{Review Article}{#1}}
%
% Topical Review
%
% Usage: \topical[Short title]{Full title}
% [Short title] is optional; use where title is too long
% or contains footnotes, 50 characters maximum 
%
\def\topical{\@ifnextchar[{\@stopical}{\@ftopical}}
\def\@stopical[#1]#2{\@sarticle[#1]{Topical Review}{#2}}
\def\@ftopical#1{\@farticle{Topical Review}{#1}}
%
% Comments
%
% Usage: \comment[Short title]{Full title}
% [Short title] is optional; use where title is too long
% or contains footnotes, 50 characters maximum 
%
\def\comment{\@ifnextchar[{\@scomment}{\@fcomment}}
\def\@scomment[#1]#2{\@sarticle[#1]{Comment}{#2}}
\def\@fcomment#1{\@farticle{Comment}{#1}}
%
% Rapid Communications
%
% Usage: \rapid[Short title]{Full title}
% [Short title] is optional; use where title is too long
% or contains footnotes, 50 characters maximum 
%
\def\rapid{\@ifnextchar[{\@srapid}{\@frapid}}
\def\@srapid[#1]#2{\@sarticle[#1]{Rapid Communication}{#2}}
\def\@frapid#1{\@farticle{Rapid Communication}{#1}}
%
% Notes
%
% Usage: \note[Short title]{Full title}
% [Short title] is optional; use where title is too long
% or contains footnotes, 50 characters maximum 
%
\def\note{\@ifnextchar[{\@snote}{\@fnote}}
\def\@snote[#1]#2{\@sarticle[#1]{Note}{#2}}
\def\@fnote#1{\@farticle{Note}{#1}}
%
% Preliminary Communications
%
% Usage: \prelim[Short title]{Full title}
% [Short title] is optional; use where title is too long
% or contains footnotes, 50 characters maximum 
%
\def\prelim{\@ifnextchar[{\@sprelim}{\@fprelim}}
\def\@sprelim[#1]#2{\@sarticle[#1]{Preliminary Communication}{#2}}
\def\@fprelim#1{\@farticle{Preliminary Communication}{#1}}
%
% List of authors
%
% Usage \author[Short form]{List of all authors}
% The short form excludes footnote symbols linking authors to addresses 
% and is used for running heads in printed version (but not on preprints)
%
\renewcommand{\author}{\@ifnextchar[{\@sauthor}{\@fauthor}}
\def\@sauthor[#1]#2{\markright{#1}    % for production only
   \vspace*{1.5pc}%
   \begin{indented}%
   \item[]\normalsize\bf\raggedright#2
   \end{indented}%
   \smallskip}
\def\@fauthor#1{%\markright{#1}         for production only
   \vspace*{1.5pc}%
   \begin{indented}%
   \item[]\normalsize\bf\raggedright#1
   \end{indented}%
   \smallskip}
%
% Affiliation (authors address)
%
% Usage: \address{Address of first author}
%        \address{Address of second author}
% Use once for each address, use symbols \dag \ddag \S \P $\|$
% to connect authors with addresses
%
\newcommand{\address}[1]{\begin{indented}
   \item[]\rm\raggedright #1
   \end{indented}}
%
% American Mathematical Society Classification Numbers
% Usage: \ams{57.XX, 58.XX}
%
\def\ams#1{\vspace{10pt}
     \begin{indented}
     \item[]\rm AMS classification scheme numbers: #1\par
     \end{indented}}
%
% A single Physics & Astronomy Classification Number
% Usage \pacno{31.10}
%
\def\pacno#1{\vspace{10pt}
     \begin{indented}
     \item[]\rm PACS number: #1\par
     \end{indented}}
%
% Physics & Astronomy Classification Numbers (more than one)
% Usage \pacs{31.10, 31.20T}
%
\def\pacs#1{\vspace{10pt}
     \begin{indented}
     \item[]\rm PACS numbers: #1\par
     \end{indented}}
%
% Submission details. If \jl command used journals name printed
% otherwise Institute of Physics Publishing 
%
\def\submitted{\vspace{28pt plus 10pt minus 18pt}
     \noindent{\small\rm Submitted to: {\it \journal}\par}}
%
\def\submitto#1{\vspace{28pt plus 10pt minus 18pt}
     \noindent{\small\rm Submitted to: {\it #1}\par}}
%
% For articles (other than Letters) not divided into sections
% Usage \nosections Start of text
%
\def\nosections{\vspace{30\p@ plus12\p@ minus12\p@}
    \noindent\ignorespaces}
%
% Acknowledgments (no heading if letter)
% Usage \ack for Acknowledgments, \ackn for Acknowledgement
%
\def\ack{\ifletter\bigskip\noindent\ignorespaces\else
    \section*{Acknowledgments}\fi}
\def\ackn{\ifletter\bigskip\noindent\ignorespaces\else
    \section*{Acknowledgment}\fi}
%
% Footnotes: symbols selected in order \dag (1), \ddag (2), \S (3), 
% $\|$ (4), $\P$ (5), $^+$ (6), $^*$ (7), \sharp (8), \dagger\dagger (9)
% unless optional argument of [<num>] use to specify required symbol, 
% 1=\dag, 2=\ddag, etc
% Usage: \footnote{Text of footnote}
%        \footnote[3]{Text of footnote}
%
\def\footnoterule{}%
\setcounter{footnote}{1}
\long\def\@makefntext#1{\parindent 1em\noindent 
 \makebox[1em][l]{\footnotesize\rm$\m@th{\fnsymbol{footnote}}$}%
 \footnotesize\rm #1}
\def\@makefnmark{\hbox{${\fnsymbol{footnote}}\m@th$}}
\def\@thefnmark{\fnsymbol{footnote}}
\def\footnote{\@ifnextchar[{\@xfootnote}{\stepcounter{\@mpfn}%
       \begingroup\let\protect\noexpand
       \xdef\@thefnmark{\thempfn}\endgroup
     \@footnotemark\@footnotetext}}
\def\@xfootnote[#1]{\setcounter{footnote}{#1}%
   \addtocounter{footnote}{-1}\footnote}
\def\@fnsymbol#1{\ifcase#1\or \dagger\or \ddagger\or \S\or
   \|\or \P\or ^{+}\or ^{\tsty *}\or \sharp
   \or \dagger\dagger \else\@ctrerr\fi\relax}
%
% IOP Journals
%
\newcounter{jnl}
\newcommand{\jl}[1]{\setcounter{jnl}{#1}}
\def\journal{\ifnum\thejnl=0 Institute of Physics Publishing\fi
        \ifnum\thejnl=1 J. Phys.\ A: Math.\ Gen.\ \fi
        \ifnum\thejnl=2 J. Phys.\ B: At.\ Mol.\ Opt.\ Phys.\ \fi
        \ifnum\thejnl=3 J. Phys.:\ Condens. Matter\ \fi
        \ifnum\thejnl=4 J. Phys.\ G: Nucl.\ Part.\ Phys.\ \fi
        \ifnum\thejnl=5 Inverse Problems\ \fi
        \ifnum\thejnl=6 Class. Quantum Grav.\ \fi
        \ifnum\thejnl=7 Network: Comput.\ Neural Syst.\ \fi
        \ifnum\thejnl=8 Nonlinearity\ \fi
        \ifnum\thejnl=9 J. Opt. B: Quantum Semiclass. Opt.\ \fi
        \ifnum\thejnl=10 Waves Random Media\ \fi
        \ifnum\thejnl=11 J. Opt. A: Pure Appl. Opt.\ \fi
        \ifnum\thejnl=12 Phys. Med. Biol.\ \fi
        \ifnum\thejnl=13 Modelling Simul.\ Mater.\ Sci.\ Eng.\ \fi
        \ifnum\thejnl=14 Plasma Phys. Control. Fusion\ \fi
        \ifnum\thejnl=15 Physiol. Meas.\ \fi
        \ifnum\thejnl=16 Combust. Theory Modelling\ \fi
        \ifnum\thejnl=17 High Perform.\ Polym.\ \fi
        \ifnum\thejnl=18 Public Understand. Sci.\ \fi
        \ifnum\thejnl=19 Rep.\ Prog.\ Phys.\ \fi
        \ifnum\thejnl=20 J.\ Phys.\ D: Appl.\ Phys.\ \fi
        \ifnum\thejnl=21 Supercond.\ Sci.\ Technol.\ \fi
        \ifnum\thejnl=22 Semicond.\ Sci.\ Technol.\ \fi
        \ifnum\thejnl=23 Nanotechnology\ \fi
        \ifnum\thejnl=24 Measur.\ Sci.\ Technol.\ \fi
        \ifnum\thejnl=25 Plasma.\ Sources\ Sci.\ Technol.\ \fi
        \ifnum\thejnl=26 Smart\ Mater.\ Struct.\ \fi
        \ifnum\thejnl=27 J.\ Micromech.\ Microeng.\ \fi
        \ifnum\thejnl=28 Distrib.\ Syst.\ Engng\ \fi
        \ifnum\thejnl=29 Bioimaging\ \fi
        \ifnum\thejnl=30 J.\ Radiol. Prot.\ \fi
        \ifnum\thejnl=31 Europ. J. Phys.\ \fi
        \ifnum\thejnl=32 J. Opt. A: Pure Appl. Opt.\ \fi
        \ifnum\thejnl=33 New. J. Phys.\ \fi}
%
% E-mail addresses (to provide links from headers)
%
\def\eads#1{\vspace*{5pt}\address{E-mail: #1}}
\def\ead#1{\vspace*{5pt}\address{E-mail: \mailto{#1}}}
\def\mailto#1{{\tt #1}}
%
% Switches
%
\newif\ifletter
%
\setcounter{secnumdepth}{3}
\newcounter {section}
\newcounter {subsection}[section]
\newcounter {subsubsection}[subsection]
\newcounter {paragraph}[subsubsection]
\newcounter {subparagraph}[paragraph]
\renewcommand\thesection       {\arabic{section}}
\renewcommand\thesubsection    {\thesection.\arabic{subsection}}
\renewcommand\thesubsubsection {\thesubsection.\arabic{subsubsection}}
\renewcommand\theparagraph     {\thesubsubsection.\arabic{paragraph}}
\renewcommand\thesubparagraph  {\theparagraph.\arabic{subparagraph}}
\def\@chapapp{Section}

\newcommand\section{\@startsection {section}{1}{\z@}%
                   {-3.5ex \@plus -1ex \@minus -.2ex}%
                   {2.3ex \@plus.2ex}%
                   {\reset@font\normalsize\bfseries\raggedright}}
\newcommand\subsection{\@startsection{subsection}{2}{\z@}%
                   {-3.25ex\@plus -1ex \@minus -.2ex}%
                   {1.5ex \@plus .2ex}%
                   {\reset@font\normalsize\itshape\raggedright}}
\newcommand\subsubsection{\@startsection{subsubsection}{3}{\z@}%
                                     {-3.25ex\@plus -1ex \@minus -.2ex}%
                                     {-1em \@plus .2em}%
                                     {\reset@font\normalsize\itshape}}
\newcommand\paragraph{\@startsection{paragraph}{4}{\z@}%
                                    {3.25ex \@plus1ex \@minus.2ex}%
                                    {-1em}%
                                    {\reset@font\normalsize\itshape}}
\newcommand\subparagraph{\@startsection{subparagraph}{5}{\parindent}%
                                       {3.25ex \@plus1ex \@minus .2ex}%
                                       {-1em}%
                                      {\reset@font\normalsize\itshape}}
\def\@sect#1#2#3#4#5#6[#7]#8{\ifnum #2>\c@secnumdepth
     \let\@svsec\@empty\else
     \refstepcounter{#1}\edef\@svsec{\csname the#1\endcsname. }\fi
     \@tempskipa #5\relax
      \ifdim \@tempskipa>\z@
        \begingroup #6\relax
          \noindent{\hskip #3\relax\@svsec}{\interlinepenalty \@M #8\par}%
        \endgroup
       \csname #1mark\endcsname{#7}\addcontentsline
         {toc}{#1}{\ifnum #2>\c@secnumdepth \else
                      \protect\numberline{\csname the#1\endcsname}\fi
                    #7}\else
        \def\@svsechd{#6\hskip #3\relax  %% \relax added 2 May 90
                   \@svsec #8\csname #1mark\endcsname
                      {#7}\addcontentsline
                           {toc}{#1}{\ifnum #2>\c@secnumdepth \else
                             \protect\numberline{\csname the#1\endcsname}\fi
                       #7}}\fi
     \@xsect{#5}}
%
\def\@ssect#1#2#3#4#5{\@tempskipa #3\relax
   \ifdim \@tempskipa>\z@
     \begingroup #4\noindent{\hskip #1}{\interlinepenalty \@M #5\par}\endgroup
   \else \def\@svsechd{#4\hskip #1\relax #5}\fi
    \@xsect{#3}}

\setlength\leftmargini{2em}
\setlength\leftmarginii{2em}
\setlength\leftmarginiii{1.8em}
\setlength\leftmarginiv{1.6em}
\setlength\leftmarginv{1em}
\setlength\leftmarginvi{1em}
\setlength\leftmargin{\leftmargini}
\setlength\labelsep{0.5em}
\setlength\labelwidth{\leftmargini}
\addtolength\labelwidth{-\labelsep}
\@beginparpenalty -\@lowpenalty
\@endparpenalty   -\@lowpenalty
\@itempenalty     -\@lowpenalty
\renewcommand\theenumi{\roman{enumi}}
\renewcommand\theenumii{\alph{enumii}}
\renewcommand\theenumiii{\arabic{enumiii}}
\renewcommand\theenumiv{\Alph{enumiv}}
\newcommand\labelenumi{(\theenumi)}
\newcommand\labelenumii{(\theenumii)}
\newcommand\labelenumiii{\theenumiii.}
\newcommand\labelenumiv{(\theenumiv)}
\renewcommand\p@enumii{(\theenumi)}
\renewcommand\p@enumiii{(\theenumi.\theenumii)}
\renewcommand\p@enumiv{(\theenumi.\theenumii.\theenumiii)}
\newcommand\labelitemi{$\m@th\bullet$}
\newcommand\labelitemii{\normalfont\bfseries --}
\newcommand\labelitemiii{$\m@th\ast$}
\newcommand\labelitemiv{$\m@th\cdot$}
\newenvironment{description}
               {\list{}{\labelwidth\z@ \itemindent-\leftmargin
                        \let\makelabel\descriptionlabel}}
               {\endlist}
\newcommand\descriptionlabel[1]{\hspace\labelsep
                                \normalfont\bfseries #1}
\newenvironment{abstract}{%
      \vspace{16pt plus3pt minus3pt}
      \begin{indented}
      \item[]{\bfseries \abstractname.}\quad\rm\ignorespaces} 
      {\end{indented}\if@titlepage\newpage\else\vspace{18\p@ plus18\p@}\fi}
\newenvironment{verse}
               {\let\\=\@centercr
                \list{}{\itemsep      \z@
                        \itemindent   -1.5em%
                        \listparindent\itemindent
                        \rightmargin  \leftmargin
                        \advance\leftmargin 1.5em}%
                \item[]}
               {\endlist}
\newenvironment{quotation}
               {\list{}{\listparindent 1.5em%
                        \itemindent    \listparindent
                        \rightmargin   \leftmargin
                        \parsep        \z@ \@plus\p@}%
                \item[]}
               {\endlist}
\newenvironment{quote}
               {\list{}{\rightmargin\leftmargin}%
                \item[]}
               {\endlist}
\newenvironment{titlepage}
    {%
        \@restonecolfalse\newpage
      \thispagestyle{empty}%
      \if@compatibility
        \setcounter{page}{0}
      \else
        \setcounter{page}{1}%
      \fi}%
    {\newpage\setcounter{page}{1}}
\def\appendix{\@ifnextchar*{\@appendixstar}{\@appendix}}
\def\@appendix{\eqnobysec\@appendixstar}
\def\@appendixstar{\@@par
 \ifnumbysec                         %  Added 30/4/94 to get Table A1,
 \@addtoreset{table}{section}        %  Table B1 etc if numbering by
 \@addtoreset{figure}{section}\fi    %  section
 \setcounter{section}{0}
 \setcounter{subsection}{0}
 \setcounter{subsubsection}{0}
 \setcounter{equation}{0}
 \setcounter{figure}{0}
 \setcounter{table}{0}
 \def\thesection{Appendix \Alph{section}}   
 \def\theequation{\ifnumbysec
      \Alph{section}.\arabic{equation}\else
      \Alph{section}\arabic{equation}\fi}  % Comment A\arabic{equation} maybe
 \def\thetable{\ifnumbysec                 % better? 15/4/95
      \Alph{section}\arabic{table}\else
      A\arabic{table}\fi}
 \def\thefigure{\ifnumbysec
      \Alph{section}\arabic{figure}\else
      A\arabic{figure}\fi}}
\def\noappendix{\setcounter{figure}{0}
     \setcounter{table}{0}
     \def\thetable{\arabic{table}}
     \def\thefigure{\arabic{figure}}}
\setlength\arraycolsep{5\p@}
\setlength\tabcolsep{6\p@}
\setlength\arrayrulewidth{.4\p@}
\setlength\doublerulesep{2\p@}
\setlength\tabbingsep{\labelsep}
\skip\@mpfootins = \skip\footins
\setlength\fboxsep{3\p@}
\setlength\fboxrule{.4\p@}
\renewcommand\theequation{\arabic{equation}}
\newcounter{figure}
\renewcommand\thefigure{\@arabic\c@figure}
\def\fps@figure{tbp}
\def\ftype@figure{1}
\def\ext@figure{lof}
\def\fnum@figure{\figurename~\thefigure}
\newenvironment{figure}{\footnotesize\rm\@float{figure}}%
    {\end@float\normalsize\rm}
\newenvironment{figure*}{\footnotesize\rm\@dblfloat{figure}}{\end@dblfloat}
\newcounter{table}
\renewcommand\thetable{\@arabic\c@table}
\def\fps@table{tbp}
\def\ftype@table{2}
\def\ext@table{lot}
\def\fnum@table{\tablename~\thetable}
\newenvironment{table}{\footnotesize\rm\@float{table}}%
   {\end@float\normalsize\rm}
\newenvironment{table*}{\footnotesize\rm\@dblfloat{table}}%
   {\end@dblfloat\normalsize\rm}
\newlength\abovecaptionskip
\newlength\belowcaptionskip
\setlength\abovecaptionskip{10\p@}
\setlength\belowcaptionskip{0\p@}
%
% Added redefinition of \@caption so captions are not written to 
% aux file therefore less need to \protect fragile commands
%
\long\def\@caption#1[#2]#3{\par\begingroup
    \@parboxrestore
    \normalsize
    \@makecaption{\csname fnum@#1\endcsname}{\ignorespaces #3}\par
  \endgroup}
% 
\long\def\@makecaption#1#2{\vskip \abovecaptionskip 
 \begin{indented}
 \item[]{\bf #1.} #2
 \end{indented}\vskip\belowcaptionskip}
\let\@portraitcaption=\@makecaption

\DeclareOldFontCommand{\rm}{\normalfont\rmfamily}{\mathrm}
\DeclareOldFontCommand{\sf}{\normalfont\sffamily}{\mathsf}
\DeclareOldFontCommand{\tt}{\normalfont\ttfamily}{\mathtt}
\DeclareOldFontCommand{\bf}{\normalfont\bfseries}{\mathbf}
\DeclareOldFontCommand{\it}{\normalfont\itshape}{\mathit}
\DeclareOldFontCommand{\sl}{\normalfont\slshape}{\@nomath\sl}
\DeclareOldFontCommand{\sc}{\normalfont\scshape}{\@nomath\sc}
\ifiopams
\renewcommand{\cal}{\protect\pcal}
\else
\newcommand{\cal}{\protect\pcal}
\fi
\newcommand{\pcal}{\@fontswitch{\relax}{\mathcal}}
\ifiopams
\renewcommand{\mit}{\protect\pmit}
\else
\newcommand{\mit}{\protect\pmit}
\fi
\newcommand{\pmit}{\@fontswitch{\relax}{\mathnormal}}
\newcommand\@pnumwidth{1.55em}
\newcommand\@tocrmarg {2.55em}
\newcommand\@dotsep{4.5}
\setcounter{tocdepth}{3}
\newcommand\tableofcontents{%
    \section*{\contentsname
        \@mkboth{\uppercase{\contentsname}}{\uppercase{\contentsname}}}%
    \@starttoc{toc}%
    }
\newcommand\l@part[2]{%
  \ifnum \c@tocdepth >-2\relax
    \addpenalty{\@secpenalty}%
    \addvspace{2.25em \@plus\p@}%
    \begingroup
      \setlength\@tempdima{3em}%
      \parindent \z@ \rightskip \@pnumwidth
      \parfillskip -\@pnumwidth
      {\leavevmode
       \large \bfseries #1\hfil \hbox to\@pnumwidth{\hss #2}}\par
       \nobreak
       \if@compatibility
         \global\@nobreaktrue
         \everypar{\global\@nobreakfalse\everypar{}}
      \fi
    \endgroup
  \fi}
\newcommand\l@section[2]{%
  \ifnum \c@tocdepth >\z@
    \addpenalty{\@secpenalty}%
    \addvspace{1.0em \@plus\p@}%
    \setlength\@tempdima{1.5em}%
    \begingroup
      \parindent \z@ \rightskip \@pnumwidth
      \parfillskip -\@pnumwidth
      \leavevmode \bfseries
      \advance\leftskip\@tempdima
      \hskip -\leftskip
      #1\nobreak\hfil \nobreak\hbox to\@pnumwidth{\hss #2}\par
    \endgroup
  \fi}
\newcommand\l@subsection   {\@dottedtocline{2}{1.5em}{2.3em}}
\newcommand\l@subsubsection{\@dottedtocline{3}{3.8em}{3.2em}}
\newcommand\l@paragraph    {\@dottedtocline{4}{7.0em}{4.1em}}
\newcommand\l@subparagraph {\@dottedtocline{5}{10em}{5em}}
\newcommand\listoffigures{%
    \section*{\listfigurename
      \@mkboth{\uppercase{\listfigurename}}%
              {\uppercase{\listfigurename}}}%
    \@starttoc{lof}%
    }
\newcommand\l@figure{\@dottedtocline{1}{1.5em}{2.3em}}
\newcommand\listoftables{%
    \section*{\listtablename
      \@mkboth{\uppercase{\listtablename}}{\uppercase{\listtablename}}}%
    \@starttoc{lot}%
    }
\let\l@table\l@figure
\newenvironment{theindex}
               {\if@twocolumn
                  \@restonecolfalse
                \else
                  \@restonecoltrue
                \fi
                \columnseprule \z@
                \columnsep 35\p@
                \twocolumn[\section*{\indexname}]%
                \@mkboth{\uppercase{\indexname}}%
                        {\uppercase{\indexname}}%
                \thispagestyle{plain}\parindent\z@
                \parskip\z@ \@plus .3\p@\relax
                \let\item\@idxitem}
               {\if@restonecol\onecolumn\else\clearpage\fi}
\newcommand\@idxitem  {\par\hangindent 40\p@}
\newcommand\subitem   {\par\hangindent 40\p@ \hspace*{20\p@}}
\newcommand\subsubitem{\par\hangindent 40\p@ \hspace*{30\p@}}
\newcommand\indexspace{\par \vskip 10\p@ \@plus5\p@ \@minus3\p@\relax}
\newcommand\contentsname{Contents}
\newcommand\listfigurename{List of Figures}
\newcommand\listtablename{List of Tables}
\newcommand\refname{References}
\newcommand\indexname{Index}
\newcommand\figurename{Figure}
\newcommand\tablename{Table}
\newcommand\partname{Part}
\newcommand\appendixname{Appendix}
\newcommand\abstractname{Abstract}
\newcommand\today{\number\day\space\ifcase\month\or
  January\or February\or March\or April\or May\or June\or
  July\or August\or September\or October\or November\or December\fi
  \space\number\year}
\setlength\columnsep{10\p@}
\setlength\columnseprule{0\p@}

\newcommand{\Tables}{\clearpage\section*{Tables and table captions}
    \def\fps@table{hp}\noappendix}
\newcommand{\Figures}{\clearpage\section*{Figure captions}
    \def\fps@figure{hp}\noappendix}
%
\newcommand{\Figure}[1]{\begin{figure}
  \caption{#1}
  \end{figure}}
%
\newcommand{\Table}[1]{\begin{table}
  \caption{#1}
  \begin{indented}
  \lineup
  \item[]\begin{tabular}{@{}l*{15}{l}}}

\def\endTable{\end{tabular}\end{indented}\end{table}}
\let\endtab=\endTable
%
\newcommand{\fulltable}[1]{\begin{table}
  \caption{#1}
  \footnotesize
  \lineup
  \begin{tabular*}{\textwidth}{@{}l*{15}{@{\extracolsep{0pt plus 12pt}}l}}}
\def\endfulltable{\end{tabular*}\end{table}\normalsize}
%
%
\newcommand{\Bibliography}[1]{\section*{References}\par\numrefs{#1}}
\newcommand{\References}{\section*{References}\par\refs}

\def\thebibliography#1{\list
 {\hfil[\arabic{enumi}]}{\topsep=0\p@\parsep=0\p@
 \partopsep=0\p@\itemsep=0\p@
 \labelsep=5\p@\itemindent=-10\p@
 \settowidth\labelwidth{\footnotesize[#1]}%
 \leftmargin\labelwidth
 \advance\leftmargin\labelsep
 \advance\leftmargin -\itemindent
 \usecounter{enumi}}\footnotesize
 \def\newblock{\ }
 \sloppy\clubpenalty4000\widowpenalty4000
 \sfcode`\.=1000\relax}
\let\endthebibliography=\endlist
\def\numrefs#1{\begin{thebibliography}{#1}}
\def\endnumrefs{\end{thebibliography}}
\let\endbib=\endnumrefs
%
\def\thereferences{\list{}{\topsep=0\p@\parsep=0\p@
 \partopsep=0\p@\itemsep=0\p@\labelsep=0\p@\itemindent=-18\p@
\labelwidth=0\p@\leftmargin=18\p@
}\footnotesize\rm
\def\newblock{\ }
\sloppy\clubpenalty4000\widowpenalty4000
\sfcode`\.=1000\relax
}
%
\let\endthereferences=\endlist
\newlength{\indentedwidth}
\newdimen\mathindent
\indentedwidth=\mathindent
%
% Macro to used for references in the Harvard system
%
\newenvironment{harvard}{\list{}{\topsep=0\p@\parsep=0\p@
\partopsep=0\p@\itemsep=0\p@\labelsep=0\p@\itemindent=-18\p@
\labelwidth=0\p@\leftmargin=18\p@
}\footnotesize\rm
\def\newblock{\ }
\sloppy\clubpenalty4000\widowpenalty4000
\sfcode`\.=1000\relax}{\endlist}
%
\def\refs{\begin{harvard}}
\def\endrefs{\end{harvard}}
%
\newenvironment{indented}{\begin{indented}}{\end{indented}}
\newenvironment{varindent}[1]{\begin{varindent}{#1}}{\end{varindent}}
%
\def\indented{\list{}{\itemsep=0\p@\labelsep=0\p@\itemindent=0\p@
   \labelwidth=0\p@\leftmargin=\mathindent\topsep=0\p@\partopsep=0\p@
   \parsep=0\p@\listparindent=15\p@}\footnotesize\rm}

\let\endindented=\endlist 

\def\varindent#1{\setlength{\varind}{#1}%
   \list{}{\itemsep=0\p@\labelsep=0\p@\itemindent=0\p@
   \labelwidth=0\p@\leftmargin=\varind\topsep=0\p@\partopsep=0\p@
   \parsep=0\p@\listparindent=15\p@}\footnotesize\rm}

\let\endvarindent=\endlist 

\def\[{\relax\ifmmode\@badmath\else
 \begin{trivlist}
 \@beginparpenalty\predisplaypenalty
 \@endparpenalty\postdisplaypenalty
 \item[]\leavevmode
 \hbox to\linewidth\bgroup$ \displaystyle
 \hskip\mathindent\bgroup\fi}
\def\]{\relax\ifmmode \egroup $\hfil \egroup \end{trivlist}\else \@badmath \fi}
\def\equation{\@beginparpenalty\predisplaypenalty
 \@endparpenalty\postdisplaypenalty
\refstepcounter{equation}\trivlist \item[]\leavevmode
 \hbox to\linewidth\bgroup $ \displaystyle
\hskip\mathindent}
\def\endequation{$\hfil \displaywidth\linewidth\@eqnnum\egroup \endtrivlist}
%
\@namedef{equation*}{\[}
\@namedef{endequation*}{\]}
%
\def\eqnarray{\stepcounter{equation}\let\@currentlabel=\theequation
\global\@eqnswtrue
\global\@eqcnt\z@\tabskip\mathindent\let\\=\@eqncr
\abovedisplayskip\topsep\ifvmode\advance\abovedisplayskip\partopsep\fi
\belowdisplayskip\abovedisplayskip
\belowdisplayshortskip\abovedisplayskip
\abovedisplayshortskip\abovedisplayskip
$$\halign to
\linewidth\bgroup\@eqnsel$\displaystyle\tabskip\z@
 {##{}}$&\global\@eqcnt\@ne $\displaystyle{{}##{}}$\hfil
 &\global\@eqcnt\tw@ $\displaystyle{{}##}$\hfil
 \tabskip\@centering&\llap{##}\tabskip\z@\cr}
\def\endeqnarray{\@@eqncr\egroup
 \global\advance\c@equation\m@ne$$\global\@ignoretrue }
\mathindent = 6pc
%
\def\eqalign#1{\null\vcenter{\def\\{\cr}\openup\jot\m@th
  \ialign{\strut$\displaystyle{##}$\hfil&$\displaystyle{{}##}$\hfil
      \crcr#1\crcr}}\,}
%
\def\eqalignno#1{\displ@y \tabskip\z@skip
  \halign to\displaywidth{\hspace{5pc}$\@lign\displaystyle{##}$%
    \tabskip\z@skip
    &$\@lign\displaystyle{{}##}$\hfill\tabskip\@centering
    &\llap{$\@lign\hbox{\rm##}$}\tabskip\z@skip\crcr
    #1\crcr}}
%
\newif\ifnumbysec
\def\theequation{\ifnumbysec
      \arabic{section}.\arabic{equation}\else
      \arabic{equation}\fi}
\def\eqnobysec{\numbysectrue\@addtoreset{equation}{section}}

\newcounter{eqnval}

% 2012: if you have an eqn numbered by parts (eg eqn 6a, 6b) this allows
% you to refer to it by just the number (eg eqn 6) if the label is between
% the \multiparteqn and the \numparts
\def\multiparteqn{\addtocounter{equation}{1}
\eqnarray\nonumber
\endeqnarray

\vspace{-72pt}
}

\def\numparts{\addtocounter{equation}{1}%
     \setcounter{eqnval}{\value{equation}}%
     \setcounter{equation}{0}%
     \def\theequation{\ifnumbysec
     \arabic{section}.\arabic{eqnval}{\it\alph{equation}}%
     \else\arabic{eqnval}{\it\alph{equation}}\fi}}

\def\endnumparts{\def\theequation{\ifnumbysec
     \arabic{section}.\arabic{equation}\else
     \arabic{equation}\fi}%
     \setcounter{equation}{\value{eqnval}}}
%
\def\cases#1{%
     \left\{\,\vcenter{\def\\{\cr}\normalbaselines\openup1\jot\m@th%
     \ialign{\strut$\displaystyle{##}\hfil$&\tqs
     \rm##\hfil\crcr#1\crcr}}\right.}%
%
\newcommand{\e}{\mathrm{e}}
\newcommand{\rme}{\mathrm{e}}
\newcommand{\rmi}{\mathrm{i}}
\newcommand{\rmd}{\mathrm{d}}
\renewcommand{\qquad}{\hspace*{25pt}}
\newcommand{\tdot}[1]{\stackrel{\dots}{#1}}   % Added 1/9/94
\newcommand{\tqs}{\hspace*{25pt}}
\newcommand{\fl}{\hspace*{-\mathindent}}
\newcommand{\Tr}{\mathop{\mathrm{Tr}}\nolimits}
\newcommand{\tr}{\mathop{\mathrm{tr}}\nolimits}
\newcommand{\Or}{\mathord{\mathrm{O}}} %changed from \mathop 20/1/95
\newcommand{\lshad}{[\![}
\newcommand{\rshad}{]\!]}
\newcommand{\case}[2]{{\textstyle\frac{#1}{#2}}}
\def\pt(#1){({\it #1\/})}
\newcommand{\dsty}{\displaystyle}
\newcommand{\tsty}{\textstyle}
\newcommand{\ssty}{\scriptstyle}
\newcommand{\sssty}{\scriptscriptstyle}
\def\lo#1{\llap{${}#1{}$}}
\def\eql{\llap{${}={}$}}
\def\lsim{\llap{${}\sim{}$}}
\def\lsimeq{\llap{${}\simeq{}$}}
\def\lequiv{\llap{${}\equiv{}$}}
%
\newcommand{\eref}[1]{(\ref{#1})}
\newcommand{\sref}[1]{section~\ref{#1}}
\newcommand{\fref}[1]{figure~\ref{#1}}
\newcommand{\tref}[1]{table~\ref{#1}}
\newcommand{\Eref}[1]{Equation (\ref{#1})}
\newcommand{\Sref}[1]{Section~\ref{#1}}
\newcommand{\Fref}[1]{Figure~\ref{#1}}
\newcommand{\Tref}[1]{Table~\ref{#1}}

\newcommand{\opencircle}{\mbox{\Large$\circ\,$}}  % moved Large outside maths
\newcommand{\opensquare}{\mbox{$\rlap{$\sqcap$}\sqcup$}}
\newcommand{\opentriangle}{\mbox{$\triangle$}}
\newcommand{\opentriangledown}{\mbox{$\bigtriangledown$}}
\newcommand{\opendiamond}{\mbox{$\diamondsuit$}}
\newcommand{\fullcircle}{\mbox{{\Large$\bullet\,$}}} % moved Large outside maths
\newcommand{\fullsquare}{\,\vrule height5pt depth0pt width5pt}
\newcommand{\dotted}{\protect\mbox{${\mathinner{\cdotp\cdotp\cdotp\cdotp\cdotp\cdotp}}$}}
\newcommand{\dashed}{\protect\mbox{-\; -\; -\; -}}
\newcommand{\broken}{\protect\mbox{-- -- --}}
\newcommand{\longbroken}{\protect\mbox{--- --- ---}}
\newcommand{\chain}{\protect\mbox{--- $\cdot$ ---}}
\newcommand{\dashddot}{\protect\mbox{--- $\cdot$ $\cdot$ ---}}
\newcommand{\full}{\protect\mbox{------}}

\def\;{\protect\psemicolon}
\def\psemicolon{\relax\ifmmode\mskip\thickmuskip\else\kern .3333em\fi}
\def\lineup{\def\0{\hbox{\phantom{\footnotesize\rm 0}}}%
    \def\m{\hbox{$\phantom{-}$}}%
    \def\-{\llap{$-$}}}
%
%%%%%%%%%%%%%%%%%%%%%
% Tables rules      %
%%%%%%%%%%%%%%%%%%%%%

\newcommand{\boldarrayrulewidth}{1\p@} 
% Width of bold rule in tabular environment.

\def\bhline{\noalign{\ifnum0=`}\fi\hrule \@height  
\boldarrayrulewidth \futurelet \@tempa\@xhline}

\def\@xhline{\ifx\@tempa\hline\vskip \doublerulesep\fi
      \ifnum0=`{\fi}}

%
% Rules for tables with extra space around
%
\newcommand{\br}{\ms\bhline\ms}
\newcommand{\mr}{\ms\hline\ms}
%
\newcommand{\centre}[2]{\multispan{#1}{\hfill #2\hfill}}
\newcommand{\crule}[1]{\multispan{#1}{\hspace*{\tabcolsep}\hrulefill
  \hspace*{\tabcolsep}}}
\newcommand{\fcrule}[1]{\ifnum\thetabtype=1\multispan{#1}{\hrulefill
  \hspace*{\tabcolsep}}\else\multispan{#1}{\hrulefill}\fi}
%
% Extra spaces for tables and displayed equations
%
\newcommand{\ms}{\noalign{\vspace{3\p@ plus2\p@ minus1\p@}}}
\newcommand{\bs}{\noalign{\vspace{6\p@ plus2\p@ minus2\p@}}}
\newcommand{\ns}{\noalign{\vspace{-3\p@ plus-1\p@ minus-1\p@}}}
\newcommand{\es}{\noalign{\vspace{6\p@ plus2\p@ minus2\p@}}\displaystyle}%
%
\newcommand{\etal}{{\it et al\/}\ }
\newcommand{\dash}{------}
\newcommand{\nonum}{\par\item[]}                %\par added 1/9/93
\newcommand{\mat}[1]{\underline{\underline{#1}}}
%
% abbreviations for IOPP journals
%
\newcommand{\CQG}{{\it Class. Quantum Grav.} }
\newcommand{\CTM}{{\it Combust. Theory Modelling\/} }
\newcommand{\DSE}{{\it Distrib. Syst. Engng.\/} }
\newcommand{\EJP}{{\it Eur. J. Phys.} } 
\newcommand{\JNE}{{\it J. Neural Eng.} } %added 30/11/2004 GMD
\newcommand{\PB}{{\it Phys. Biol.} } %added 30/11/2004 GMD
\newcommand{\SMS}{{\it Smart Mater. Struct.} } %added 30/11/2004 GMD
\newcommand{\HPP}{{\it High Perform. Polym.} }              % added 4/5/93
\newcommand{\IP}{{\it Inverse Problems\/} }
\newcommand{\JHM}{{\it J. Hard Mater.} }                    % added 4/5/93
\newcommand{\JO}{{\it J. Opt.} }
\newcommand{\JOA}{{\it J. Opt. A: Pure Appl. Opt.} }
\newcommand{\JOB}{{\it J. Opt. B: Quantum Semiclass. Opt.} }
\newcommand{\JPA}{{\it J. Phys. A: Math. Gen.} } % superseded by \jpa below
\newcommand{\JPB}{{\it J. Phys. B: At. Mol. Phys.} }      %1968-87
\newcommand{\jpb}{{\it J. Phys. B: At. Mol. Opt. Phys.} } %1988 and onwards
\newcommand{\JPC}{{\it J. Phys. C: Solid State Phys.} }   %1968--1988
\newcommand{\JPCM}{{\it J. Phys.: Condens. Matter\/} }    %1989 and onwards
\newcommand{\JPD}{{\it J. Phys. D: Appl. Phys.} }
\newcommand{\JPE}{{\it J. Phys. E: Sci. Instrum.} }
\newcommand{\JPF}{{\it J. Phys. F: Met. Phys.} }
\newcommand{\JPG}{{\it J. Phys. G: Nucl. Phys.} }         %1975--1988
\newcommand{\jpg}{{\it J. Phys. G: Nucl. Part. Phys.} }   %1989 and onwards
\newcommand{\JMM}{{\it J. Micromech. Microeng.\/} }
\newcommand{\MSMSE}{{\it Modelling Simul. Mater. Sci. Eng.} } % sxb changed to 'Simul' 15 Mar 2011
\newcommand{\MST}{{\it Meas. Sci. Technol.} }                 %1990 and onwards
\newcommand{\NET}{{\it Network: Comput. Neural Syst.} }
\newcommand{\NJP}{{\it New J. Phys.} }
\newcommand{\NL}{{\it Nonlinearity\/} }
\newcommand{\NT}{{\it Nanotechnology} }
\newcommand{\PAO}{{\it Pure Appl. Optics\/} }
\newcommand{\PM}{{\it Physiol. Meas.} }                        % added 4/5/93
\newcommand{\PMB}{{\it Phys. Med. Biol.} }
\newcommand{\PPCF}{{\it Plasma Phys. Control. Fusion\/} }      % added 4/5/93
\newcommand{\PSST}{{\it Plasma Sources Sci. Technol.} }
\newcommand{\PUS}{{\it Public Understand. Sci.} }
\newcommand{\QO}{{\it Quantum Opt.} }
\newcommand{\QSO}{{\em Quantum Semiclass. Opt.} }
\newcommand{\RPP}{{\it Rep. Prog. Phys.} }
\newcommand{\SLC}{{\it Sov. Lightwave Commun.} }               % added 4/5/93 
\newcommand{\SST}{{\it Semicond. Sci. Technol.} }
\newcommand{\SUST}{{\it Supercond. Sci. Technol.} }
\newcommand{\WRM}{{\it Waves Random Media\/} }
%
% Other commonly quoted journals
%
\newcommand{\AC}{{\it Acta Crystallogr.} }
\newcommand{\AM}{{\it Acta Metall.} }
\newcommand{\AP}{{\it Ann. Phys., Lpz.} }
\newcommand{\APNY}{{\it Ann. Phys., NY\/} }
\newcommand{\APP}{{\it Ann. Phys., Paris\/} }
\newcommand{\CJP}{{\it Can. J. Phys.} }
\newcommand{\JAP}{{\it J. Appl. Phys.} }
\newcommand{\JCP}{{\it J. Chem. Phys.} }
\newcommand{\JJAP}{{\it Jpn. J. Appl. Phys.} }
\newcommand{\JP}{{\it J. Physique\/} }
\newcommand{\JPhCh}{{\it J. Phys. Chem.} }
\newcommand{\JMMM}{{\it J. Magn. Magn. Mater.} }
\newcommand{\JMP}{{\it J. Math. Phys.} }
\newcommand{\JOSA}{{\it J. Opt. Soc. Am.} }
\newcommand{\JPSJ}{{\it J. Phys. Soc. Jpn.\/} }
\newcommand{\JQSRT}{{\it J. Quant. Spectrosc. Radiat. Transfer\/} }
\newcommand{\NC}{{\it Nuovo Cimento\/} }
\newcommand{\NIM}{{\it Nucl. Instrum. Methods\/} }
\newcommand{\NP}{{\it Nucl. Phys.} }
\newcommand{\PL}{{\it Phys. Lett.} }
\newcommand{\PR}{{\it Phys. Rev.} }
\newcommand{\PRL}{{\it Phys. Rev. Lett.} }
\newcommand{\PRS}{{\it Proc. R. Soc.} }
\newcommand{\PS}{{\it Phys. Scr.} }
\newcommand{\PSS}{{\it Phys. Status Solidi\/} }
\newcommand{\PTRS}{{\it Phil. Trans. R. Soc.} }
\newcommand{\RMP}{{\it Rev. Mod. Phys.} }
\newcommand{\RSI}{{\it Rev. Sci. Instrum.} }
\newcommand{\SSC}{{\it Solid State Commun.} }
\newcommand{\ZP}{{\it Z. Phys.} }
\newcommand{\GRG}{{\it Gen. Rel. Grav.} }
\newcommand{\PF}{{\it Phys. Fluids\/} }
\newcommand{\SPJ}{{\it Sov. Phys.--JETP\/} }
%
% More journals added 8 Mar 2011, below (sxb)
%
\newcommand{\jpa}{{\it J. Phys. A: Math. Theor.} }
\newcommand{\BF}{{\it Biofabrication\/} }
\newcommand{\BB}{{\it Bioinspir. Biomim.} }
\newcommand{\BMM}{{\it Biomed. Mater.} }
\newcommand{\CSD}{{\it Comput. Sci. Disc.} }
\newcommand{\ERL}{{\it Environ. Res. Lett.} }
\newcommand{\JBR}{{\it J. Breath Res.} }
\newcommand{\JGE}{{\it J. Geophys. Eng.} }
\newcommand{\JOPT}{{\it J. Opt.} }
\newcommand{\JRP}{{\it J. Radiol. Prot.} }
\newcommand{\MET}{{\it Metrologia\/} }
\newcommand{\NF}{{\it Nucl. Fusion\/} }
\newcommand{\PED}{{\it Phys. Educ.} }
%
% More journals added 6 Feb 2014, below (sxb)
%
\newcommand{\TDM}{{\it 2D Mater.} }
\newcommand{\MRE}{{\it Mater. Res. Express\/} }
\newcommand{\MAF}{{\it Methods Appl. Fluoresc.} }
\newcommand{\TMR}{{\it Transl. Mater. Res.} }
\newcommand{\STMP}{{\it Surf. Topogr.: Metrol. Prop.} }
%
% The following journals are externally-edited and the usual IOP Publishing
% submission guidelines do *not* apply if you're submitting to them.
%
\newcommand{\AJ}{{\it AJ\/} }
\newcommand{\APJ}{{\it ApJ\/} }
\newcommand{\APJL}{{\it ApJL\/} }
\newcommand{\APJS}{{\it ApJS\/} }
\newcommand{\ANSN}{{\it Adv. Nat. Sci: Nanosci. Nanotechnol.} }
\newcommand{\CJCP}{{\it Chin. J. Chem. Phys.} }
\newcommand{\CPB}{{\it Chinese Phys. B\/} }
\newcommand{\CPC}{{\it Chinese Phys. C\/} }
\newcommand{\CPL}{{\it Chinese Phys. Lett.} }
\newcommand{\CTP}{{\it Commun. Theor. Phys.} }
\newcommand{\EPL}{{\it EPL\/} }
\newcommand{\FDR}{{\it Fluid Dyn. Res.} }
\newcommand{\IZV}{{\it Izv. Math.} }
\newcommand{\JOS}{{\it J. Semicond.} }
\newcommand{\PHU}{{\it Phys.-Usp.} }
\newcommand{\PST}{{\it Plasma Sci. Technol.} }
\newcommand{\QEL}{{\it Quantum Electron.} }
\newcommand{\RAA}{{\it Res. Astron. Astrophys.} }
\newcommand{\RCR}{{\it Russ. Chem. Rev.} }
\newcommand{\RMS}{{\it Russ. Math. Surv.} }
\newcommand{\MSB}{{\it Sb. Math.} }
\newcommand{\SFC}{{\it Science Foundation in China\/} }
\newcommand{\STAM}{{\it Sci. Technol. Adv. Mater.} }
\newcommand{\LP}{{\it Laser Phys.} }
\newcommand{\LPL}{{\it Laser Phys. Lett.} }
\newcommand{\APEX}{{\it Appl. Phys. Express\/} }
%
% SISSA journals
%
\newcommand{\JCAP}{{\it JCAP\/} }
\newcommand{\JHEP}{{\it J. High Energy Phys.\/} }
\newcommand{\JSTAT}{{\it J. Stat. Mech.} }
\newcommand{\JINST}{{\it JINST\/} }
%
% These are the IOP Conference Series journals: again, if you're
% submitting to one of these journals you shouldn't be using iopart.cls
%
\newcommand{\JPCS}{{\it J. Phys.: Conf. Ser.} }
\newcommand{\EES}{{\it IOP Conf. Ser.: Earth Environ. Sci.} }
\newcommand{\MSE}{{\it IOP Conf. Ser.: Mater. Sci. Eng.} }

% 2012 new option for twocolumn output
\def\ioptwocol{\setlength\hoffset{-0.5in}\setlength\voffset{-0.5in}\setlength\textwidth{6.75in}
\setlength\columnsep{0.2in}\setlength\textheight{9.25in}\mathindent=0in\twocolumn}
%
%
\pagestyle{headings}
\pagenumbering{arabic}      % Arabic page numbers
\raggedbottom
\onecolumn
\endinput
%% 
%% End of file `iopart.cls'.
\end{filecontents*}
\begin{filecontents*}{iopart10.clo}
%% 
%% This is file `iopart10.clo'
%% 
%% This file is distributed in the hope that it will be useful, 
%% but WITHOUT ANY WARRANTY; without even the implied warranty of 
%% MERCHANTABILITY or FITNESS FOR A PARTICULAR PURPOSE. 
%% 
%% \CharacterTable
%%  {Upper-case    \A\B\C\D\E\F\G\H\I\J\K\L\M\N\O\P\Q\R\S\T\U\V\W\X\Y\Z
%%   Lower-case    \a\b\c\d\e\f\g\h\i\j\k\l\m\n\o\p\q\r\s\t\u\v\w\x\y\z
%%   Digits        \0\1\2\3\4\5\6\7\8\9
%%   Exclamation   \!     Double quote  \"     Hash (number) \#
%%   Dollar        \$     Percent       \%     Ampersand     \&
%%   Acute accent  \'     Left paren    \(     Right paren   \)
%%   Asterisk      \*     Plus          \+     Comma         \,
%%   Minus         \-     Point         \.     Solidus       \/
%%   Colon         \:     Semicolon     \;     Less than     \<
%%   Equals        \=     Greater than  \>     Question mark \?
%%   Commercial at \@     Left bracket  \[     Backslash     \\
%%   Right bracket \]     Circumflex    \^     Underscore    \_
%%   Grave accent  \`     Left brace    \{     Vertical bar  \|
%%   Right brace   \}     Tilde         \~}
\ProvidesFile{iopart10.clo}[1997/01/13 v1.0 IOP Book file (size option)]
\renewcommand\normalsize{%
   \@setfontsize\normalsize\@xpt\@xiipt
   \abovedisplayskip 10\p@ \@plus2\p@ \@minus5\p@
   \abovedisplayshortskip \z@ \@plus3\p@
   \belowdisplayshortskip 6\p@ \@plus3\p@ \@minus3\p@
   \belowdisplayskip \abovedisplayskip
   \let\@listi\@listI}
\normalsize
\newcommand\small{%
   \@setfontsize\small\@ixpt{11}%
   \abovedisplayskip 8.5\p@ \@plus3\p@ \@minus4\p@
   \abovedisplayshortskip \z@ \@plus2\p@
   \belowdisplayshortskip 4\p@ \@plus2\p@ \@minus2\p@
   \def\@listi{\leftmargin\leftmargini
               \topsep 4\p@ \@plus2\p@ \@minus2\p@
               \parsep 2\p@ \@plus\p@ \@minus\p@
               \itemsep \parsep}%
   \belowdisplayskip \abovedisplayskip
}
\newcommand\footnotesize{%
   \@setfontsize\footnotesize\@viiipt{9.5}%
   \abovedisplayskip 6\p@ \@plus2\p@ \@minus4\p@
   \abovedisplayshortskip \z@ \@plus\p@
   \belowdisplayshortskip 3\p@ \@plus\p@ \@minus2\p@
   \def\@listi{\leftmargin\leftmargini
               \topsep 3\p@ \@plus\p@ \@minus\p@
               \parsep 2\p@ \@plus\p@ \@minus\p@
               \itemsep \parsep}%
   \belowdisplayskip \abovedisplayskip
}
\newcommand\scriptsize{\@setfontsize\scriptsize\@viipt\@viiipt}
\newcommand\tiny{\@setfontsize\tiny\@vpt\@vipt}
\newcommand\large{\@setfontsize\large\@xiipt{14}}
\newcommand\Large{\@setfontsize\Large\@xivpt{18}}
\newcommand\LARGE{\@setfontsize\LARGE\@xviipt{22}}
\newcommand\huge{\@setfontsize\huge\@xxpt{25}}
\newcommand\Huge{\@setfontsize\Huge\@xxvpt{30}}
\if@twocolumn
  \setlength\parindent{12\p@}
\else
  \setlength\parindent{15\p@}
\fi
\setlength\headheight{12\p@}
\setlength\headsep   {12\p@}
\setlength\topskip   {10\p@}
\setlength\footskip{20\p@}
\setlength\maxdepth{.5\topskip}
\setlength\@maxdepth\maxdepth
\setlength\textwidth{31pc}
\setlength\textheight{49pc}
\setlength\oddsidemargin  {24\p@}
\setlength\evensidemargin {24\p@}
\setlength\marginparwidth {72\p@}
\setlength\marginparsep {10\p@}
\setlength\marginparpush{5\p@}
\setlength\topmargin{\z@}
\setlength\footnotesep{6.65\p@}
\setlength{\skip\footins}  {9\p@ \@plus 4\p@ \@minus 2\p@}
\setlength\floatsep       {12\p@ \@plus 2\p@ \@minus 2\p@}
\setlength\textfloatsep   {20\p@ \@plus 2\p@ \@minus 4\p@}
\setlength\intextsep      {12\p@ \@plus 2\p@ \@minus 2\p@}
\setlength\dblfloatsep    {12\p@ \@plus 2\p@ \@minus 2\p@}
\setlength\dbltextfloatsep{20\p@ \@plus 2\p@ \@minus 4\p@}
\setlength\@fptop{0\p@}
\setlength\@fpsep{8\p@ \@plus 2fil}
\setlength\@fpbot{0\p@}
\setlength\@dblfptop{0\p@}
\setlength\@dblfpsep{8\p@ \@plus 2fil}
\setlength\@dblfpbot{0\p@}
\setlength\partopsep{2\p@ \@plus 1\p@ \@minus 1\p@}
\def\@listI{\leftmargin\leftmargini
            \parsep=\z@ 
            \topsep=5\p@ \@plus3\p@ \@minus3\p@
            \itemsep=3\p@ \@plus2\p@ \@minus\p@}
\let\@listi\@listI
\@listi
\def\@listii {\leftmargin\leftmarginii
     \labelwidth\leftmarginii
     \advance\labelwidth-\labelsep
     \topsep=2\p@ \@plus2\p@ \@minus\p@
     \parsep=\z@
     \itemsep=\parsep}
\def\@listiii{\leftmargin\leftmarginiii
     \labelwidth\leftmarginiii
     \advance\labelwidth-\labelsep
     \topsep=\z@
     \parsep=\z@
     \partopsep=\z@
     \itemsep=\z@}
\def\@listiv {\leftmargin\leftmarginiv
     \labelwidth\leftmarginiv
     \advance\labelwidth-\labelsep}
\def\@listv  {\leftmargin\leftmarginv
     \labelwidth\leftmarginv
     \advance\labelwidth-\labelsep}
\def\@listvi {\leftmargin\leftmarginvi
     \labelwidth\leftmarginvi
     \advance\labelwidth-\labelsep}
\endinput
%% 
%% End of file `iopart.clo'.
\end{filecontents*}
\begin{filecontents*}{iopart12.clo}
%% 
%% This is file `iopart12.clo'
%% 
%% This file is distributed in the hope that it will be useful, 
%% but WITHOUT ANY WARRANTY; without even the implied warranty of 
%% MERCHANTABILITY or FITNESS FOR A PARTICULAR PURPOSE. 
%% 
%% \CharacterTable
%%  {Upper-case    \A\B\C\D\E\F\G\H\I\J\K\L\M\N\O\P\Q\R\S\T\U\V\W\X\Y\Z
%%   Lower-case    \a\b\c\d\e\f\g\h\i\j\k\l\m\n\o\p\q\r\s\t\u\v\w\x\y\z
%%   Digits        \0\1\2\3\4\5\6\7\8\9
%%   Exclamation   \!     Double quote  \"     Hash (number) \#
%%   Dollar        \$     Percent       \%     Ampersand     \&
%%   Acute accent  \'     Left paren    \(     Right paren   \)
%%   Asterisk      \*     Plus          \+     Comma         \,
%%   Minus         \-     Point         \.     Solidus       \/
%%   Colon         \:     Semicolon     \;     Less than     \<
%%   Equals        \=     Greater than  \>     Question mark \?
%%   Commercial at \@     Left bracket  \[     Backslash     \\
%%   Right bracket \]     Circumflex    \^     Underscore    \_
%%   Grave accent  \`     Left brace    \{     Vertical bar  \|
%%   Right brace   \}     Tilde         \~}
\ProvidesFile{iopart12.clo}[1997/01/15 v1.0 LaTeX2e file (size option)]
\renewcommand\normalsize{%
   \@setfontsize\normalsize\@xiipt{16}%
   \abovedisplayskip 12\p@ \@plus3\p@ \@minus7\p@
   \abovedisplayshortskip \z@ \@plus3\p@
   \belowdisplayshortskip 6.5\p@ \@plus3.5\p@ \@minus3\p@
   \belowdisplayskip \abovedisplayskip
   \let\@listi\@listI}
\normalsize
\newcommand\small{%
   \@setfontsize\small\@xipt{14}%
   \abovedisplayskip 11\p@ \@plus3\p@ \@minus6\p@
   \abovedisplayshortskip \z@ \@plus3\p@
   \belowdisplayshortskip 6.5\p@ \@plus3.5\p@ \@minus3\p@
   \def\@listi{\leftmargin\leftmargini
               \topsep 9\p@ \@plus3\p@ \@minus5\p@
               \parsep 4.5\p@ \@plus2\p@ \@minus\p@
               \itemsep \parsep}%
   \belowdisplayskip \abovedisplayskip
}
\newcommand\footnotesize{%
%   \@setfontsize\footnotesize\@xpt\@xiipt
   \@setfontsize\footnotesize\@xpt{13}%
   \abovedisplayskip 10\p@ \@plus2\p@ \@minus5\p@
   \abovedisplayshortskip \z@ \@plus3\p@
   \belowdisplayshortskip 6\p@ \@plus3\p@ \@minus3\p@
   \def\@listi{\leftmargin\leftmargini
               \topsep 6\p@ \@plus2\p@ \@minus2\p@
               \parsep 3\p@ \@plus2\p@ \@minus\p@
               \itemsep \parsep}%
   \belowdisplayskip \abovedisplayskip
}
\newcommand\scriptsize{\@setfontsize\scriptsize\@viiipt{9.5}}
\newcommand\tiny{\@setfontsize\tiny\@vipt\@viipt}
\newcommand\large{\@setfontsize\large\@xivpt{18}}
\newcommand\Large{\@setfontsize\Large\@xviipt{22}}
\newcommand\LARGE{\@setfontsize\LARGE\@xxpt{25}}
\newcommand\huge{\@setfontsize\huge\@xxvpt{30}}
\let\Huge=\huge
\if@twocolumn
  \setlength\parindent{14\p@}
\else
  \setlength\parindent{18\p@}
\fi
\setlength\headheight{14\p@}
\setlength\headsep{14\p@}
\setlength\topskip{12\p@}
\setlength\footskip{24\p@}
\setlength\maxdepth{.5\topskip}
\setlength\@maxdepth\maxdepth
\setlength\textwidth{37.2pc}
\setlength\textheight{56pc}
\setlength\oddsidemargin  {\z@}
\setlength\evensidemargin {\z@}
\setlength\marginparwidth {72\p@}
\setlength\marginparsep{10\p@}
\setlength\marginparpush{5\p@}
\setlength\topmargin{-12pt}
\setlength\footnotesep{8.4\p@}
\setlength{\skip\footins} {10.8\p@ \@plus 4\p@ \@minus 2\p@}
\setlength\floatsep       {14\p@ \@plus 2\p@ \@minus 4\p@}
\setlength\textfloatsep   {24\p@ \@plus 2\p@ \@minus 4\p@}
\setlength\intextsep      {16\p@ \@plus 4\p@ \@minus 4\p@}
\setlength\dblfloatsep    {16\p@ \@plus 2\p@ \@minus 4\p@}
\setlength\dbltextfloatsep{24\p@ \@plus 2\p@ \@minus 4\p@}
\setlength\@fptop{0\p@}
\setlength\@fpsep{10\p@ \@plus 1fil}
\setlength\@fpbot{0\p@}
\setlength\@dblfptop{0\p@}
\setlength\@dblfpsep{10\p@ \@plus 1fil}
\setlength\@dblfpbot{0\p@}
\setlength\partopsep{3\p@ \@plus 2\p@ \@minus 2\p@}
\def\@listI{\leftmargin\leftmargini
     \parsep=\z@
     \topsep=6\p@  \@plus3\p@ \@minus3\p@
     \itemsep=3\p@ \@plus2\p@ \@minus1\p@}
\let\@listi\@listI
\@listi
\def\@listii {\leftmargin\leftmarginii
     \labelwidth\leftmarginii
     \advance\labelwidth-\labelsep
     \topsep=3\p@ \@plus2\p@ \@minus\p@
     \parsep=\z@
     \itemsep=\parsep}
\def\@listiii{\leftmargin\leftmarginiii
     \labelwidth\leftmarginiii
     \advance\labelwidth-\labelsep
     \topsep=\z@
     \parsep=\z@
     \partopsep=\z@
     \itemsep=\z@}
\def\@listiv {\leftmargin\leftmarginiv
     \labelwidth\leftmarginiv
     \advance\labelwidth-\labelsep}
\def\@listv{\leftmargin\leftmarginv
     \labelwidth\leftmarginv
     \advance\labelwidth-\labelsep}
\def\@listvi {\leftmargin\leftmarginvi
     \labelwidth\leftmarginvi
     \advance\labelwidth-\labelsep}
\endinput
%% 
%% End of file `iopart12.clo'.
\end{filecontents*}
\begin{filecontents*}[overwrite]{The_Saint_Who_Left_Was_Never_Free-cover.pdf}
%PDF-1.5
%ASCII encoded figure
1 0 obj
null
endobj
2 0 obj
<<
/Type /Pages
/Count 1
/Kids [ 4 0 R ]
>>
endobj
3 0 obj
<<
/Type /Catalog
/Pages 2 0 R
>>
endobj
4 0 obj
<<
/Type /Page
/Resources <<
/XObject <<
/Im0 5 0 R
>>
>>
/MediaBox [ 0.0 0.0 1174 750 ]
/Contents 6 0 R
/Parent 2 0 R
>>
endobj
5 0 obj
<<
/Type /XObject
/Subtype /Image
/Width 1174
/Height 750
/ColorSpace /DeviceRGB
/BitsPerComponent 8
/Filter [ /ASCII85Decode /DCTDecode ]
/Length 108831
>>
stream
s4IA0!"_al8O`[\!<E3P!A"3Ns8E@W7T)fH6rHTD81YlN@<6BrEZdbYG%kIr$7K@9+>Gf61b9b>E
b/rsF*D&/A1Sj6/g)l.1b^I'6ulj[A7]X-%13OO9lFob+C]U=FD,5.;e9AH$8FLo3=Q<kBl"o1DB
NFtDK]H0/0HZ%0HbT_0fLj1@r2Y^:186YDf0!"B45UfEbT)9+@9X[Ch[br$4R>ABOr;REc6#?+DG
^9FD,5.9lG/lFCAm$F!)l1BOr;cATW!'BOu'(6#:XOE[`,?Ec6)5BQP@dEbT)qEbT#s+EVNEFD,5
.@<?U&DKKT2DBNk0+CQC:Bk1dq+E2I4Ch[B3+Ad)s@UW_^Gp$O;+Du*?DffZ(EZf=>Bl8#8Bl5&,
BlbD5BkCp!A8,Qs7WN3SA9)BoBHUl(C2\A?D/XH++E2@>AoDg0A0><"E,oN"Bjl*tDf-\+DIal&E
b/[$ATAo%A8Q%$@<?4%DBO%7@;[3%DBO%7AKZ80@<-E3@<-:/DIakE3&**>+B3Q'+DbJ(A926hE+
O)5Bl5&4@<-I(@s)TrEZfR;Eb-A6ATMg/DKBN$Ch4`$DfQtBBQA6:FD5W*+@^BcARogq@q[!!DKK
e>FD,5.Ch[`tCh[NqF=n\-DBLYn0ea_lBOuH3G@bf++D#e2BjtduARlolF!,F5DId0nFE8R5F<G[
=AKX9K@:WqdGp$[CBPCpiFD5Z2+DG^96=k1RBl5IEGA(Q.AKZ).AU,DBATD?)E-622BOPsnA0>?,
+EV:.+Ad*$F)u&(DBM;eDg-(A:i^8f@q?cPEb/cmEbT$!+BN8nCh[<p/g+,,ATME*FE_XGE+*HuF
D5W*F!,[<Eb-@mBOr;\Df0>5@X0)<BOr;bARQ]iDIal3BOr;P@U`no+DG^9FD,5.:LS$_Df9/=%1
3OO<+ohc9lFob+C]U=FD,5.;e9A[BlbD>DId[.@WZ0lCj@.3+DkP.FCfK3Bk:dc+DG^97WN3SA9)
BoBI@jDDe4'7Eb-A%DIal3BOr;qDerrqF*&ODBk)7-Eb-A;Bln$&DBMG\Ec#6&+B!,p@;^1$@qg"
$+CT=6@3Bc-Deru;+B3#c+EV:.D.PaNFD,5.FD5W<+D,>#F`M%9De:,#+DkP)+EM+9+CSnkBl8!6
+CQC1@<?X4@;I'&@;]Us@ps=f+Co20Bk1dc+DG_8D]j.8Eb/f)BQ%p*H>d\C@;I'4Df0--+E(j7@
rH1%F`K)W8T]"r@rH4'Df]W7Bl@l3@WH$gCNCVBBln#2@;Ka&FE1f#Bln'-DK@EQ<+oi`AKY].+D
u*?Ci=3(ATAo%DL!@DATDj2ARfgrG%De,AT2a)+EqL-FE;A3G%G\H+@1)lFD,5.@W-@1Df$V8Afu
2/AKYr+@ruj6ALnsGBOr<1BPDR&F(c\;@;]TuA9Dd$F!,=.CM@[!+E_WGFD,5.DIIR+Dg*=FFE2)
7+E(j7F(fK4ASc'uAKZ#-F(HJ&F<G"0+E(_2F`V+:@;]^oAKZ)5G@>N&F!,R<AKYf#AooXJ6$+5k
BOtmu+CT5#GUstZBOr<.Bl82=Anbn#Eb-A3AftM)D..<)Ec5DuA0>T(+C].q@r!3$F!,X?F(o*"A
KYT*Df$V+AS>dqA0<HHFD,5.Df0B:+Eh=:FD55nC`mG0+EV:.+E2%)FEMV8/g+,,ATD?)BlbD7D]
it?BOu3,F*(i&Amo=`4!8*XASrW4BOr<.GAeUF@;0OlDJ()$Ddda/+EqL5@qZuSEbT)qEbT#s+D>
1o+E):7B5)6lCi"A>ASu[*F(&]mA0>f.+D#8/BOu3,F(o/r+E(j7FD,5.D..<)BOr<-F^fE'EHPi
1FDQ4FE+*HuFCeu*Dg#]4/g+,,AKYo1E,oN2F(of/Cj@.6@<--+H>d\4+E(j7FD,5.F(Jc&D.R:#
F!+n/+D#e>Eb0*!Cj@.>Dg*=;DfTQ<Df.0M;f?Ma+D,>2AM%t;G][V?+E(j7FD,5.@ps7u@<3Q*F
!,+0G%GP6Dg#]4+EVNEFD,5.A8,XfF`V+:F*)GF@ruj6AKYo'+EV:.+Cf5+F_#b?F)-%R+@9XL@<
HX&+CT)&+Dbb0ATJu1ART*lDfp.EDe:,6BOr<*Bk)7-Eb-dEBl7NsDJ=3<+C\nlDf'',+EV:.+EV
XHAKZ&:@VTIaF<GL6+EV:.+E1b'DKKT1B.b;uDBO%7AKYE*@<,psATMo8De:,*BlbD<D..6sDImm
1/0K%JAKYhuDCcoEDBO.:Df]J4E,8rmAKZ).AKZ/1ATq^++DGm>D.R-nF<G[D+DG\$B5)6p+D>J+
F(K0"/0JtEAo)$gFE8RIE,TV:FD,5.E,Tr3EZf:4+EV:.+E_R9G%G]8AM+E!%16'JAKYo'Anbah@
;I',ARfXkDJ=3,Df.*KBl7Hq@psIjA0>?,+EV:.+EDCG@;I'*F`Lu*@<6!2+CehsAKYE&+CQC6F`
MGABldi:+D,P4+EV:.+CehtG%#G7BOPEoBl7X,BjkgoGp$R-ASrW-ATT&=Bln#2@V'D!F(K3"DKI
EPAThX&+D,Y4D'2P]EbT)-G&M7@9DiF6ChI<s+EqL;+D>1o+D"u%Ch[I$+CSetBl[cq+D>J1+Eqa
ECG$`2BOu'(F(f,-F('88BlkJ3DBLYn0KC[,9jr'VAKYMpF(96)@V''mBlkJ>D]i_)EZet7Bk;0m
+@^BcARo@mBkq9&<c;eiD..="3XlE*$6hbYG%#H6+D#G$Ch7[/+D#e:@;^-t+E(j7F)-%D/hSb!F
*),4Cb-]HD]j7<DIb:@DJpY7Df9ME+Du*?F*)>@D'1DQBl7HmARlol+EMXCEbuqCDfol,+D>2,AK
YGnASrW3Df'&.@rH7+Derj(Bl@l3AoD]4FD,6'+E)./+EqaHCh+YsF<GC2@<6N5F(Jo*Ch[Km+CT
.u+Dl%?ASl!rF<GXCD.Rp)ATD?7+AZktFD,5.F`(c.A8,po+EM+&+EV:.Eb-A-F!,@=+C]7tF=\P
PD]j+7BlGUGDJsV>AThX&+CQC6ARQ]uDf0W<ATB=E@;]TuBl5&8BOr<-@;]TuDJsV>AThX&+CQC%
Cggjc+E(j7B6%F%F"AGPDJXS@@3B/nG9CO@Ci"/8Ao)BoF<G:8+EV:.+CSu)+CT.u+DkP&AKZ).A
KYf-DImWuDImm1+D#_-DBNe7Eb-A(ATMd+@<?''@;]TuBQ%p3Bjtmi/g:`3%13OO%13OO%13OO/M
/Os94`BmDfTB03Zp"+/M/P+/M/P+/M/P+%13OO6$79_DfRl]7T)fH6rHTD81YlN@<6BrEZdbYG%k
Ir$<1_]Ch5XM9lFob+C]U=FD,5.;e9AH$:\rSD)reY@;0V$Bl7Pm$<1_VAM%_4DIkjO0fUjA/MT4
B0F\@f@q]e!CcW\/ATDWrD@HqPH#R=U+DbJ(A926hE+KIcs4[N@!"0#;#m^bD#mgnH$O[CT)A*Mk
%hKj"(_[f7+t+id+seKa0Kh3A/NGF/,rJ4N3Boet778R\92J5I:H_'563#[>6NIQ'$kNm[*t]),5
qF<55s[e&5s[e&5s[e&5s[e&5s[e&5s[e&5s[e&5s[e&5s[e&5s[e&5s[e&5s[e&5sbOd!"fJ<mK
AgX!?qLF&HMtG!WU(<)ZTj?!<E3%!<<*"z!!!$$!rrH-"pb1H!':_m!!33&"9AQ.!<WE,!<NK,!Y
GDU0`r%S;.1Yg";<PAOeK#`6GX?\;?kP-n.fF&1NJ'm6W90J_uLVT!<E3%!<<*"zz!!*-'"TQC?+
;"VT!<<0'!!<6&z!!!!"&Krk1&j]B'E>&HSg&MN`!<<05!tbMt!90RDF2f&siigko/2)t;l-?7I^
0IQ113%^',.UU?\@utM`3Se/nY1<l2&S:LemcMb1hc&S/SsH!'TZ-l6h\MFcfpQ.-Ohr6_E2+0jj
ZM8*7p0X5,,afc.m'KSDR.mo9;Gr5IM>,HG*^<g7$#SX#92]KTrU9bVNG6+&4t7\C+j-;VGBF<$N
YFlf0c9dO8Tn"7PP7CVAl<$:[Gi=bF0D)qI1GlP$^l,d_+U$.OgP(pO\$efLRAr-QEIb56%hDoPR
a`D=q/Qc"e3EjE=9aHDfb;0_,<1?3aoV5[WOdYtg4X]63/pA&Ob)?h,LhIsu&>:\9?:`sgjH+['Z
gO*Q7RsFO\+QpS@16#.L>i=s^qR]A?+^s[:`/T'6H(E-9dOmCOT@PR>VPoFAS2])GA@[um;p1W#,
%92f4`BL?[36`Dmu87AY8S&e=)EFHUEjVUDg3[TG?):B`/V6j,#`qH*32QP%WKq$`Ho`rAP9,F*a
I5)_?mL()0XIGn&p;NWr,!;V='V[aK="[m,O^?OJHVNq[;6?_?^ddO""52&:koeFfu[5B4[2_O!$
ei:$Vi9h=TM_'9mEZZ,Jo$8oH4>Kr`DpA>X4C3]uK=Z/cJmGE)Egn!=V==&&S.4!190bk.pP6SB&
-D_KugPf_rgaW;[=mddTGgIE=T_C.qfcaDfI%lJVd9q(a0GJ2GU>VoQT:b0/R`/M<&n]nA]<n4>J
`S><&Y<:Ys"5PH5=,?Mph*%se"^:>9Dth)QW\RRd?44Lc?Grg9ONC.&Wr=`,mC;NSGi0GoM,mq.O
:0`k`b?Y=m9?]#;fD&_K3DR'hK35Dp@h_h0`M-blEIBkBN.6nrUBqF<;G]6&j&mK\BRFt##=85Vq
]CS8ts-*W\RI\gf$Yu('tpTHsu=hC,b'1XD\,L>LdXlO*YT6S5,S[b-tWYC.:,J`8O_r!,B1C`r/
Du@Wq(Kkeg<c\)/(&:g$fflBBd`k";VTGoWF,JrXg#76J;)hp?C@^!$5iS>\WI]Fp*eHT^:`](u-
11)]@1eUP:`.3%[sN)3tIN([pNNcjA<Wi=OGX$b'c\lF(-7b5brj$kT'LBsssBV_Z4c!l.s]eg5=
leYN`c\%1O#i@i_O41S,Vq+TWp9"8Q`.*`U%-+g7,^&k(p)[df]LgXhT386`-O>S%4l7*U(J'ZKe
hlFCbPqR^/MEH9RjORZ@aJf_=H"i2DHqYCAZrN5mfFj(^.ce6>#J_(FDD,$*Q5N<E(Wg3qH"I3Ve
^\@r[OP.3YUNkP#@l1XASB6LTL@Nl^YNY1eF9P^%7$N#u#`XXIi7i*oGbm?_Fmn.q-G"+/hphdn@
N7]qAEN?]99[U/SPEP+2L^jK<Ka2bBsMd8b'^.hiBC3<I'(0>m)m=3lpFR')9'dtl_\8lmGFKBFU
nT;tI4c4JX9^$WrFri;d-EtVr!eepq3X,apl@a$I3A-&6:[f\+@D*NCN>-6mJHQq32HKR%1/7WeT
NG/rHGe<EC`GZVPf,sfNjVX0)A$'5sX5LZ0`9a9XrUJ;mM+PBkO8um54<!8C(16drH`Lq@15Xm]`
r7\<l@Q`jiMZ*EmuB1:(\I"%RQWJkUdp<X_fpLo2:=3;-nSs']t[/!'KqGcO702[NRaKEDX^7i?F
*.1-dTFG<>@XY<3i&72iFLk?tNl^FQ1AG3B/fUU;'IY_UosE-meDD.#j7q)$5RUD_S8b-#o?=Lsa
KsDfm5kS>p5$fe-J"MPkT$^Es.C:h&6<o8m4O`nK9A'"%(BH%Y3K"*$a^*&nHi)d>$%LlDr4Jbjo
.0*b3Q[E,3igq'oOnWk=E)/2(eB\M,&F]eB]:g]LN-qO!l<,Z+9Lb^8Sd<j'-i9[U(P!H7dh[!HS
iObBoUnGS&(L#&qZ")tk`.kYTT!1Wp"MA9ZDVptV;!mZqm)t3\8iqC0GFEnc]j[hE2"0T9`_tGHQ
OQPT>9tcrGf_GE^PoK:G>YPh)9.rR*6hbf\:luKisbqf&Bk?Ck#i(fdO2]TOXAU__O4p0IA,Zhed
e-4Cr&18<I[!a5+g7=>EC[n/RJ5labt8M^DP9h9r(gpAHi94))m7RV&nR'W3"XuM&e;?ml#+)U"=
WmS8l.5M7*#@NICC>,rei#n;5-eh@ZmPR656ZD];1RQm9T+\Ao]#)5IT]3JOrS4ta0A46$T(D0j'
VX`LT"59B?Y>Bn?,>%l=0N."bEOlr>S8FAV&OCEdX@@#^gbo!*<k8Xm(mj7#pXZ0U?IjT%Ip$VlD
\+>tc'Q6`^DZ2SaFFgT6l;lf$f>L6elh/OSO?%NF#4=,5+>gi/r8b+Se/H_QNF^1h)DC,&,U7=ho
:5*4=eD^e&pQe,*!]e0an5&1h`o[k(faUP;]G\i2oT2>>m4YtCt/l1o1Cc0>YXaXPq]IRfgC#Bp@
G')#_f]E<J5I"<uIu6,[ZNled2WPj.)d)9>@!#pp1kcT3,*YRCJSjD<3)M?_-1O]ije_W0jZsn`-
-GaWdq$YnHjr0fk'gS%OWCU78d2gN$T6o!+(+ebu`s*PopeWo$Z'f)0*S&8tbBJtuD#)_W/-)dbS
k#/dp/gquVOfkQMM]M">a;<&Doi?bPA^J>(43'2M?qr+liphTZ?=1NIRipif"G,6(XNF+ZE\ng<1
JQNGlgPJ:GhESZJ.b<@mm?54U`]a\?,p_@F@OZY2=_/INENBV+<PEU=!(X[2[Gr37CISYHee11`r
r<pXJ4U-_/'!T"]5)oA+?V%ANqE4,0DES>G0RR$]O?2C2iJ\n[W5c<^F>L996^*K>4Ps-QE$_s`Z
kQYg$:F]1472Wj"PFV:=?IldC4QYf:q/]S1'M<Os>Rf7QFc!(p1.4NUa>.#[YVE]iPChO/BJ2'h!
H-@RWp/^m/UI=G.:X7uY-*8(mKtP[r-iQR<nT5elUtFVttPgDf^tUW_uC&j=)dWm(jem:iE0Ib't
$\X;nR4C^Zn/[H>aB_;:nrpl"=rE[q4P)qRZGFQ.0kn%(U:^CaSBScC]p?+4SM\[MrW7VhkY`2>O
FDKJf,rmOW53k->^#I1@6A.+-<SY$d)7s^hD>0JSjQ63gAkB0dPoc;^4*4sEh`Pfm6J:I6bXE;Ab
:'8rV<mB8D+4'A,_kO6Y`q6;bt*S]l3l(00MCfKaC9JU$tXGLpr!96.9G&tNF_e(-_GZUCR976[H
Yja&;q__V"]CnJ\#*8TtOT6B'pnQFc4=A#B0Sh[hG_X)HV9u^JG1#BWK3GA-XL%AQ,=/aWHN&<NT
s1brE!&DX_c1T+1f#38mG&AbNtWgJYMI1rUCi5>O;3+p.`!mgLQYm;IA\]ngImPEZp0brI-Y]Xj4
Mg2.7*b_sOBA=0?C]QhH^<Uf.I6VF^b@p\b3@L$es2i+N7H1?WB@n=755D\.473$38NpkF8]EC<4
V%Y[nrAq&i.\14`/N[VS632MaCCa<e6@Sa)P!;[-k8/ZT8A*:j1CGUj[7e/MX)$HneQ,G7gY?C_4
B_fb$bqXLdTZg$AC"=IDKHo(8i9ta0-Y*0!5;H_(]fNa8a"Ar"?[L@c3nlsehGJ/Ln#N+=PSI7Nu
:D-]*ce+P]j_7`W3e2RLPW$`K(^5n'4%IQ^>uJ36H/ZifB&H"]P&I52:ojm]ZpJe/tTK0d@;D'dN
73`CUUM5<c=Wl:oiOoc\&'=aX[590`8b3E1ffG2N,6QYlQ].kjGbW:B1"gUFXpRoQMB]>^:VX04M
[QT'/9HYujaH/:^k&@o!%_SEMY`]`5*(jgTu%Ghe]l0'V@:=)e7??3uV+E%[F%g?AQMmHd91u_6W
Y?-_;(;=ooU=BIKQE[>O0TU.20SaIu8(Y`sVQr&5N'`%Xo7:f"YudN:A(/#)IlQqNn)!lADTMrYp
W[)LmF'O20fjp-Q1Fdo^%1P#Ed_'5\Kj5Q-R"BD>]-&'nm2qAOD01=`!Fn/n!=TaZ=km=3'H;.4Q
Y%%mDY^Ja#@.;J99eSFIi9E3f@J%c8X8Jh?Kf,n//u&8HkCNM:CeFe"8YIK-R`fpVGV%X3BSg$mM
2$is`[45jj]PqoXU"<W(d>Jdbsl0o!5Ooi2]P_]*8UbVIb%>]0XEHG<Ssh9mUeBu>q[RmSTa$^?a
QG%4jRgggUFM"\)W1LZ68`+6s6R[,J.4Omqr#i!P8B]h@^a#)a"^Pm'b35QQ3aMc)kY&rS0D>Rlh
)`(>>G*euD&1Y6B(r0tBL:^Hsb(Ydf27!^dY%moD"B@JRO$FK<gDCXhQKrgPnlH\uMCHdP_4rJ:@
HP,pNWmg\/\8:WP_`Dq=Xp;cf3I6*7/3jf8'=X82l[9\'8ZA>pOn5j\2T\,2C>Z50?IPkURa@.`U
$:2WP*A?MVXM._9LJjLpuO`nmbH(^#e>[+\d[MG(I@jk%KZRkk14k>HMN)#/@5o09^Z"BLe*n;PI
6Zn!YtSit-E^jjYQf;HSI%Pb9&rfK45"Ii2ld)]e85Mm\]HEMFPAl1C-3q9;YpAMi*0@S")tGEhA
CHE6Xo+V-I(bH$:c-(n`@[I8=#]=LYZ%'Go?[Y3K3Y&HiXN3mnN&?C3:m<`O+p?$HBV=bO90_E\W
nVGqB_h93lq;X$9&<3Ynp?&>B339"QEc&qqIp&r$A?9nK^d^ON@3Opn^)4br,?TMD)lM7<@Zu7#b
1#)!8"9f!HHfWMTRN_7nV/nOOB(n:@fDh.itR2<Bh'sZ1+D6V]['<[$h]j4@;2q"S5\U3=e:o)gN
qGUickrYhD5tXUrV`>NSemXfqBf//LIf&gWP%#R&-MR@sq>@EWp\gO1%ei>J9>SftGEbDXVMu-A]
p!&9\!Q(nX,CXaF]VddEokY8-:\X$=LhWh"Q)dmgPKLQ-3n;5%[J(u/D.J2(`o\\*!]VHf1pEWHW
1fuB(pX9+G#E["2&iqTi=jeONAA\0i,^GVL4V;^L[9+fDs&S6Et1kc:gN3'mZgl\9>269stZ$^uD
=IS9!ZMh&qXrkA'U>CoN.T.!\D8G,a.ZIJ/l&"K;%Z+4*o6G<Ul?@`c',ZQAg\u0k=Vq<cgR9haZ
]t:.,r)0Zc.LRRft=M%[.rf@UJY-[H'NHn4P%KkCU<9*isjJ_3DlY=jFUa*)jR!KITTm%]#t\g/)
,V.KDELmOn2Nkb!P2[qf2LL](:e:.0[I$Q9\mlC&P:01>o`b7]#%EGcn$@2(o8%Kb[kT@Lpq;`YR
BMG\ot<WkoBG(?iCcrE:m0qL*bO!c@U4O\+'im^pHEp:ie3SO1J)<\4`1C[[R^8ih<oHc8PLfrF_
1mG7MbbQN)0fp+q'>dU#'3$JJT+M,`g,I(Sn^Er9Z#NlpKg3knpnU]rDNuR\VO:<-lV@UkF<QFib
cG=`b_EqK[f?R\V=Pq0nL,#GpDnJVU02K'T=hATP_8]NW$3(!3F+.Hm\G4&5r!::IEP=SI+kp754
(*\@/&l%)\B?WqWNfN`8];UM54\(VmdN'UJT\Vq/!=j^W@C'm4)-$6UFFnX8i&3.cLWNEM%,g?"!
6OOG8YkZ-pN(=-e&g>mX?rO,Y!"(S`LURm:CG1L8et?/ngLIf:48_f'\7P)n0jrrEs.V]s6&T)jo
$t@2lXHZ.Jfr(obLNXO6tq]#OSuN7LggJl]U+A5N,lDoQ"liOcr]UAT>Z+:M$"5fiZ=;Fh]^V4pB
u)DW_]<R713AJ1IHa`RiBCr_#]?V2;dY]V0Se=s/CbZLgk-hPseilI0V6r&6%Tc7A'brSh+r0Y9Y
I/(GnNM8ObS?.F000J6UAG3:.0(K1&ZDe@>j*%g.^Pj:<8a@j*+p<)ek.OK1o_O!UhnH@Oh^@lX,
">I;5ILRZQ*IYJ%X&X9Hsr7F7b#URL/J+1UNN_"ad?g>Wqqt:YZ5#t>aifYS\\P#JsJ'1pFdnQR(
BX0[,E$@S7rLT.#i4J`[=EX_AjQ].5nXA$R`qn:/^KDpI%*GY;sTr\)GnE:Mc/RK6XbPAahlH@R?
=lm\Ja!J+gu&arQi&kWS'<FPRh6&JjiB&r[G190[7*H9G#uA/b42(T@o#3c+>EIfl6'P)k@h"(q!
O6V8<fRScMb23:F7rAM;t)NqjO^L]s<jQN_R^+q-QO]U*YqPcG/jHW-&?85RdbN?)\Sk!kWb?f)g
`(k/3A\+FqMqcPq#t:]7WOcoB6uQ8OOZ.Kd?`["d*oPnfAIZ\i.fJE1Vtjfr?>OM<+Y3$+RdoWJh
)OeBOH"\Y,3<7LQs+Q5@S2F$&D)OQ!'5$L6%h1AdJ>g]L7\u'217Djc<4&,Oa,88Kt>,ekKZs;O=
adUOsd-N,YI'T@->T.@&H>,,QhDOqo4..Vi2`ZEQn_9ZRqhn!Hm5]?8.?FbMs&gXZg^To-/\P'ON
_Yg,3,VRZYtSS7)1tenkh=[A;(%FL>8+^_[s!GQ)P-&7->>[E5L[V@B%4jD=RTF3;.fCm8AM=8k"
PNUa\k`9f&F;'3u"B^6*0guZ.Cg\rp0PZ[AS7q^SHBYo`:H&/h]@2]j@el`!GQ+<&5(U^Na=2&[d
7oo*0,W(?JGFTAjci2U(htGm_.ihRK2i((33Mit2DV"U6).'@h'SfKA*qO?u=O<"rhi%tmQWs4M-
%Sqn!6&3jD;$'K_:rmP5McYSMhlfN(2q]rWr?tG7s0E&X!M]04CEW+EK`(N=pil?EZB)3A#C0_[F
_M,!S=M8W\Q#u(UURB`_sF4?[NND='*-RPV%;7AR'!O)N!4NAKq\,HGLf%>[D0=H['nHRLuM?:H1
nrShWrF)*:]ZG?j#P#^#<4h`(oWqFm\oDO5X='4P1G_F.s;e/s21L-7D2hqL*4#A;Iuon2h6I]qK
D,At0o!#`A)'=?ZB>SF!F5FupadqMq.@n>np*Mc$.2KGU`6CZMFJol"uc=@r6S72J:K3W"K`H0*G
qU_7RW[^4'X=\cp[K\n"A(j#R1H83a4l>jl&\8N:?iW07iJUd3gq;:!<B'*"$u/\,fZ04_j*^:d8
-\-8WDOO'\U:fK+nUYaD!BC>^BQ&p8j\)d-//N.;\E;4mj]$%Y%BCGjS+sJO^W(6lRNkVWc19D\H
6>06+?hLF/2M1PK0b"_U:KrS?`QI(<fp,f,S>*:*;&emJ!$Ef^ji_o^f_+mT.raAh9sGdjs\Gk7u
c?nCJOqg4HFk5&h_bO^lBV^sJ'1kZ%/2_mW:>=,^XdJ+^W$?`NST/l3SP.)fkrqHbo;4i)O)m:C#
AibEF'G1H/UhDD">-ud0OAYW?Ck&=XA8d(nY\Nd/9f6++dPX\(glZbI4S5[F,/>=;8Gi!27!Mc4I
!pPi^(bE%(_i+mX6dO$^4g)uO8cVjugD2*hlg96TP"kD(DV3]OIBfS:O[^?2Pdh?uH@1([Lt<C9c
$uVI4Ca"\mIoVr>Y]7iMXk.!>?`k=/#*#J"DR97N08i*=t/Kj&hAC,poLO*CE%7;/DVo]U2T<YIH
<[FG@5g\dkd0_prAP[g#k6*,(&"D`>o;7f^[Wb3:W$/^O42$L7N2$,25J-=P]lmb=lV$IsciA,@=
Gs1F:).i6&8L=VZ5(]u1:@prc4!0ok6(QR=3);YQ^p)(F=8&cFq1aH()FaP[T/lF^][DRLFpDg.B
ZT9[_[rc9EYlR<TC9;f8r`B0mhpA*Uc0*)2.0-'*+r<$&K0`_ENcA;,Go\#frXO855Kc*r/N'^$e
S>oX32gI6D/Z&Vc[Xl11P)jVZD"u&/c*pe5n8dRl\&S&K=HTW@S1QQ,.uPh1m(A4F)V:JX1/jAM\
$k&7Y4c<tVD!,^-RcUg+'D-&RU'o;k5uWgX61BDKB,&7!`,`e:n%Ca)5@i4pZ!_nl2-scQ>p[%^'
P(e*^*uRMa;Vt4ehl<aab)9!l2&oT6\%Cl+]EeaOSg@CGb&agXjG68$j&(MSlsHc9,Sk#6+]"@'j
[Vc4I:D8/1&PQ&X8;)"NR*0r#=pdW&#N&#q%YCMl"7P`YnFS8tD>?c/AAV;+r<RZTEO?;<3A>7*(
1i6<+f\")/'Ff-("L7^)>:=V"9/llW5NR7\*3B9O)\)Qur\e3u0rc(-&o663g?.e1ZqI`i>H_67H
7&jG#NqkGON*-$98'Z4!H<l]:1(kcf;gCPV6#s@iCDKJVn`)#VET?c7oa%i@?U&AG=Ido;2sVa<!
S_KrZa=R(@hCsEIuMWS*B7'Q2sVn]+aUuRC/r*_XCQT:H--et-W<d?o4@cH"C<lrjHs__,tHfm,o
gh:Vb'7@N.,i,C:<$7]V-E+`[qfUhJrC>B%^"]`t+*c(kb\n3FH^B8o)!qRC'!,?ep`).UV_75`W
nF`\^?bSBuh7du>qZ++/rArr@Fcmaq2jRh=8l1h1Rb`eO`T(d`oaE/))H:>X_L\;8Hq(uRKMWk$&
X]AU[gj&nafMLDb';JNpc0"EjZ<ZAD0h`4"qX,=hQIm3GH'_Uar,m_ZAf.9OKg!Nu;XJebb#\/<g
6_8Lh^Nl;K3I[DZ@%gAWmjLH^Xn-d)l*ad,XHor@9nU?>^K:F8FcS,ceRK1%rc:u7,n[mnN"m^t_
qRh3)o;Z(Vju89K`[Lrm=@QuIPm&kF7E@kF]Ft[BJm5`ZP,hUS9gh?X7:i-8)oO/'WbH\lE;&Pg"
MbCX#4\dpR/6YD>C(J5>@F.Z&%1XX5DeM.qKcGRC/oUaL?_dbn&)TS$qfa2bJD7"b(u:gr*Oh``n
:2=,P/;(82a\iZfq"J_f9JV16c3_l'ZnI-.-']pV`o+t:/5M(%n_,1a0``pJ^fH)Drj98Za8hB5_
C]>q,0Hc!J10CP'1!_tVQe`$HRS`"'tC!,Neig+2<P!7$[0Lc!ACNj.*4[plt?48V)?T!Wd8mLT)
0+Ll6;V)&re\Akh!"iTde<lD^qDVPX162A7N*b$_f9[csqIsjb?+Mb<'^HWk\&Y%MWXWnWrFVU8W
u=.VP3tHaN*4m)Fd]iUVs3(2ira[[GsKrLeftsXk.]H9n`_LPZe*30E;E:Vit3I:lPI8j/"1F2+]
$SsQY%e\=alofb$%=0YBZ#0lu@_*()OVd7)!j)a&`!pLILf@nam<U4*=$_dpZ8B&i;G3k9sArGg;
G57c3RSas5s/o:Vs3]Xl5-*c3DAM+=IN<4DaEN(1du14g(3pZE;s-+SuR8IKlTU\E"7iOs/5`il5
B^-!ThZpjKE6H1%O^OGa:l#S?$nmf[&`B</)m0MD`/J*5Ti;N2%dGWS6]=YNHh@te]j+fH:7$a<+
0Qne\S?H^/]U6W'Uu]Q_;D_kV;.Oh0.j1*qp;2f/?7R@s;?XRL(U[/Xo:ceRDir#+'!.gMEd@K)`
nft;4LFEMXMe*3I/1gfM95l#iT@@!h:9&RU`",KnXHqG2>lH2%`l\P]IL?+=RqC(N[l,Sa2VbW%a
qQUbpo`=_SPX,27q)NF`ubL)Fr/_HG!5.\Nd`VbPH_K$?h'5c$"rb3?q,a$XKC%rrD(#0dKk5OW_
A2bJY'QCLqE@Sca707`7H!E-NAJC:&=eG+MqF(/*g!7$6`N5[*ur[0+!s:qelG%r?/SR``db;;1i
^=%r[,<@)NAY;r<r]#O-L]t`,I"IFoCF#d]q^9&1N!0#<WaZNm1<]gKN`/&VTLF1IUM9C6#WJ\'<
YB`Z/\@ST/@kfL_o9Rc5]"7cK7S'bBX1Om97\.^tY-:AClE-XuTc%+[o>hb16)+1M.5"m.?Ek\)p
ZEp0i5VL[D]%_nl"^Ocb#KMOZ5@pco>WU<NnEm2Ir&7UZju+F.)/Pu=.NeF)?!##XI?,N8SfjTYU
5\c4\"4k2C1;A$HMJq::HI]Dk#=^9604aO+!kTia,HgNu>C,6Wii7e6=HBilVu(h\id(Io.*tN)%
gE^#_S<&%2tt[AL5Z;MsAWN)hJukJg[R:90>LU&['`K99`O[aPsTY]X0qR6V_NbK3D("6Us4oh5I
:n1Y`br@@'N1kD!n0PWE[C>6iE,^.NM?9PKLp>FPUFbZ@EN7=e?&O08s]L5=7`jj.pf;G,IDoH(R
dPR`4[TE!TKVRh$b"=er!#iZd0jQc>'9k%3F$Q(5,4*k?b'0k,-.a[9+AS`#2(:_p6jIW8S><W@X
O+]<0)XeWPNf`1A6>cCWm7IpFI+3r5\f2Hbuj-iInZ;(35dk$Y3S,1<8Go-o0]GJ[BqQmSVsXi,-
K6.N,8g9NS7!(B]jj7UCtOgQ7K#0dY,)RXOs'XKUlJOfsukfoP6r7d1*\.6B;gn7!"$!R'ObTL)a
eZ`31&lAFRVKo3S[r)B<3C5-gqa^A1#d6LeP#dDe!Ir8lEb$K!HOR/O?ZqH=AM6XM4BPQEJtB>KQ
>n`.nSG<?^P&^k`^!,J+GKoWEhg#9sKXVF1a'AB&_AFI,Dm=j/E%![p#O.6[=SPbFg?RLe5.F'X_
H,"VoPmD'<3293;4Q]5;GfGPl<EW/SV%V<up;K\TiLI]Vk;6as2#GMI?kI`k@#RohIPC_#DUmSV=
2JF.]mt%HB&X*L=R:%#G/fuo%7$2%[B%;W+hnlbX.ZP'EKNYu?F*jlk9F=nHX"3JEmcmPYJ"Co>I
S$e*7p0BRPILC`KXZocN[\>lAKT#1r?pskN=^P".ulWWm007b=F9e=2kQb!h6h4@OPH\ij!&)-r&
Q*Z)bXo"L\)]18mF:O0.cg3S:q5JoPXIljon?BZMt9+WAoQS2\(&KnWiNc/3U5B"t9:qRNphA$1:
5kT1E8?lu/U``l+=VVU!%B_9rV<o_0'f,elD\L2<UEM>rt@^WJb`(IVt;!G+5#,5qc3m?,FJ8VNp
UE8;636(?ClL61M:0i^9]C"b.]\Po>6SQV+Ds`M&D,'-=P.?A&T6P+H.T>@GYE9l>b`@C(dt8"bX
*nU5c!ft)YFgao^CTa?,?7b,a,G0WJhNWBVJ+5W'!M?=X'3X<%In)g*e1'.dt!L%[7U:..[D:mD0
j"0CFRo@_,\3sk+k^#F^D:852@NP`(k*=miH#.P1f"UW/DZm;2IK,M)A#nN2h;-`Yca%deUj%e_@
0;jh4U]0(aKfk)9gFO<+p8-3OZ>E#[o=_NhrE(uRt%a,JT+()ue4.`Za.o#t(hCBT:![^9Vn5Ojr
7Q'\-+/4JHT*BoZgh=6p"#f_khpT/Ve^[m)c,.Q.r`+2sFapAunCBUBWU2odZ\SB8H.iY8=oe]Za
GrcRaC_APchC0',bNCrO'@B]54%(RiQkAk#,oa:hqR+<oUWr5_AXN8\fNF*Oc6'Jf7*%3^[t[>i=
2m)"q.gLSP[r.3[3!n$=BrR\Ujdnjml*.@]UC')8,%DnLJRp*q=g-9q%&uY]s5bS&8tdj.+X"Jmi
!Fapr6h!+7SC>brFqLrF473\"E2Vo4<K-;$upQ`M9#%"A:BeZo*4#Q6YKN`0>I`[sGhm<<usE^f%
&DEVj6Sc89%e2cY,[:(aG\&$*sTj!]`qgUpZ*Yc>r_,/o!!S3<]b$U%ulHl][l*uJQ*7j!,Ar`#L
b]1^"1YCd9_Z'Me#O>[iW43M;>_YJ#m8Lpqtd^lRU0pp[%5<58JPR,^lM;$W[1:)iYh?O2+rr@!%
1$i@#Y=[mhLsKYh>I3UB4Q/=oD/Yo_)CK=&CGX\cerDC3FNWbR%"k*]Wg(8'Fecq_7P)86a(4oKW
7E[BNR55',<5Ef??tBUSXS^4USpR?XHag5TAE<(9*/J!CDui'NsLT8W]A%d<F0`;$u]Ack6Fm'[4
e(^,)5@B'<!92r\HU<E]UfI^`&PuUE?<W`F@G4C.32HFr6H>O5DJ?a/2o'S8&QorB&b6HHeiGN/a
>LeeItj5NaJp>qZ)GAE9N^X''6F,[q^r=A#U+0afX\XDZlVZ<i-BWE4IU^THor]j[.K&No"p6s%Q
o9(@?4,ar&sS?W@%Gg.WqHUgO_*9#^`ceB8ASaQ'sMCD9XO&*/oms)'/M"ko.;g5l*G=f".N'`N]
mtdsO44PB+KpPQ(a,t\\d8&D"ps&b(9#johX`naEMrME^C8m5PPX7d04Q`at\Od^GpVA=dh^&o=q
k)G=m3=-p)5.UW8*tjhD$@uH-DVqpS&Ob?htY5rQF2/qG.6:b&gV96hn5D_l?tm5Wuiq%i`R.#p$
4>5,eIa*9#;QF`\b;1C?s.^Y@shQrr?2VQ0E+bg%1[LV0:MiZY$+qRd,`<NskI:7QRWH%Gl&G;CH
h12u`mHMt2Ss_I'<!UC,&4A!@c/!$8F3HT9/;a#+FUC%@9H?`/4d='kEk+\QIoN$*Xh&"Vj>D'.U
l$dXdWnc8Vg<K:nPX,KFeA&UE7h^em<.a4_EC-7lZ9&<H#&,0+>,(5'P*5b9,7N1C-FfukU0GRUb
iFCRZWm,FhNS;e8i(fEAe*A&0(uH'tHjsL.8'#1.k8TDgX8qJoI!EfV"Y0<-$u`PKm?4K!c6(6'5
iWb:3Ks'9I`S8\b:ge+4,M5\YjUO_*jLMMWF:UomE5JZe4]fA)7:#5NSD-/IooI;p+MrI"a7N9ir
t`_a-h$hM,gu:,#GaT*"jgcDH8V3pf!,<)s7,^B#91*S5q4V45NtnGcgc2S$n#(hCTD+:tCsJ#LD
-R1KGJTZJ8B3prPocDF*4TNWj+V*0%m`>r_e5=B1mP;A$FZl>==iY20*XS?G^nQWuMO5,,($<):-
ak6ZIoWqj^5[&`$Mb]ME;BoQ";]6PAP-j,sYm4B'%fYi,*`ebMKl+oD$2lqHqdI.Cj5+D)#6jAk<
=DC^"[5*s1nRZXi0aT@#rr<mf,T$,VlR55XMXucp)G7)E^tuK=n08k"Xo^OX2lsV0d8cTb.9B8D9
=@>1@PLR@$au/=0iq-?YY25AXZ/jPbc(f-UuU<BZceA^W/0-P#+]SNS7kC(Yf]-lA5]ur]>\0Y-0
Hl]o`c^LRP-LRp4SECZ*`h_Z/lhW]7&\Y%qt%O],/NZJ;d[B;3[!X]<60pWhW;ZA$A"H7eq$KFOu
Lf*[R\_b"&@")f[OIcR=?e[B*gcSY[BJN,NTJ*)Te+!!Y=6G`5SiN/jM+nR`Y+@om$o5r\G>Ds`*
"f"QFP2aD-o]:2_Mm/l:].gL'DVd:l(fFbAIc3i7Ir8lj3>^i5TIh1!uo/7u/8Op@6>H6!*O%Q)X
htb.+caIr:7.nQ>F\'Z$LM;(\7>,0Fi]5-USrq*pI,sUU,N9kO^@BIWlP9P'e#NB%G]*H?]V/s*p
nq5@H]<aWAs(m]f?Q<"c=L+<W1t&_..ZI#&rlhOP^FhW1=+YT$hI[hHrd<NIWEhil6hWPZfu:DMo
lTiSV8%lm&8@opeQi<e(*KPc.[%YBptE?>/rrBds7c*5k+aC?3KtkI+K4kZ52YKb'"eWN\BL<DX7
:A0?(1C/`#aD@`3BMB@CX?V\2->7GV9@Kio&^3+ZcSL:`q2hje.6-M(&TG-(E)`ZVk\Qd/Z$e%`=
GZ9,)*PYNF17Z.*$@CLV(AVfa$=d-?pWn%;6EUJ;.dYPBGE8`[n2rqtg<*T9AOA8ZTIl3-4O$:%<
P^]*IpPIC#5+ILu!k7XD+Wm`-XtM@?l;ckH:hod/*cjl6G=D%*`?MDjY$A1F#**Xp20'>e7Kh?qh
^?'JHae+m_\,>Gp2QpEg,]YbP#:Yhf,s.\kRjS>/h,7D1^7QF3dMlnUfNYXa)7n(mhSc2Wr7gu!_
,+0[oh(6dqX?ah2cMU1kW3LGf[$C.%>^6f5]ZN+0-E92uYO,P`8=R:6^k[\l4Qd4Yt>BH?%A=fo-
[O#L7'`Ub#k+3$eW7Y<I!4T]5E>@L[R#O1&h#(e20<i_D)Wr8jY+./gD`FDHAB!)d+heg\./Mm)R
u7i_?@e;^&H]OlIhNqg1Fa+uA.7<*]`J3.[odKna#dO3^d+mHs,CZaM6p>'Joq7KpnpK&6'>R6VN
Wk9?<XuPGldmUlf$4Z<(;9sC*3J]i%]OhWHeupp51#nmr]n%"'?WVB<7mkR3;>NfW_Epi1IWIQG$
f,+g!*Q32hn3WEkc\*#3io&%GA?Khe;T%INt\sG_i</80<DE.&r)/TVJZ+uQ^T89XINenh.<eof>
l0+E(cEGN5&4O@qsVs;D3*uembWjI[V39qdSn9=#ian"Q"%D6CY(%rFqu]C7D@#c%i33KT_;H=%#
NeQ%8TF`X4]<G^O-.O5"oq_-H,YIe_F-0<E[)M&j1Y4<WqeAr3ZGRj5o7U)G3CnXf1*$j1LhJ?8j
@5`?mGL76d\o8J=EhIAuok9+*Sh<8jR(URs!'\Vpn,[MGodG__l7S*&LNf=>,M%W?[B`[*.!.!`K
(#?:7`nsMoaJEg(Z@l`ciTXJVOaH7.La."J70pcpFSob<'r5i2Wj_hQiI7AHi(e89=Iu&0#e%4'.
.;:cCb^GYg[?F%@r:qrpP"Omb!coe&EPj:+]p[)YDdAQQ),?h,\l+@hW\gq=;//f_LHgX$[k5/UE
;u*!6,p??c/*RUY9If!_BAXP[]u\8Q!dY>_6cp*j@=jh`AqI;*f*6fG2!#1?b3LHFhr>h14!pjT;
lrbeSA&BS]J"Zek2qqT<)bPp'Aoq=r4e<jAAe?>8dK(u7fNe(_(@37+@CN(0?VJ#lg$?)7f#S%)Z
HeT1K+PR7KbS$uDrjNt#*\*FqD[2\Xr"`\Ef:8]2J'V1eqX04O[$TP>[4\!ATK%8m':ERK?>):<X
WP98e2fYq%j_K1X`\bBDHjS(T?2^us/`-/dG7F]h08Jt):m'5%c9).TNd89F!GF8(PRc0NbaKR,d
%3#"8#Kp/jTb.m3#LQ%jS6MWh9^5&qP#fQ&WQ5b8KUf$mdXs:XdN[:cT!:>FgpBLW6,m]%+0sbUN
p&1ZEDCS#YOm,))O!`75_jG#B13iNlM)!gO=+]NXGN+Xf.k#g/6lGOLl%;>.^@d_T%WOJJ'm`^*=
;KY>011I(/^)q!meB]SI0._e-1<X7=1oDq4RSNS@]Yi'E(=#h<DW$B<pIFKl!InoK\\I(mH^m>Yq
$LTTZuSV;>OilpeL7eP&$<+K-EnD_;G=8$[d?3\+&p>6j9^+r;<)`Pft0bdQOXJlBn=c!#)mIsbc
5j:NZ8JO0cSO6se']XJ]KNm:-_TgIY:VL,mFNQbeSMe\?;l>%)]!"<fJ*JM'&UlIHM-NbRlQr8;L
MjTq_qRUi0j<4ThjeVs]Ri;C>p4q`NDoA3nS*qqUW:#c+XB4fGfG[FD-P%RfF;6UVG(N,8c`?jo8
?OC?ZR\bk@,6qVE7Ci.$n\gWVF9Ij+l1HaN-Bu/>@=a?];RVFAoW+F)Hja>c`M/5K3afV5\o!.g'
p<Hc0DHl4<&.&go3bR]Jt-1t.Ej45`APAg+&955=\Y)/G_d>2mk?/JF@FURVGslBQe'XV1'*^',f
1\EAp3ZLB69Wf1[<NG\9TX0!Ike=>:e`M6Lo:ghZ,^\>X\&9C6,YrQ^-bX*L2g)#>r:jdO%.4Y#'
T5DTm?3$-oNB!4MHH2-)'NP>=rrCT',`pf)Ht(q!;sfUb`I)`XO))uMna/0MW<&=BJAtA-T5<>M(
[>oW0j(^6mhUSWqtp&G#?=g_kNOVti<SUa4lkr>BC>d"X>2C\X5Upd[2)A]pPKD<2=l\[H37Jm4f
M)onY6l,o?ID%;b,Rh"KU!aac)NLmoW";kN=^(YIN]?aIM1/[sQ%;Y9mKi4QPLF[b!QU-'lMF;1!
n"eenKm]Wtab?=j\^h<,>=^M:_4^ReZTJ<n>:Ja%\D[&hh&0ckha&*EH2NjMl+pP;5B>HQtn-gI#
n__^NfWr8=O)1%@h5@qO/^2kddVC9+6;FbkRCC5JcEF,nrX5$a:=drDO!muN;fn:e\g\^$#Y8>Gq
D-SH"^?ik=BKrRLHWbO\iPsM*TJKc.Yp_i,PbVU"FBM10prsuLrAhDc#aPR\'?SOA0\3_:=XpeaZ
@o[bQg5a'C%I`LDgA']GbI[`S&K]%g@^`meuo6\OZ18@P]OtfQh'',S*KZL=[p/91a<sP,#X$G7i
WAPJd(IC=]$p,mD7<3q=\DP<u5djHEbrr_p'QQ`]U^[U\;l$VBoQT&Y_`-.6C0YA\*::5A]"bS2>
Nda6@Z$l\qo)X62H;h(#+2acAl;pr&)jUF6<3SD7Z*^KAHJ<HpI:0r!9Ydj?;H[T'+UlC=)FjLo&
c];Jbq(i?onf$0Y&7T>5i2g+`\/OAb25Fct?qa!SO)Fk)Z*0Q'_S?LZMDT(B,E.8\L?uCNfc"\.>
b.ZbOrnf\'"^:e47X#[Fq9,Jkh_)FYNt$SaK1?k7d+(1u!PrZd@_L=2^@IfJeR>D@HH;q@o\V="1
h8EJrqQ%60#R(7(tmqpR\^)SgQ[Xr=u1?QlNo%"V^tH`>_a"o(rj$gE3b@)<YnU`cb-i>"7^u^;M
U%f#Sc.c=M=.FblJte'CV><rrC.8_`]nbXD!BBhE[p-*\h2&p[aVB\K5d4,q/hAAq8ZTf9ZIRHDc
^=1>\(;W*1[s3<nWTqd:T'B/P9X=HpQ;nLE@Uc6uRI,5kTL*>`EurSo2pj&U/Pl7UA1@7lfZ\Dn>
Z%Got`1!\o>eRM"N(Qs;G)mM(ij7QcLCF\Wg2-Tuu-Rc[?:UVbR2o]"s?",O!/eBn9H,Znu\BtB'
.pFetUJ:2*_6PDHa[/Kai@Q^P_`4Vb\!0QK\A/i+^K=t]OT"VeZ[C:?d6=4$%gD[(.fs+>"ue>B]
.f@^UQA0+g*c=NZJ'Hh\\fTX07ttkC!_.J"<08,keeUX#&8B%3?2^NPNODC.iiVt.AM*J'(/7RYW
Y.#qdeJ!;@?[Im3QOCZB[q`=fO"T\,+8NLrYY7(p^:TS9+p"0/A\4VtP*sV7nZr.+@n4'`:gc(Dr
4I#;]>AN?0\7iuZ@9<5:r[+!s&gYjp"A7*mi5GR46r2OR$lmKW/\8T,W&#.V.f05^$A``a;07GsT
9*u_rXS,hrOS[j=a(;F['d@YPD@mkJ=i5u[&e&=P!:TM"m^RVNV-Ehd8QF$p(9eeXrH$4,pS$SU,
<QdmP>0MU>Q474rkuhoJF'l@0*T%js'ZSFIfG!jf?8RK9)7SDS^DQ<!O/fqhd'2;@DRK)'M]UHXF
-I0hhqG<Q>LtO@C114rGk6N+HQ"D[:=`XQ^T7FcpAY,<@,6AQC>;4j#6n.dP3T(`EK_?Dh?(r3%5
`k/Ma$U+*BtCb>Il9?DlDpO\D?8q$P)1N2pO2oZECn!#Ij>jNu4b^r+$H*:i,#KL%gB]qp'd@F)e
"ol<e8Zd^bJ]UcXJsqCm(%<j/%"-ad9"7O&MW:Dd&;S&M3b_N-jZQ/Er@GcVO9[]&8iF%Cng[4_9
p#u2JflR?':5!PQ@$!,8:[p^.g[_1W4@WQ+0p`gl[D-UKGqJtp7P"[JcJ5E9&;Kg7+Ges3*8\MWc
Bc.]P(.R_n5J2.r?$!*sG8HK=7rn3SMT"-()2n/lV:Mt@BrCdNd%j;l^GUplFd>L\=FaWLbghOj2
nS?4k-;];m_!&EVh4Gg<(jrF56W=7NlsO%5j\k^mhlI29&`6cBp5"4`(nmQCW:87LmC0bA#P4`!;
g4SAaJrUhGGOm9(qk-J48&e=.HpbF[c)D=T8B3CE!=s`SZ=)#!LLAhE2A[D[+<I]S8BO`K(#upgR
-Jo8bQQq`)F@<;uq9EWa9r?@FVYh*\A*C$pCMB41`>`_O&N-17oj?Kio?^JO>CBA4^Y]ZmiN[9qK
%h(eJJA$N',J`FS*'9kgL]'&)dY"eorN3?(8C>IgOON^]R*U;Cq.$<=_nDeEP_l[A7qREjK,$;=p
N$uh0e@_-4RARatnRiV%<F/c?I5Zi9j!VIi4gdak,ZD4cdk(K]Y?7?]]rbiJ<iS.-NAff5>ge8F5
AX4S?%'&R"dN'H0Nh>jJKXBHSC-T+V)"0cD-,W?3a>^^Anqtb_d@k4H;VEuFd8)3>5RIH`3eR];&
P=)S=#jT;S@(@$ubNsO;uerbZ0Z^V]qq$IBiWgR\*TOrFKLBjG,eZ%&nSe\&%)brr<s2Pc7CFZXJ
V[PA'[bhJBN"J3!:el.7pNn,@@`+?nKAqV$J32W\Eb:/#E=a,TuDRq5<5\8'6XMfR0.NO_kNEe73
@*jU9m\.!B/`a'#jWI:Wu))8MtF$a1DVTlX"l;?8,9V06:2JB:aG_:/iMAIe9(N*Cb.0uu4jb=Rq
Ut.@q2X`=r]),mmN7Vlr/<HuTnW\uZA$0c?,\3?[kM@"[\ka0iiC>bd.7Q8NZQqJUdWs!#(iiuR?
R:n.>.ue9XAj(4,Ll'p-SG!HOU4*&<LKFO@@LbNojbmcYcO4/4)!d4AOqJTWN:SH\W=_tbfj8%oC
!S0CUr8gh5aW'XP+DDSpfR^W/ih%)n198bI4Is1bEEINu3DQ./n5+hV.!.+=k@s`f_[,/'%%:!qP
<G\(2LXBX`qh@2!a(Mb>DDb$WD"'a$n]?Y.1.=OW^%4+<YpK3;<)699%T0fppfI(=AS%\uU_>O1,
_FNWbs3&f[Jo_ZO*58c$W4ZGnIjP,jb\L@F#Wn__;Y"D')b?h4_p/#sVIbCFtHd06Fc^`AsfrN7L
G=tC&j"s)fE5@c6/3fY$VuHae!A/QplQ(dB9M$=)^Y^t25F`r@\+OMj;cV/F,\d+&3MLpm*+]*PN
]guAa`uFG,B*@o`DSh`gnJNTYmW''Q9$Vn`CY7sRcM?aIsISiA>lIJi1aF/B(:Ftf@q4C5N`<EB5
M"B2eUT.ET\@SVRGW7l)MCj9Q0XtKRfLAI,g322o`+ad4\f[-em=Y^_YV?q8Y'`d/1W3/$0:H&G:
pG6ehK"a.I`G[[^ePC/SC'c#O^jHZ`%ne")9K*\jq%+5;B%BVC^gae:4n>'"kAWA1B)oeuT]h\Du
p`=9^kc>YUHX#+S[k7;LMSo>6&RXBFs;_',[XDZaSD]!g.Po@9M?.Ur<EJs4aj,DFm_i%g^0'U%q
Qk?$=Gg-lD$pX3!#_[VRk.Ut`p5!S)\%?PV@nK0Vr*kl#K96QKY?nnQfK,#VDOe%5<IXHd'3M7TV
XFoHfN:=Wjfj"a'uO[N-(u`!n"1G*''\QcVOcbbVkgFYEWf[WVi&NI6-n:(0Il@sNl+o`jN8*+Q<
dmn[7t.4ctTQa8i;k@U$#RHXa5D(P\g2+om4XE\Bn9a=;-o?ig-V\+NJ\9=]Dnkm,7Vi5>bER^U?
j9FV^sV`p`b!4l^Z0?`.3hMelpJGq/o3=&jLaV=7ru,\Ur])MTJP]<$P`\O8F6@RDu`F!5a.h^10
f$MM#<*MeQsF*.I>q<&/qX&ZoZdu)/6b7Qq3g*\IYnnHo(A'S,WUjZi*>2L-[E(h!^20P.,='9\H
l.-/_d!`OCTEDq8K=4$O=ts-7*uGhV;`1M<kiMDo3;G;]h/DQR0YKP30`hNCTn*2dp@unc.#XkDK
j*N^Z!^f'jL+ZJM'QS1[?OC6>#/71$6"-S18JA']!#5AUf?A)YjN0PR(Z%R6f4sN8:-\d[uE\bYr
29uSKj4#pri^mnQO(g:EYAt7Pn[f&WIT5E^Gp!:Desi""Kl-#>tZ1DI;t[iV<EB^F;uO0r(1C3\6
WMo955W=1Z2m^^DUgifUemQ]OXi`ktC2(,43`f,<([[CoeG:8B_rZWk5#:+VrS2s%9OpQat7a\-E
iS>4Q+h:=MYF2bOkWq2F>.Vc.=e;r.^LQ%W7]8^11P\)klqdR]2YBTZ+OVJ/aT21W:=OGD`UN3a_
ilH%-=0^_'Z.2lI';s;*qOM#@HdaJ^=a.S"<[/g]=age!gG`5tJF9qi63L3!qopofQ[IDD"GAUH^
M\&'$anBW*[Mk'=FVsX2D/`4,7`)O-W]Kff,'c`gkbq"Ee45EO^."a/srN[1PaIaPj+7`H):slN"
Us"`F,Ncn]jsfC<$\qK_o?[U(S5!MqFQkmlD<(NmAS,V1#XfeOfUK\&P^Z?O-_$DT/du[hl"\J35
Hj*_RL6!f:jaN'hq#,\-`HZREbIe-iJ;3HX.ikf)t$g[UaPU@uhNiq$V`MiL>SWli3(E?KRJShAa
BEN5gFG4-IXY,KI^%oT0g3dm]LjhIUQnY!@9H`<uGDqc'#>iCYJVF?VBV%>3SE5a]V?f!8rNDen>
-Fi`"=R<cW_c#,GkLL-5ZrrhR%5'6+p4beY&.^c<bU2BEDMm"[8Bn/uQi@'"DO;m.oDWc(_qn;1/
`Xi0'+_C2ekfl>DoP4<V<<eGqH?RXW=I!XXnp9`U;i9:Di@HVKe;V\d@\hsZ72C!ViuL<)hae<IH
RJ6nAD3Y;bgIn`FMe(D_=r'*K"JoIrP-@?F*6'!-#?*9N^T:NL98$4g"#r&?Cdk0sH&f[)GeDEMn
ajU.SI(\BKf0(9a6GR2-+Y.8=,]8IA2[^"lVHgF=`B[s6<bG,'3tjh-fqAgO/:lI'Fk=WNXRkLNM
1CS+bADl9r8`c>ocQ>^lD339I>7kK"f*$UNJ0^5:[nX>Qg1(Vj;I]to#eDcHE-b\GA5d--FGsugJ
p;E+XB]B9pD9EI=/MHU)6V+6?&VB7PUjk\*ed4kZ]"h>tMlR6ug?ktU]A<DCMe0qh0Fn\_X;pa0!
#_3\%#eRBEeWmYgaJ)\oo/UIT,96H$b^Hs'4SA6N@p+V_'7$OSCV]MXH7$>T2&ok-M(m`rrB\f,U
4)KU2mccqGtr&)BW&X?_#JiR5BsAh)QF6/!)<hqQ/*H7rrQF20->M^/HGMJXG$XMFmY`Onu."3+^
?D!)X[`EHYESGfD2:?]uB/<EdHZNG>[$^)rkV0B3+L<a[4fmerXiG`LRZBY7NNYFM=qg/[OXU#)g
\qh)qZAhcDKZ9\trL$%.[is5hR^2a/Q04i1r'ogngqrQ>GA%B;.A$cVo$EkH+4A1Z@5"l*Fm^sB=
2;@T_*Nn/XbXAN1nqfH.)E!04Y8?oU?Cb#;McEJYhZ>X@GQ.Zg&g_L4q8s-jq;V]qEqeTF&nAtSg
C;*d$@c1V]s0\pr8Ecc(%2!.#g#Pn^Pp@/f!47aoo%)\c4lhn9!YXZB_JM!nn+n-Du=*4[QK#B3K
Z#WXtm[ImFR;'ek2gREIj%jl\nYC>ns3m2_bomk)\bTGcZ>JAcf=SHGffM0$QCW/'W'd[f"Zo!sE
B=!7!-*(,t&fR`]f4FC1KUm\qM[ij+8R]9T];N?dmqa-\K^nQ;L>Hi6s#'<IaY3l*T2^A.G^+ecO
B=hAp'?gf+rRt=t;?>.[eo8HQl+au-1)YRrL!kl&9HFh,\U$%@LENqHS1"R=(^bD)$],fp(6RQ!!
rE[Iu7NThrmu4Q2oZLi.^r@p,MD:Jr4/20FCd[Vu@TWJ3]Y`P8-TLek=^&ZEYDbacX'=F%:$psLU
(f[r==,,UF8l6K=\t9=(.WTuM,KDbPjMThi?Gbe/'9lA-!e84(,Y=(j!r"&YDt7`&`IW<aeHJHCe
CgNLl!KHAl'WM!4pqD`L,3O3&*h!+rrR5k6FkMWY`DkEf<QiRSI*")HM/i;3=,0lD`"]((NGc&%8
]-M0e2=Jj[&ZfnsmH+]Ks5/A$mk%Q'Q4f#%P$=NC'tIuLp"'ar!S4,mBt;::]kD,fTOWerT5_4eM
)]P=^mmsiN:!X04ci@oFCNDR.+.=NtjPN9o]&r@P$i5;en]X8-ec/=]j8ufS8SQ*3RlFGq$/sUQ$
3RW\$3?!7Rc;[Q=YPBm.Ec#jT-\D!OfGD'R3Aj^;l!<o+Bu+=k.coALWk6'/N2q<FKHVc8rrC=mI
=?CcF@nbh>1X60,@u/DTmn]UEK)TC5ASF6oW,tgMeV53/u,!lg!+fT1?l_9HH.33=3imgF'9kF,.
.9X]s[H<ZCOU(o8OBC"EetD"T&24m/A_$8l"[O!1`8>douOg-e+e,2)-F@c#6b!W:JeCEN*76RpE
X?`2e\fhplT0'=F\f6]]a.A4lt7dWmjtG(f8Dpr2M`OAU>U``.Z)oA]u4h'6?2.`X:G=`s$L[brH
uc/G'bNW&BH]jDZ#ag`Teb?)`Sfg8uM^Lh;GEMitB=#g1R^MI#@ZZtMCN9:tY!.bhWXlHbUQCm=f
$<%AAM5/og:=YP'\AIc^`=Ed*3P!LELPrKV+4]M_LhNj%+r5(cVmj#sE:/X5ppUeJef.Oeih1L@[
Jp4EXfcR5*92+OUC,?abK$]kIFdp4lUJD/CA>r/N+H47+/_WtLgAD.Nb`sE=VO[t4Sm2T$E3t)34
_`Fm>X1o54ARdD/o0%$+Bb<:*\TY0tDP)k9!&'MQ3f;\ATh'q8@XJ<M:h\+W(F5`[7`-gAC0LEM7
fs:@NG^i_aj"^J64+P)f,3iA@WI>Rp\n@2^TB<K(t+Wm,LpF'+kF45&pVAOn$'B$/JX5-PEtk-F]
9mKi19(9+FH3)Ag0`u.Meo1f9P;*08t*DH_/PUJCjM_hm*Q[(n+8CrsFc;?$o[SR7Qn*_,#(!V,7
#i/j+3'lrH2d!keDd6`#FmjT@Cac4MK21`WH]d'd(gc'.rr>6%"5>&Z8R,`0a!Iae0rdG56EqiU?
8kR*]?IL8.K!ckrMUNo%12>/P5b=&6B,#]#ln1n<Lq0"brBr!p#d6f>1m5%5aOBs6'P+6^Gm]Mab
?cSq\WLYCFU-t8j(O.4B$$tifDR+L,BS+D;eYmR\f%"gV\VX%2kLfGf1PqMmF1(YCssp6h?_,dp=
e_m?/NpDX5oeBU/"l4DXZ3#O*_KhJX<X_',TR!.;\$<j4?R'lT&K68ac,hsHgV/gGn'(M2nFmIQQ
O7*qsiqSbOCD;+j1W%X_!`b+Y__bF1U2NZjp.Mo'%GkK&tS:6l(@?/V(>i?%')*h]6lfSnpP4GkL
7(UD7C\3mV7OK+ug#I?MNEN>.:V2rj#_\t#2oZg5]9=:5/`Qa`4[5h!BUnc8Y9jHQCqlLE$u??BU
0Su3N!4m7r,6gmFj1aoMO"(t.3%;'Bn6o\CG]L(+;^9'l.?g4#tN[i(KB<68h.+,-OEqJFCc+1of
f[9ibQQ<dG"B3`GV1MX^,6G2LK,mrrC>]jA;K]e&ELL3NR"!mCY1TZRQOReMm3H.%Rhimi43^SC?
_Xc#0hCasUlnS:'>7\M4R93D"Ep3qISR6Z'P2R&QW2,oF\f]pMpFEIW@M6t9\Tfqg0a\B$&n@A2%
UJKF78^DQrZ0C=`)`AG0I&F"aMkFBtV0qnS,b,mLGkMla5Xo)`DSt-3ho.bTn)EBP3*Lp-q2ZEbP
di+UAoWl2:eiZg(\*j_j*=%O<YPkLSLeI#/VS@L(;!sT[!\ARqh?F1VlN%9XGMDl`Jl-$sALNn)@
H5<Dd2>AiMl.-r`grW2fL10.'PX41/r]tE"bp^17Z\$5,t0>sBmuWEIPPStB@*+eed&,F\)I"G%E
lf3$>`i1)C7A)[#JiX3^QsIo"@SFkf&Kj(l12O;7f7!m,V8lW[jWG]/PB5LrG2m[U'H+Xc9i;BmY
eP8%AjnEK/]g>7NGH;d`3;ecI=WrrDq-DmiWsCW9fd?^u`K7;+Ge7)r1]\`d9O9rm[iaeE(K.j*r
o1D00S)/-l8+D.-/X4`ZdF0GsX"`]F<GL,/-W:$d-O_C?D[:IZ7VYjmfW;4T[W?YN.Lph6YBDBk%
$`Zi?AjabU1#j5-Y?3+"[6)*D8-rkNM"@&im4-^mel$J>`;A\M'?Ro:5pC=;"8'aPEebI3=te/Bb
%Q%BNd2=Fl#S%,rr<pNa'Uhe%dikE7/K6IS&H/lONNfgI,`'KoT'll)6/%#^Gq&*2ihuU'!TLon!
q?.P^YX3P\qXIBu]pH(UI1m`\gUY]mO9]'7qqG1bNnDG=,$,mjB5WG!9rm+Xk?.ed9<P9dd$D$<l
VOdQX3m;9,#g*3:EHEV$7Dkt0thprfp=X]HhQM'=k0%WgpE!\iWMU8p1MOZ$3So4Fi8jMpKo.2CB
@4aY50g8s<pC7Yc%gn[iga3H5L0>W@+/i_$AqoonH>Ld$c\'GpZE7oJ-/'%Yi35+89O*QMkV]_hh
c=jY5egf14b[;4aqT=aIrp'*6bLj<Qj8-ukMD>G$e"Od?G'Q9!@"]2&389LaT^)!(CQ,'5)\7Q87
18/41eCCNlZk3Z%cTc_mDk<b"dn*ogo74&]OdHh6*.SbkNTsMC:2%BIk=P]\pW[2:*b+%GOHKA=$
C\.):ptFItZB5djgpeZKYrZm'uKP#b92i@KsB_SDs>sY7Bp(ir]ij\]]NXp2G2oZM/Uj%1.1?%2"
M0AaKa\o[#RoG9iDXDc6+ngn[s8pVoWH?.so.j%lWk_Ohl(=WH[S3@e5bK]'M/XGb-!9W<)oa%Qd
=!:c39<r8G.jZnJ#4=78%6A1XoNHrNur+)4D?W"!Dd?A]ZH-83CPpQG1n!d3sC9%NkXl27g7HL9I
D&`@]hgenMmdjCi9<VUY!8eMn6Bo+=o.dFKmEO:gVI.>"*iaCe#J;=#eRG7@2NRk<GE6QHg9A\r&
gM,K`]ffUD>M*)odPZQHujutBe2aNH]?X9Gs2XC'I`WWFTP`n^S`!K[D>82gMj.G(K]oTTj"uJ:#
UlR]I3!]*1'`l6*%\k%P3@;X5b7+>2/K_j?@MoZlaNlQ]0$)UGEUbIWAV"7[WYa$$,rRp8\'\#t!
ae?]:[);2Ifg#E-3!HFggH&`:6W6C50Pj2M<<k"aNKk:kMBYcT*4]_LM5,BcSK>dnf[SQ*)0ooM#
*o\#$*JG$!`ZA/=GJ<2nO,ukb9Xu.$a=OLUa$FZ4T[&/k(Ti*.M$MU!j?!"*Y%aU3+/?&GWPUp;t
79B)5NHYW&.3%Z_pru0$=q\Bsr8s,O^#?/pq&D$0a;2LI./RQ9E8Fb@*_]-/i$D7^_R*Ze?8VuXJ
gcqRn9W/5<(U0]J_qI-8um_$<)ikd'>(\1=SKtNpq\B3mdS)lpsqJ//h/3m@,lD&<q))?*=&[jpH
2=EJYLi/E6?87(h,3&iQXTcr>+K&*F0^B%uSqU#*%1H4?C2q"I,4]R@OhqmD`nGACiq'#dKc;p;N
,BQ8qPgdJJ4,AaJRo82Tg^`pM;WZ(nqfn+a^_4_B-L8g,\9ZO_1E?g,BUP)\&]S-O7]om&)-Vqp(
7CW0d,>P!Oih\u0>>^#pa:gttkaaib!hpgUHYI=+.]Z@Ykq7tE.:P%(.+Dp2jDGp8gq908@g@)q0
&#;/Y_""=*66K0-N#ekDCB>l+h\*SE\)Z%@[euJ;B;<q7#r^2Dc+(&pL13N14CW_<-mE$U3[G=24
8<hE5uk"tR*Qg_]R43T[FHN!@3Rir[ob+6#O6<%6A`.PNGRc7A;^t)>]qt`_!`X7Rs.rs&t"0$O*
IsU'YbMX\CpcWi$(7-p6i$K%lSj[@SE5mQnBY29#[*r>b5mM2U0o]+'?M2X2,X&gg+p$H8!;[nn<
Xn1]"It@n%7q?(BqM-.OF42b:LTZi"<7#":S/iTYQj#Fr9b/9KG>ah8GpL-L/J09e5i+/W_9<L?h
JV.!i5m5u7Mksa?=B2l^V7:`;<dmqr68UW^c]U4#13pkI=o8^b#dofjfo7B5&_c^Zig8:o-U6Mo5
cLX*"84Jqa2nX%&P)65ej3@nl=[?9.fYlSPmk3'+O@[**RQdd\r+>:+$]c/'+;B*>0.>7dKYXW<$
Is&#8%F`;F+?Z$`gb4@02gm=--nHo,c]``ei]]mqS\Uu1@t^@`V5`Orr<;*Ah&RYn!nW#nFpJe$R
<+#(1*u_][jgLP)>lJ\mf=gq":$^EJQ>"efcBjf)9+b78"Qoaj98*=-dNP^YZWadPC"^^mHnU][d
6p&HcqLj3HZo2rca4<7JS@<B75-igtSZfRuYXrESGVgi[i(.!-NY2?D&?R^"<];@)t3I/(sQN./,
[o:$B3GZ!F*AnS,.ClGr_dDb-9S`_uh/5,6a;D,_#k8eX6P_?EHN'aX^URVDNYt7r2%__"dG@FTN
5HAGM12W'j`eMZA+bZ"9r8ZH"j,B(^C%7GUd&l4FZf;LRWmnh<[H!do''K#/R&<YKH8<hPbO6+A'
uf4PM("_9\B>%H$iu3oDHisaPqqbH^f3\5?Oeg?m:hh#XFF(08>4l<nr\s?-iX0oN+K&KKRL7mmV
'3/ACE!dg,]i:@omLgpKk-<lB@KQ(KteMR1d..Gfd"HLf+C'2Y=Y,pr"Ecg>Cf-*p34oB"tFr<OL
iiUK6dWa[`B*2gEd:#,Nl^Qu;-7K>/77J!YZ4lT*gq.39r;ESdr>MH?loi'eiC\8tD5UaHEALs*H
o`nX`(?.%[VZh6nQ<UU\Tb0?944-EQoqdAgYX%U.e`R@%FY*qM!`glH/G%Q]9;D-sYGS";200,tj
=06o>.hZ-WDq7N)`0+Elm?`?b']Vr'7TqhXc8bQIfqSr/rmul#H#J2P5(Ud^2o[[td9*Y_pF92$-
Q&e1nnEt-l\rjZ%XmF8f`U9ji;E#1Uj!ubEpQeO2.(`62Kr!T35:jqPXG)Qd?J)J]iJcrcj8q<CF
u.q@f:`mp>k<rP^kpSXM#gQ^,XA).ld-_LK%isEMY*>UPdT1<+F`;V[EUBEQ.!=p!0t1l"g;!G:Y
ugR*CdB.l,8ra#SpbFgELE)WWOW&XN@QQK!b]`Z":4CM$:>'YDbF[a$RdGf*ulDXu//hEQr3V+TG
?1#5jo4Q%Y/?;290k=-I14Q2Wd)tSs5lDebNh6FFZgD6cm]P0-]T-c#&f$&[X/SV;(hqoE7/\%Rh
mE'(%"62/%<VOLBf8L@<p$XNkS>$WRDAA9C(Kb#+[7@mcTbY3@+H]m+216m?hK3h!!`m'PW'D1@E
:@U%C5P!q:S'>gS>NZ%KRP^KrB0\)X_rn:;:E_$&r@oM1lh!-_b7AYn4a)bYXjS(4#4/t^Kb5_Hb
/'.<[,W@rn,W_J]6Xhq;_1g:sr_*.>c9JeR0Wp2LOUJ=)U$PNAg*mC19'n-"cOACO&N6T=-[sC#!
!%`8R^;.sO;j&)SCe"7I")rB%c=cmkr<-7JT<iC:uQ"a9PH\%?!`\q<^n0?tK<-F7mF\Iu-#)V=A
105)`'=Mf`2qa4f;DGdL0BPjnZGh5g_=.((E=N:DBo--jgHp$AH:>PkU(r0-3LoWl2!)[$q`bs2e
.k!c_.tY@gibF"V8m"?aY,aE*X2Yi%h*,m@q8a(bEhDM$8"<tGO/od<=.J0542*TM@oj"ZjiW;Pc
@CI3`GYV566nj?q7P2B#(lDPH<J'FRVF,Q:3)ulAQPW`V;HG\qdek6]6dU+B5+4gE0("o.3W/!j5
V5n?)ri-]t:0pl3nK9D!g5@2/SA/ef*A9pR1kg-OXft5aCp,138*'G23VY<@k/^bXZtTh_U_M7q6
j*_+/EZI.;qgQQpa]Di;9mXHl`;SG;m!mYmoJm]J7tkiDNXK2Smb9?OP2C_WmV=T8AbD&1IX*Za(
MX.]H2FCeBL?3f=rp96b1,RW*8"D=\Mcg?_AS>@jiNu9BRG(BC.AZ!PUQ6!Z;S@F`Eg,7'n-oUNi
on0a;hn<UC%8o?Ml=f+WdYMm;in28`?/lNS%HJ1h9K_<rMN-0JQG#:=+'",*)Bce2=Tkn;IH/IM*
(/cB>FG>eJ42ShE:jLi-]T0\>*Nme1lL.,o0[W[??LCG\+]]E^e2Aqn('6TGPbr<1$u?m&?YD.^,
fB/@aPeN7Nk^#Q5VB]A]e"Vprb=4^6di`7O]TU^`>Tg\AMOTiu0+7HM<C;mC)T,&jn0_e#Zk4mep
]B=jF<n(9Y'a?n3g#4I(`TMoB0[PCcV5Mg4#D9'_EV[afBfN?MN6hsl[6PL#:lPY$]b%fL!IIZa5
63<k<(%ldeD*]2G?W($ePr3pkE_`l_tqR2^Y/I6DpWH43e&r`(_^"-fG#VU3]e/>TfU.THDh0@25
P2Q-DGEd]d)$X-m=/Pe33Cd8Mqmjc3qJco2`eC$]f_`gL#]"\4\M[toQ^<0+qItpV59/gfmAn'Q)
KE,#XJ!RbUj]N-q!lU:MK@&E_;bRN7_-DJOtSF@9bB<rZ(iPVmcXLU`M2o>Gl9(Rk@qMMJh079Ir
#O1.E=6>'^Bg&j4nG1($g8aamkQa[cZL-rNUKW3kA5!D2iOhT5'%\nZ:l+\VPY\5>LY-mTb-M(F5
1B/Juk"0nV(0YB_P!cJ/4U=hV4CrT)MJ@m8R()oa2ZD,5Mkk9h$grNet5Fg&87G-^=XJ=]cBS(&B
Jf8mWtG6?[@DG>rD<&A5"mbRBa]FjGNl>A`B$MpdX^qX3r1`>T6W+gN4OHu5MDgjL7l5NNX]J$a@
qduT[m$mF:^NbhUC9)6Y9"V9,`b6<tO,V5\h/^(3_-':'3CLL#GGEt-B/c,@&<OV5)p9nqZ!O%u[
t_q^A0ZH7j:B1E[Sbq7hG:-$AM,C7(qrS%[2PBm^1MB@9&I'sIU`r>5C@q@NP+?&`?rTd)p#CkDG
S70T@=3>Y>,rIPY\2^rN\tkVB?OCo3%hhQZofVB+]b`?T>OKBVN6:E)XFf,YN0k5,g-;Cqn7b.5l
`Uqo@q0D3,J3!'-k6Buo\iR"T"hT6iq*Lp$Ek(sq/6MJ[c_J5@9>>9)%%Fa^[rj,fTbpT.gdf/8a
'H>U[oL@U2VNS[3(,"+S]d?]hhV_fGN.'n0:Mr'4l,H3$H8F`'MC`6PlI\Rp;n]?HP?+(96h)LRT
arl#J&hPYKE`fWQo,V2E":7+bmjc0skh[e<?QO/DD7ufAF6tF*m^^3W?//NMM;"jkF]k4kh=huJ3
60NOPFPsNUD@dAb!No)CM%I!%<kJ*(lWfUfpWhnmC2*]S/5<l"gWZi^bI-L^+D(m0mL>;UV"^4X$
*c@idea%N.L"kP']eZqoae0#A@,IZnYO.@h$K#0/o79)ls3p-QgVgVPWG`YFId2og.-ZmhA3=S=&
]Xn3SAmT5:H/-D0`0Oc/SG@_0D1IV]Dn)tJ4n=XSV&3:P<qB,Bk9)GP485\_Z(M..L>pQ$he8+E#
o4F+mM,-<_QY$80f;=,Wqai7uY/htbTR*[TOD;`:0'b^LbOnss`(sI(HKM4d^k%2AQLUNf5XmY_e
`_e0g.E2;e(h6g`\G'5-!&tO=fhe3[UZ1gF5K5=I4CbPfI5ld-0fS7ul:;H[Nlo0sR^]kVje4']\
\f5@o7ffFitaSuA/RI^iri;hhE5?oU?nt+a/3tCQ/@3@R"8;3H"b6r(^fkd4[?1&JW"\E1Mpl;%]
!N<du$Mu-Zd%JUa3&$Z'2IOMAXj\6_jI?,p#5.l9YYtRiU?&bQZldA@fHDE:2jP:F9Q,D^82MV=_
bNL$Kbe,B`qDE!_HqZFY3XM,;FLF'i`a.I.lrGiKI:31m)DcC3#]O>IbXaij>i;tfgB-':j%:>H&
4Hd))['(/EZqdDF`hD>5?0@>FgU_M^F[c&@J*?=mo-O.[THtF2+Uf>+SB;C[XH1kd_90?Ad,@I(V
]fCpb9e3nhV1+t`r8LC2?`BAU4l^;:Zhe1WX!GHJ'=@Yq$2#MdT"Zp*>mcJACO&(Ip!klhSC*00H
5BYPVQgE*SGq2e8hPX!D7%Cjf5M\[S4O_u\PXI=Q(48g[_GV>9dTSp/l/k.fA0XQS3N7jbBJd\&X
-IkpbSX-8BNRUn&Ul]*SN5X5h*5(^DWL^ff"-Q@`Sc5?<EKo,J0Osc9*MAS`JfL7<JH;`m:MR?Lj
ej?@sUDjY%Hr'$khfdi]0tHl4=g43P510Q=M?nDbBMXicjFJBQV`UV"aoe&CVd\l.Og-hTNE]k\:
LWiuY*X8Xiq6i^1Eil3ljenUs]p+_A;:Y%)VEkGj%h8o-Xit$+M)QA43HQ#C2Ief_<4T>E$*jlU:
YD4r`P+I*`6%AcV%;X(ElL,GthJ%Du6HYX$Gq_jG>[E@I$$U;*R;=9,GJI@eGe5X=_s1p[1l/<DR
nQKCkDZ/j//u?_')G=<hE/=ap-,_f'",4;k9'm3?Z-e?D0G':S"VQ.b$@Io61^W8QT3cjC&3t^^&
#hC7T'R87<+*n?IqJsZaXLQ3G0a'0%03YfhRqNm(](DY>OeogLaNbC9Fj=O.<CJf>_ing8bp5ara
KXN0mT^`O`/%13f@`?_mFW)nYLR`+[PY5)[E'QF*W:M`"ke*P<Qq!#T0`EOGD)TY8TTIb=TV9!.^
)@MR3Kd(#O[0oHLkgQ[Z^h`$A*A,Ki6\AX(tPLu2)X>#pA]]<3V4;dU+2On/1S_IC^Zn*fI3#2D^
8R2YcS>^]0)LaB'qY9(R>1;X`Y'4NJabulNgA/t-p?T2[IP<P3Auip&jF@?FQl(bWS^-(T)6,U[f
6O*lMpS&CjXh&*`J0.Crr=QYHJWQ"\%?g)N<JVXi*\uP=-X<FP\A[,0k2Xh'uk-Ym(jOf1mZs*(=
&J_b&j+%IWBg,MgPH/_@\7tmkq)(e@+ZV65N_FLU]8ePNNCGXF]7FRC-uqlfe\h&bDKg?+mNnS-T
b7ho5o#gF[u%,+<gR-T.>a+f1*!UYd\J??6_1eC*lUpk)5g.K4=hmVdm3-0VSbXJ^a[Hr`e8QJb8
_22)Bdgg\8p<c>uU')L,tq9:Zmf_0,AhCifqlX(NIUJ?h3@+PYp]E&8B/Za6AZ4@GhRZ#-"GOq9b
+D^_kQt.a2SGL7iFHT>0MXk43W;N:E;tIVGe#NmkMcgT5j"WTjjRn6<k!9W\AisQ-4Duc;RTj!9i
rV9B1I-pr<kRC%=kmpJQs\%BK0tYg-(J>4,`^UMI,bgMGCH%^7#=YqO/,O!@I4oE33j-5ekr3H76
WLumuI.1H(u%c(X6Kaq]`niil2_EH.G)^*7C+<HHC_ApUnmu.]_EDpPBZRR8eZ:>m>>Ob?c\6i#d
:%0J[N&T/jE@Bm_lCeZ"G8.NNl\GGT`K\s&(JimZN)9bh4S9A,M?%fMZFVB*)@DjTW.,^0f9o(gE
7,!_d@`aHQ?16>M*J#lj9\UO?39Ej8,=:;T@WR=#fAMn>_1gON!I6R.7Xd=Xr09:EaF>ARJ(a8'D
KqJAFUb<6NEFhNT\Xhi'JKXC/]cJqXG>G;JBkqL^#>X[<30bmmWVs'!Eb:T?dq>c<.)#>cEp.u!g
E#67]2JBag+062kiXFnrn"X9;Dgo<_#p$bN$chh%r&`Vfc0=A3\Mm729%"?\I#@B-mt.i>P_"4A$
0t70Kl;R#Xm'1*+OT?q8W@oCGreJT5:Gt+f`<j5rmsD>[?*C$X2+Y.7GH%hItSDQ"#olQ;BRCn?7
eu)J`Qt2UQ\X_?Ca=n?\;M\m)XfPa>'Q`T%5TbCPbqobu,"b?ZC10St,GJk395:)mn3g&D&d9^J:
sPO=s0JKC<ImCXGd4,=tPZA](@f%1,s7sF*g#.$GBDPtlGdF5E@VgZ&/`Xg8'?FAhM*r)=*j!kg0
cal"l][H%m4:@E39Ob&X\-bsmp-d)/ls`H]>Kl`ZP\>*dmHR('k6sq[-!0Q'pA?SD.Mkpr>Jqg#.
-i>3O4W#gk6K?PQYYpD@t-(=H>J'dmeo.Pf_8OTP7W6Rq;WAjM\]AcamkT58%L18rF[rVHBUjXX!
1HiKi2uN8%KBNnD];F7$JMHaRpf1KOksfI-4:5[!a^3jC]=G](;rID+Kn;[0PtC.\t(F;I>\>TaY
UV@i/L%\Nm<d]=bAQcdK@0bI#!4EV`L)\T$k>r8EJ<bI"9Se+f642_-dF9^\(bVp/^<2,WecbjN2
nHHc)_a7VWDk59!/B,cGkHtMJ!])H]PUNbbD:(ej)DHX*#nf\kbI4>,(b\1;NDCs&.e7C6lE3hV=
^7QJb(Bm*]I.;eU/ml"o,^/KRrp2RD-]7(<OX^g"J)^!D8%\?G#MFl1RG+8ojNTl],Y+D@('J"OT
PIcNOP[aT1:*ZSEc_(jB1'Zc_3m^Ym2g8$[<Tf8];b3M36O`P8(kPI6%;[J(kB*055;M@h<'rs''
f/&[1'EW5ILRS4J"9:)G)mYj)j<*<%NfW'D=1eDgsaPKTWk6pQ_@>MbK&jb_/`0j"^!I;oZA7N[=
S],NftQ$fddgjjp0$c<k,"420mMN'C60[!B__WisrUDp9-UXmkGD:p=P7@u<P=AFOPHAhgt6OPa-
C;0Q8+ek[c?_'HTq4_$=N+&eU_m>8.IQi\Q"F.ICYhfP-sehYXZ=b>b0Ml"F<8q&P<R!97PUAiiP
7V!d<c^`\!jC@/F\kKk4f7P6Y.p5usoioY[Wj1<-EM<-i-8^%gO#^'XUWn4ka%l.MSCTCKc/hBq2
J.oUFRLO2JO,92&E<7qX7L\T"eh-Da/ZpL];o_+YoObd+9!r`D-1`SiX2:fAm@;0]0(#OBZ[(?HH
@/=rN_tQ9%r2WWSI!1d4VUM)j0ZO):reti=<V(X^U&CN()BqR,^R(PVobImAf1VXK-e5AA:c:\!D
q@*7m0f:D1H8X`"W&jSo5Qac^5bUO's6\@6pmk+Qrjj)P)!S9@PSYXZZ\?uPI>1P&;j;(j^nKlI/
kiOd-qk`pO-C;sH^WQ^iISnf'HQ`gdN;Q@*3Ble^0j"_Z14,XhZEJ?d]<jJHg_'7#foXR0KMVl'l
$@J,aB,Vhac8tKsbM19CO,OLqeL>_4S%9mte!1Z)hHH4Rm/!,>8[$kahT4R*ei]['`_2BFcj>;&X
_T@-V)_MVPY"=g()n>uBj@ArIeg>P+t9t,7L\1Ng!0<7`!A8BJj3**40LV3qdW^0q>%iTUSY$J5?
"E5e$/62Y;e1TG)#T?2c$;eeX56Y91doNNo[`WT2&*p(X=`Ed5Yp3[&U-r6DLIF@kj5DQgA;_%3t
.1qdL^R?"G;so7;q>c2d+V.2(O>U8qoeU2Pu'O)<M=Q3NQFg&5BC05H&RbGul"^pOi%`e*Mh(`8-
3nDaphC7JV3i=mlY(sp`FXaFTbRVBScl*\W/WS8AQebeDBHB8/lMoCY7)F'aQ;.\5P4;^PiI:W%+
?d:<#nV=a@DgoDU@s\g=ZB9g%0!?[*8a7pi15o\o5JCeSB,eZI/$46jb91K6Y2]`f]O/lbp5NIDV
n8I61jY:qh;_VYhu<[9Ol"Zm%(%UHZ>/?op7MROekM:`"3tL[=+YB=3;lET4?59!NCrus1h9Ut:,
6emgV8]?7M`;tEl`ea`]&f`'gm4W+l/Am65mHhhCKg71[1t3RQXR;Ggh=Q4$GDt8<c5GCk_6U!1=
UE?bER:NHYk!/OrNq5N1W<97Cnr\d<%Dh:o+nLk9R=`hMnlWbq:pembc<5AooZI`:4H4hD-Xb3SG
Mca*qsj[(75?bBQWb[l0.0ud0Eg$=*,)087Cb,-XoF0#L_DiAanY;gh'EMuNQM+$LV+iAttF]$nQ
+`S_9T1sXg5tI+YS'4Ka)o(DjP#-=C=A[oNbiNW5B-VQn_*m;W@tbaF*d.(<'VCX+b11)H'OB+ek
:',q<_o[fJ`cdPfP=-Z(aYN+F^lS^*f&T$)0%'K%PQm5VNL9B=&I2l&'Tkg[B=N9PS?QHonrCola
5]t;kc]<7Iu@\1'oo`ol'b:l:;uRHoS*JgK3n7RQsPYpr98E,"0%eRVu-%[Y/E5gdh+aI_<jAoNs
T)C"5ki?>g9G\XpOcOG?f%:j[-2ANTdi[Oc:2UomrMDF/C!JsVC7rb$6)'V<MhUFlK2QaB\hoX+,
cA$$8VSoS[;"XL5%N+PeI9,Ae;ANPB7rr<oH[DPmbc/*P7Sh'g@Cl_YiW[*YsRDma7,'/amC(r%6
3A-Pfh8"mI!#Xe(T[iS!!j_lj)fGs^:n4-5;,.91lAG)nLj.M8^g<AWG5hQN0:[6aVDamp`YfS8.
W<mUjUtJd24C/G9,T!6*M5IPmfqgMj"V(CP^RV2e?85s*%?jV-0eTP`E"N]!2Ek>o\",N`f-3?VX
hq+B#7-NU;cfN4@CX]0bKUF(u;KL0KahFWFIhclk3!2'[(+e\E^&u4T&aO%cq0`ja:V4M]bViP\[
aW]k6'-,CVM'O*1+?+?\4&I&5(S\COGces^)!@"+WdE^mFr&,jZ>\HSX7KV6Es5ssLE;Zq8hpOYO
gFU8V8Nj.j[i!F;6R=:[lKr:GJ0>Tc#,`b8h2J^l@^6!oDG%B^cAUT`J!@@\=]@iUd>W![NSO]Y7
Pcc*u<`\rflLYsjW3'7_p=*qCrr>'LBqh\+UV#\)@2]Thn.R#lCKWise*8OYA--$;l><FRd:VW#k
8H+dp;WRak)P't^#l"Yo&KOoSZhX:0g';1Ghpq/@SUZo=,g[bN6r/Ieh!ln/DKp:WnG\f-2P8oO^
beD9t.grpYX'ea@<5JLV$Cqq\_h%R<"R2MTq,%dQQ;hf^tGWZ1--$b<qhg)![_G<8LTSDsrEC(m0
7Ba+Y=p(/*:u#k)]W9IcAW8'p3);pVWb[b^LoM]R8;0Xptm@8-d(f#GAtbWJ_\I\Z(rW+bRJXDtP
\Ce-AkH,?It4QR'DT/HRY%PPFA"uW`9nrOpK[Q+r/]a9Zc5j7`q;B!bNV;hT2Kd4F$2D*k!24/Sl
e!*eTCjdWmA6"+cZ=iL4K.,B;_gl4Mp9u4%1j+#?&eR4%^%8=>$HL]YPmgndG%QoA\gjqLA^Pb/%
bn-,@n^R[TtYO0/HoTf(A+3&Mp2`+:n_Thr\9u4IZ-+A;&jfs=0-m>cm2UoX>-%cIp?I*3ib9Hn?
T<EIQZ]?%GXMs4C`KPM0o,/Rc:U4-F]bp)#c-",IqaZj3B3MdV?+2+>WEeIsbJ)/7eOp,GHm?Xk+
%>@rS\pE-TLs6;iVu?_$NJo_N3EEVo9k*7=1`$!lGsF`b2-N,:@9U:CNTN2tI[^U@)G*)mp)kldU
diube_2k++@nt@B9P%f&q*h7]Y\!BRY>Ma*%nQ@?>(djbK=)k.XACn34ae2hMJ4MopDoO*sbTpKa
LVJt<nV6aaX=24MFal"TG*Vse=qH/X,BXTGF,`O+bf`b^'PdB0T.pRW;d"3X4.;F;M^p#P`X\fcn
]Qg7I>262h:*/-(UTHCN(>f(]O`.h<,XXm1'pF*=VnkA'*Mp\27e3DS@i0@<$Ds<5&uCRCQ2lC&a
g0B]Ol./J+PLg<G6&$S]-$KRU&3JH,0)&bPj8kd:lh^j&=dJ1RGebaPW%\lREK?Puh.:DZeLFDO2
QgmCRRs,q],sS@XH>gCgK)g@CU!&VZ&)F*-Fd0e<(Dmk&/W$@@Pe4m4r>B;Bh/j7WZ4-"Vm=PXM9
QZ/ktN([gl]\7<9uam?[g-S<JmO#Mil/QHiW%r+SiQm1'XXAr-r/,L71?D<btqmD]Zg97&c-RjZt
^R\oa#mJ&qo3QfU,;qb\dsK3oe?r)2d:"8\41;IV\Eo1(hOMN[b4L]h*O1=SJWJ6ZHYL1lZUf(6o
n0j3[=V)!FOdpV"b4fS)m:poptDL,(*lt+epFL3WGEUHR%_ug7JrN4][J(pMf`k+Z3]>7Dq20$"[
DU;2)3?TBY\\N)&m:/`])"k>BK0L.\<qpQ%0;Em]_sXR.W!COW>1)[hflQEP)QYD1=R&4gnb]mu<
uIF5!*Gim%o]2L2qkZ&Nhf%rks"Z?^Ks=PIh<?>tS-D=ltS?7rnKQVd(SeRgF&jZ^9o+`6B+8$G8
KX2ssHI88ph3R.Bn^J>(=8?53*DGcB?rqSaU<G,47kBN-S$s8NCjrHneU@D8nMZ.arqg;J6c?Ti[
h&b_4^K>+lH_UI.S38diSa!d:Ofjj;T/n^e:c)JEXGj6R6goam?`UX"bG_5r(@PL9]hegi^e*pn]
5JH`NU(`4!58ItE;*1Dh0Hs,mO6ug`uBnNbe9oHN*,VFk98E#oRju79OTLef8NJ.<mD_(-&X1B<K
e;70r$0cmAi5`,">j(X&mRfPbfd!PBWqb&2r4U^ZRCEj6,8nA4FT:@s=(n6?@\:pD"]gH)Dp+P4s
l/C=esgT@;P>%pr=t35brADg<DfNCT$hb#2OL!/nlD<_Z?:$;8"KCNSqBE`NeKG>Ff8$JT'f]H2#
a>soQG>l)NbUT#E^4F<IhV>gP@',oUT.q[+Q8AV9&q^ki*3h;G2+t"uoaiZ)&YML[l3`7OK%SqA+
`^l/-m]Qjil@OFFj3@sQ]B"pjMoO^fSg3%3eU6L*0)kYIXZd^q>T1:`A_X+QK-dsh7Jk112i@8[@
/"B\GkDob)%RS3SICc)DR)+(gC>-I6T9#>DFSu/hi5,jW)`<CCSk/oP[3@^Uc!PgmtiPISAT?M=9
2CPr\7QhiHWM`]Aj^;Zb+NQJOn;?]D%SLYKLRu:OU9]-p=j']?>hNCIS-%>H%e7iNbNL]\o$7l>L
%\9Jeh_9TAJEgYKL%prV7V5&c\fDUTG:UE8FiF*-ObA2sD(Pb-`UJ;cnM^O4Of>nN2?6[>s_UcOH
4\]pS73^,SSVR!luDg.>%DR+l[,[.Ih_2_f-88j6fi@EI_:OC-8pZj$Hc.3O%RUo]QDK2edUF-iM
]O>Tr]U12_Pue>7k?@]"n$'Rm*p^fG20T!"39Lb$jQp1Nn]0q0N2U6qANIt":,i_'QC2N'm]iL@c
/(lgO]?O/R,LPL/T[HVCUUM,ZK:7P;"35o]j[qtS.Ur&<U2'VcT*jPa_nc_NZHO%Y8JQbanElH51
p$]TtPs$)%UWR((YQF35fXeqoT<OnRiU:!'"CoCsB,TViELE(muLNF^B-5ZDOS&?^s9Dh`8/4EF*
Ve6?d`.X&rdoq:'MXE?PnE^A!qP0!8R$9!!hm-Fb#1'meJjH'H;/3ADs,h94j0B%;uN=6ere.)r]
"`^+:lSgu6-`=[@<$>@dN;XV#OgjKX^iEE7bN(M+<h5C.mLPK.aWfca2U%jJRgfU0>L++CSDHj"m
eT7j%($P:WZ/+8?+a:r4LS`KEg8UF2<;KIir9!doERk1p(^k9YF.`%A75onc;9/\:(tb5>B;;iCi
s==\[il(V41`.A)!Y.WGEULoaaO,g"%eD6DRK"o.Zq"2a>/-OUQ`tP.]7a)RU'%Da'X4#"%ZPE(P
/kHikG2`R8E_(3eTpZO\`Po<[6,a)Ri_')3QK:pPW6^C$-J29p4%C_G,E!NujneKn1=qCjP8[WAi
_pO67:J$D>Q968Hr?6tCNsJmmAN5J@`])SoK)Nu4oro'7nr(N"k6M]uM,+6*S8hV9XSJ+m)49dK#
EO,Q4[7BIpTPNpW">Kn+P'op*OlI&`UCpb_ZIH"1Ljje[tM"#BE(bJ06KI!tHA&G0Ne`4YRVpYXZ
A]-uiUiZKG89UD'6nm5OR*nMOaL;kdH[a_m3fRV"psqDkQS2Id^2V-N.g2U3ZTr7[#E1?.Mf`Qt:
K0IR`K/L?Oun-JAo1ZQZ8.<P[:Jn]4?Lab2;Q<SPbSk%KBgXM'a`XDH`6)!3B@p&ps,J(o+4F3/?
"ZM1tK]&2bm&A6a'g$GcqNZpr"$CE$%:6R7*F.8(!>2bbA6>]PB4+H2]NEb>+SFAXDqJLj,l@/,7
,Cj]:o0BgeKK;1DGL^#]G];fsR0i3KM,RQ>AFnW\urL\Au2jW+c/20_4OmtqQVeUMk,rBpHi0YPT
8bu,SeYXVjs!'&BfaJ6mn=;S>l*\a;ci^t!7rNZVgK%2f6/DGG6fZ)P.`j2XHlSE:n7q5X;[JH\.
ZDa8Hf0!U?Osq_r`F^cm=FR%$+<n;UE&/$SEUPp&2L:c[Hd2ALHBh\01A3,=R7/RVXQQS+8<%D17
P[0&N*/Q-%b8t>RAuh(]^Ae$],d6[Or6Ka7>;ipB!)=B0d9(8n*Pbl%?;',Z).G>):92,^DT:GT_
>/?g^khFZL+`g+n<S_!aZ)5[j6^6;f_&[+`abi]#@]B;s_08A8/f&I-&:ZW1Qk<=;\0_n%jT]^6V
N+?)?F"4l?F@7;^qs=O0lDh5Ik%8j0IHf:4,9U5U7ne-NNnXD(:,;7-L4ng[6p]G[#0>cBa'$Cst
j]0VUP>+YfSb#3rfCiXsUEUDpXiE(>oC\Gf#Nr*NRk8?/l7OiBG00Z4/5ml(s3XHhsmGU%@?WCfg
<[9!H`4sd7D]E87;KJ9pfb(JA]boi-A#f_pEJkL=;J!K#2ca.f!Ac3FWAVkV8*IRlD_T0WTZKF2'
7Q#pa]3n`qh*0"51TFKI6XAB2<>gu<$ct+d6-<F%Z[[Ba.6UsjP3M#Qek:!Z*LH<*12oSJ)I'EPK
j9A>*!H'0*FA1cdf[/Eh2a830Kb(3a:OTJYssk8RcUNZR%6"R>QIgrPIslIi172p<T[#jcM-/FVq
b?n()0uG2<`"^7rh<@65*l$QVYVBZ>YT5atO:nY=SE0)kpQ\r/5l]oO#Da)^mKIPQD?pispTqeHa
b:0+i6p<Ck)Z;F14P5-dP7f-cWHHI^EF4l&W4(NK.\B&m:.;lh1AN"PkpW=+@'\W6&9;sF'93aV_
h0Ja?-$tJ-Ae@=lNs5ok/aeR"ac;X1c;HqXD/@4,KI-GaMdE]Xm3;_Z/l".gCt5'6-4@<C2'Xtbm
1rdX2@-a*rNg+32=^iE<Y(f47P2"-`N<c4eQP"`iu/YJHZ9S14:U3>Q?J%Zo8H<s.d*<ICM*0/+S
1U@Ns7#`nAR=t54BIDV1%CYX#J$AHEDV#lC/b%l[@BE(=S5,kIgu1GRZdl>XOQg\]9[V`lsT;!me
SY,;,S&S:SS$43&,HhF*[Ho&NAI][d8rqoTH;A/9rM(t"IIai6bOqNWm"A&[?_`_JD-^E79f/PYd
$!5m"#c8lItQ"I4rCAm/dP*Xn:UBIPj-Qb?DS1Zn83QRanKCo0DnV?^`3BL.Nm_*Ss3Nb<qr*Ci&
46<Z[%f!3\Ht:c*&M(AOpfscFc)7Y",#5Lmi%=MdlQE:k'SUnUZZtoq33nn"`_&C#kt,[AZ=:e;O
hQ+RdctV,QZ^H&j-l:JR]gUD7(&pE<4@<H!@YA0na?UT<@,-u'W_5>8@MB%#$7OpWAT(#%\;"EWs
^/!hG"c>=0'JJY:Ai8A4]a,39)DcHehipQ4#iD,B^3M!BG@/C5bfL$C(f4"AXG_W3g0&N>VUC&t_
7fk;RKMPFaL`hVqZskj4QRPC`&>0r&YIbk2h],cjXJ`R$W7<uj*^ZA[\W1#Z1Y3?=[n%C&S!N^lb
i7a>60R?a*`=jIBo%COB4HXeYc_o<qNXcM3d680@Dk<%19<'NAB!1;:b58eYY(4*ajMdEY9f"Pt(
Gol/)rn';tmAp2l/Lu9NSBR_l.p/LVX:,[[o=p798.h0D\'8Nr<r?HN6tIUo1.G8DJL4"T$Pf\^N
D!cnb5?#QjV!>@Ue+S,N+0rIB<5KVo];+(S^BQ.'C7PU\O37F_^ce!#P3rVQNUaGC>Nur??KB#Ah
9)I.6>:t>2lk1+?6f)2]MP0@dO'a.soh[So?FWJ13jH3@9Inf.pMh#-h`!PDsL1`W;?3gZKs?&RV
&7]TuM4%W9q?CJW-@mA!<c93AbhR^TOjmYT8eg_-5lNcT$`P#t1:^+UcimfG_f;dZ?WYoW\UnWKl
=gbmelZUX%D==)]bM-StoE4-dpGE<c/T\B22N&.6;=q3%f8joPCj&k6Zpnotn([7*>8`q%uKDol%
;M&_lk?(L0:)i99'e'9^@qIi4%Io[I,G_CgDo^bd<(p:\Ee21kc.Xls)&J@f7^GrC/JL<UiSPNKg
#I=?\EgaBQp"%7T^)I5HA1t'T2LMp^DRODl<O>hSE6m5UV$hR30NS#rrD0\!4G>]7'kH3o:kn1;S
1l2lJgXbalh`*Dd0eY^!ECnU:p(_Hp2T_)Mp9)=Y!?O!\C.(lcsT=;+ch[a[f0gN.0DfqM+]b=]&
%TDDG@5>/68DZ<,)*;.MN>G]4ps?(X;9_51Hl9D&%T_=%_r6d?sQ`<#/I&3c$KXK?)kh]L;<RnPc
d7TW=gqb-=MrrAS<9.NZ\8h)`*]jYPNZVSkP9K.&r2Q$-/PcO*2qrHcjpF>XW/@psU;/0Im)[oS-
Rs\[<l)%IXmkqeY\m-Aer+,p]`L"3\do9=JBp+F!Fb2\._2M#VrC=4Tq86\lm;(=TA;^e\mumgJ7
m,nOdTF\J>b#8b]DLVaInOWs=:XQ2ho8D*TJ%D$DPp>!.dN%`a^(?1;]C/@f-@Z&aV@q$?QU,#lL
$Q@iiHtKH,kTh"&rjidoY&ggP<-F7+UVA,uL]*Anasu?.Z3'af8Kf]gdX6@eWE3:Ad4d\bGi5T3U
<0E@RY\0[k4>#rbWhnWQS'+hnLc8Jr,?p>fj8m$(ur`FB!57S+AfT:B"5+^4Z[\PS`5(&Z$g:@hb
bq@f":X-tc=Z\uENj9A)WDhZ&#Q*6]eBnNBp+"$N(Yu7[\]I]DmpWXqs-c+gLTL[G]@+072]MYId
Q,[<2,WEEn!&"4Non8:)RX)_k?RMqoU[*p!L@61'=0tdCrj#F7HH1Y.CnI%D(mc%FYXo-%p;Q;n=
I\,UP)XAM.@GA/=_Kh<&e:;XIscp2U&a'gUH&FY4Co6p[GFJDnt,Aa1o;keN/mnfnUspkd(JD6h+
dn%./osrBi!>5:stnJkKXJ'JO]IJNUe_dEf@g$?Q9r_&FjK%kK[njP^.[sRlX<,nr^s5?Ejtn$uC
-X@s([S?Put>]1db3\]YFM9kB/V00cmC.%`W_:.t2DNfT.I%s3ePS&t+$\)T/I]]DL9OB9HlUHK(
$^/(:+N)_d9?p0ECWu;F,VY<U+5L\AqNu6@cVO*q@.QZ`q-[is_ps<"oBOk*crEB_@P1<P:8HdLE
g@kKM6_uFWnAWtk[oeM(b.8B'P.^>F-Q9&.6D$6noNrDtMe1Fd@8lDT`\^NWf*,2g9fF@^ffLYND
gJKJ3g<S8mmE<VY1rQ:6Dp:C@Hg@_Ds`3o;kKUk(m-,nlgKI_(LOq9@XB.Y(csgOW^Rsa$8Vd2pr
8!Rk91Ha"MXZNf)=MhaSW'#<)<q&]KJa<-+=\MKtK#A;BR9p_hm2aQVJ\O^:XVbGf?6*Dk:sCS_F
%TArfnq9V$X%.blrn5'$(L7UQh1T/])/Fn:\.PK"A+K?gc_Q`3h2nU`o#]CWYU&pm(qh\We=[p3<
b+,UnK`B(G7N9,pRYKdf[1=NEE#>l^:ANDX735:MWbB,SL\G\Fqa2V@S7s4IscT&4u0%`cTk9b<c
=(Y7?Z;G/0B/sa5J"-s`O5hZJ\hW#VFe"Sa,p#!KEfXG3''[k7<)R!s_nu4h^GkDdS%![h5,IoPd
Je+-68V^.,.\i")_^&]9^3]Y=jJ&UDY@XO/QZHZAV6X8[FR9[]9]WI1^:qZ)C=W"_AY4];7pU9^J
/d657\O*W4W=1WufE(NSgnm.)@gNFff+s&qJdZE6;$,Y91J$,f3%2f^l6cZlS"aSE;0Hd9gd3OK$
c;C24_K!NQkd0*h0,*h(*#goM[P*S$ae#0lIYFj=WYWY^/nC-MDCARH9CO>=a17^7>]iFB(e^_IQ
q>Z%oPVi4YRLt;*C1m_@?$Jq;o'2.c&MT$UM*cg)*gPGL^aS=!3YofiK8eP=[AJC@2[Jk0>(e!!5
"JceVRZi?Qj\`OFb9bZ+(,f@(NA;uSe\1Db\g2joQ#=9IG7Ml:]f,J`;?Z9d%sY9cX;[;Cc;G<R,
?L;i<t,uiF[^A?:M]W>nX,lBq2`U+P5bNmfH3'@1mA``;$+MFVbYm/?"1D:^)YO_,>p<qVA`?M]7
)Gua&dMkgfnO_`__I9UY]QOm6W8*GutH_+af\1khV`BEi2ms[k>Q%A%.;B@=g6*M)bGhPZk_heEu
W0CBDpN+oC9j*cEG%JpfC.`NUu#_O>R2Vt\JcW&45(Cu/3:Uii3`8;sP*h_S`[d\CA3n73m\^?Ca
;XLe,U<X(_-4R\f[25nE*G@Pi(ChS#3rrA!'HQI!O\@_1=frfATC9i7qK@H`Z(@cVd`O=JUOmIq$
DcjnXl,pu?AK/nDPX.N/\B%6!bnsT:6O&Lt#gn,=Z^YEZXB"li@a-ftcEA@(ocMA>o8:en3:i:$P
;MX0S8kR8:QJYXb#@*]:[RTQ-agc>OcbLr[rDdoc=ZdUN+MYTS#o_&@r6U*LJh,pc3_]47^<AZ/s
KWY7T%o9SEGr2<SP@M(mO)t/3fZh)6-;Q/!KjtZ'`WWF,;?9;i4q%,RVf9>piI1hOZq+^""DbSEa
k+S1dc<eJF)o:$0&o0p3f/I!`L/<"49F9=t/W[nUHhW!-5^ql_t(g\1G_ri'6+ETG')GBd91nV"#
&>)ZqWh@?#)m8MP]_'qu#O*43]i3KQD,-PD)>oTjc8tLd$S;TS/F']=rBZa?1aCP)?O^fk*8o"ZR
U9S:;eh?/WRQ`lHr5fT[0C>0C7Op+q@$O>VkW$h'N^M&.LXbpaAYbi_X:KAVkK[et2Xd.=p&>$3>
t,+>Mu.9F4f3uHPrk/#0TqD1Wkp\kUgNErOXo99k*j9#$82;<<P)%,nVG;9OSI'<]2lCDUtiMm_u
R<>:;,d+[Z]J)AR]p=(6(oaY`67IM-q.1WEXsmC*pS@X9@d]^U?u*>6V(HPZlIH@HI*9PBuhO.$B
K$5l/WIBZ[qGg,\"KCn&.+$<]5tC@_R<4-KR3c&Sl%5J5CS`g?G:.2k)[Hd1&AL#Z8ZG5hP"GgdA
TC&YDp/ZO9q0pB#K)*fY%n<&_/@cieJS6>->k;GOs6l)C9F(XO5_rQQEGiN@GNcL7(pLbi<J`E9(
:M0CDN>;FIi7/l3j[7FtjWdc1R"mc?8-`4]#KD]VBWL.sh3GWm5bL6t@(1oMc'%gUW_;Hj^>*D>;
Hhmf,*G(=D^Q7Y,.R;U\u1LcD%2r/)BNJ/in/7S^1L@JXCuM9G@(8(KXJIJJ(%fc7R^[:hF1?/f=
V=EBuS[-^[[71GO:g;eRIo'?L_k%XA[QN4lZWJX7MVaWDrJV:?%0d([/K?i56GGX`8[,U=#3Dmu1
=2]OVc9V\u%VZjl+3kK^d#FVL$/,\L`R`9pIrGqi>OH.?gJ]$jcpm%\c")(5D,=K)`E%iDFRg%NI
EgaJUS@P\bMbQAo_CL99p_eZnU;(?6p=+k-nWAW]C/B[SX_/7`\\J0!lb#KI#H#F"EN8uhZ^2\Le
bIDM1-o].Z)u>2I0W0t0j)6HfBZm")50OT*[eGM28L_L$&P]Mu$T7AUC,.(g[Lr@QD\bjmHS@#h&
hf2]QW*j^OX.CS2=!$)dC]jMm>:,u-.J*%-JVtP>Dr?_$-ToZ2u>TMDY8a\_-=tXaZp>Y]]Xt9)o
moOQ'FQO;-<L/>:D\r%i9B2A0&A$cEbJch(<"l8Lc3<o)1eKXkjPK7WPlmZB=CN<DUVD[1icgd%>
]`84K09GEh?6=Fhh\Ldhs`gRF.gVfk9EA3.NH[4T&(So$?T^PJ+a3^g;#\mdFUe%tb"OXf[0(?K*
V1[W]shfob#A5[^(!/tIb:/0naiVH&X])D5ZQK#LDe`2*kC-[-Hc5#T`>!l_W/aUamTAg]5FCetW
WE]#bQI7hY^D2L/&c=pEXaT]bXW"7t3>C:@etQ^3UDTQ,qpm:KK1bp\pI>,sc$7fTVnS8)q8Y"Ko
9bZ``W#ouGkDJM+hMlaJ8Y\T'Y_9A-hrSN(VH*W/il)S^Lnl_BeX%kp,OeM,c>B+[gXVjC=qg\][
+$oNA=ee\'GnjEt3Cr3q=C^%cqH\Rke`]rT*Z,B=*pHl^[.*G?*r+@I?6t<b+T4A\T+O77]`XC<:
LA`rLkBf5<ud0;t(HYan<O9LdKi._W&e<F>9\U2i+bGffU>)Pu%DA%dkG86b9&+^-&cJ?M1b,\,V
ELX'e/];_q4!'$N@,moF2=fGQAojpQ]^=AD2[[nNhc.&(jXDHpAH>dlG\GErCqPC/i`Ho[U@2/sP
c6mF!XjV+q+>lenQOqU1[;O;MPfc*cLsA['DuAY;OW_7k-o6AUX:^d%Dbllu[-Cftq8l:1`Q.1G]
?NpKg:pNq@sZ'qLE!sH8#pVeqdTh]d.7(Q86NqfCcOA/>,pW*>,8XI!_?<-lS?n/,1>>;bSG2"N5
kePnEYOa(_]OC?""#*^T5q'UFo%+AMWWnF+E)aoc$[21J+Y(#bVl4Y)HA<T)SM>[#QU*M5:PqhD$
YCRbsX*\reU#3TDco9kBukXB2A8qK]^\G3W28Tf[JiPn[)>=6g!sD@e([6]9HVg-ji3o9kJRi:8g
abBE5:]L@lm$2i^-G2PO^4jcK1GfZD)nbZrQ^$0j_\e<f)oT\^C8DY/jiS.SYG;idaS7DMtLW4,_
0ifPhp3=Q97RB\E>[We\V6KsgGL]VniuZ3$\n\!I9R^r'AM2]eY.cSZIGI(:q(ZMH$J39l!csH+6
]+W:/erp%%s2u1c6QSMD&Cib2sHTso3"(K8PuT2EESj<=*`-fYBo#fbj*9NB(i[^N(mb&Z^Of1ks
m7DC*K;AJF\r!V`%b;hqY9oacbC-/TFnCI7p\Q<Go:ZqTM^KX2qT3p=jkp27IftIV#WY?8"1NrG`
=99=`FEbnfD6Iibm1U<&-[YP_hd&,=BnYE.>[$R%M#AO9-n/+*4BQW*/A.fM-r:jIM1<jDki_O<"
29C\!@`O>9Y/>YcqVp/Q:`ZSH3Cea%$=oDjGU5anA<!52'KUKM;GA^g\"u`u.]NHi:SCUq%)FkaW
P?**j!p[5j_&8P=(`TI08F?BMJa2s@kPSf^,VsdRE??\]pi9qHPb%[?@mY<OCF9J'hZo*'[8`Q'=
%k!Mes"]'RB^n-h.0q@9LW2u`=?h3C;bgPQRZ0G*uIbLcoW56ALt>]N*@.4Z7%i)]TT>%VY:I4nn
-!c>C["Eml#&u5%JhP.r(=;+o^LN(:02sS8GCi#p$Hmk9Fp"*:RN.O-\R::DNWG`=-FoH-1_5C#6
V+Uf[B6E,,cTG`Z'iKr^!@Ofs_1EFJkI(rO'9]$&HIe_?N*9DSD&mc5?HD\\@71'bGO"Jr8DLcj8
=`iaPf-rOn#o2tLEXuFWV*D3ilI,Tlf]OD,9_EM`lB5?VhGcI+Bf<0+?\CW_U@SmF?Z+=[jhp+\1
:JN`_k?(5BQa2SlCBR[g]#bTq=&)XL4l=!k$t#?f0dN^=jk^hs4)]o8`KT0EGuWjq]\fnP06L#4e
Pup:)-<DK1pqW;nUn=bI.N_^/p[&?7:&A!8lD[2)6Jj3$G6p$U:BrdQHj7QBdC21h_\tDA4JrS-Q
?-(RADEpcb*quUg8\/$7$opj5VI<$+p1X?]Eu#A;GImitCdkds/GCn7rJmaPsEF2'/=a?Td<I!,E
N3IeX[>0&-Q=rr?pmO]tAa]k31Fn9`"-F;ht7[!h5F#)K/1EV)RSgn%bB+ELc+Os=O>pQsf]@mm(
mij8/[3=j>4I/5Ot@IJDk?MCi0SSKek@%5cfF%'V9,'0jMl2Hk4Y&^8`J.]#\I-5G,\mgN6^ulM%
,26>tcV'rKX?1Pomj=Fhr'*R>i?/MR16Ca\ig]+60!AFB=rVLQ$(\]7gQA<W#r:[6N1glkk3e)bq
plhJ2rcOKe*3#EXQ8^>UF=WUj"E$t:HiPSIhC1fIa8lb,G1m\$Ys#%Jk!m_]jWt*AM)g6M;Js6au
,(]\CGioPEbLYdGW'0`N&c5S:b+4<-`MY]fD3)g9'\2Gg.OSk;f9L**8XbR:dYM=OG_,gnPZRI--
ZZ=&^9u5jCe'!^<,=Qpe!sBQci#)WJ'93S<W'Ret"ChZ_[M[O1O_*CS;?T=*_!IIS7\j#K\r]=b4
BVD9R_-&WM1<Fadko2q/.442>"M^:['E)3MP$qE*QHhea9l,;jYN-i1ElP$udP#.O)eIgDrZt=<V
[rg^&?Pkl:Au<=ZQf`KkWfL7GVIM&3q"9WHIB.-:S>;oZXbX!E9NapS4&>7g.JEJ`LW4(Jc!dHPn
V03MW9]omS^O3ET6Ws^mCVuM-N7!7PqlnTrqteGaGdVk<^`A)<@UnpXlC&*=Xrb+Ra8(9e!JY8Y-
h@mZB6=i?I)b,Su8?2o9C0C4PjFr,Maf,-.TKE:m*OXrFYJLcoWq(R)eJOQ`REgAS"GME@9=npg,
^SQ%a/4jSL.E9#^=cQW/GH*i=b)IJCWBUHSG2a=otI`FU^/8Zrkl_G_YRFIt\SR,,=p9Ug6KUR]f
IlaDk"6hKGnH1?kXYnHMqkEOVRQ2'm#Il6V:jMOac<o3ap;.<(7+o-%q/Y980LM3%"3Y*ZC"^ms(
l0"UJjO-:Cn5$I]j06pq=Hu>5T#Y2>h]jWiY8Y3M,1Y;X;cSl[<,2WJOV"YGF0^q=l#C2VQj\0dX
.ehIN>R^S13d&m^!3/aJ5P@e.Ofb$!+HUdrr=ZcHaDCYmB)o;%1,4Hi3+GQ_/\DaAnHe4CjZW#[_
'QHBr!E5.F]g@J<F`Hr\>nob2h[UN?l7(TZ)52I9o)n'2i(4PSjaOGE3g4`f6V1,9_8SRk]qu&h&
Aa2jR%4WRT_q(LAmZeSR<eCm?#e_`:8fkM@[.Rh3=P^GAQu8X#Cq'(S`DHlkA6;K=5_Xi9c(eRYL
<?WDT'<CA?Tnqp%CATg*i;)['Bj"&UY!;Lp'noH=:cETd6)[_iib9jZ=c=mE2G=S_i8TSTLaJ%'l
X3*M/H$'u60"%dE._7%k(m_/WjM9?EEn3FDqR?4@ROHarG\j$g.NQ-oi6<#q;<*U)_Tjd?FQSbKn
cLCK4+DVlpm/Q/NAU&Dij4(e(c9n>(4!eJMKqm2(>`:0$RL#jaokS,V]*f0rrDk52=F!cm4LH>)U
j%IX9V\7d93;D#-ZL4.YFjj'smACEd^]X(4b!rS^D#Nhsl;*W]qlDYWLDP:?s6`0mNTThp`akNBK
3p)T<08Y1jbA1D"6*Ys9VK<]GH;k5^"BhsZL)qLZEAmr]Dd_j'5nhp,%Z8_mVh?DH72N\XE\7s9g
1l;K8O[nJ#.CQE_9G2'[WgnZtcU!D6gfcW"?abWr/_l10_rAh355$Ss#N@sar='-tBnXsIrS80Ri
pVcQJa%]kBEVa[KY7<+4f[/(6b(m5c*8R_&@.Rtb:Bqb604Tsn?\`p+3NlQ'Ct;#\YB+i>&R3$[\
,4`5H)g$E(m;&n2^CLo&/L36H/"QEr?]kW)CC43c!h;#qdADF23eEF]\-Vr@84I\.l?C5"#V!,qr
C]MG)]P,bOrSAnV@o[]Q7M]>r]b8,ar#i4:hTi-sa2[o&0S<=B=[/2cP42JpsE%nR38&?R@m>)e4
+!/]ts>ACnkl=TuB4h[&okdsi2DF>W9n_>/scM$P51[u]9@Z&7C"$ZpntKTiU>Y7n9EQ2&Q\3op%
!'N%j`R`FglPgT#j?<g^CMGA1Wp0Wesn0Z(smq,;'c&NX>2S*ZTK#_^@mba3bHWIf=I5,%,f]9Y6
R\M>N.=Neu(Lf&fOrN)k?d:d2<ke>5(,a'B@9q=\9??bZDn=/kg3S/)VS;8JdH;+&ds$&`4&hAM3
%@"d^ZuSi[:<1r7(GDVAmgh.gF:0o*dm$LKS[q2-!"!8U'iS9ejj*rI9o%rpS1V@Gp9&Q)Ej+'rY
P*Q?$F)^NH(I/5,[#Em2A,JDp52O]_JRo/:1Y4UXtL*T/geZq]a%]4Qi?alGb%rC3+ZFonpI4nFM
en\a#MfpUnr?&rt0.S;1&bl`p&7r'WKMm>3YuX43+<rDe\7`_dNN`_DM9Ddp3L`K<&/')ll38nJL
)<nl1NA#ASCA%EgnB:N;!l#u.r2EeXPc?2S-0>,&+];ViBAKK]KGL7A,9%a.Dc7p0g;I@^_o\94/
>M_(I/I1[/UbaTg'^G0oX9ueoN2$H,?F!#neLStI$;,XP3T&R4olT%]#=2G[..-A=_Z1+"oc0qT,
JUZ^6I=7kXkLMo!GtR>k7re2pr`)>qrH<TSfVoSZB%]1kSu6l,b4\CN2Gm+N$OgRQQGJ<B#U:3.G
_iXA_S.A;i:q8pM;h+?.$.40r?*HTi,CFQ00*mA8q2$="`(B_QpoA3VA^VdR@[%gdV#<!1XIm&A%
U<o"("uGb&O&g.8tj9WFV/Sat2-mGEkE"Jmck,Tr%aB_ffq(7=I5:78f'U,@8;mKA\Z:UubI5BY#
!H"8KE8buZAiC+=:rifJPq"9Si8?OUj`,ko`I9,[bD^4(0kMB6L*ED)12"U8'/.]0H>-2103.>hP
W5eH^1</06AOMr73bpN%4&7l7$Al%*&E<"2WQG[?f:4*f>:XsIHXk?8`m4#5^+1a-!f$Vue!_jGG
H=0V4+-1$"VU[Hc4n^0mOS4pJa$n*nRXAJms,EC!30C,S3^ja,HoX]P"V%sN0E&?h=6u54iNCiOR
n]&d+5bLn;18#@T"5cmkhGgV9HbK\GlR#RLO+`I%:q3_<]r!d_8I3L=]8:IPKk$#+fj3Nqiq238G
(F!TeYXbabii!uRIuH,Z;c\Djg8QBh6uQkDltr\0H0Iasl^cXn^Q(V]V5HciMF4T+0,=4Jfk-Rht
oYBUi0#1iI&W^b>nH#cDOL^Mo""Jh\;"\'l0duq(/]#ahW[PK\-0,.7J'i-sX"'X]h"1J40g[?kF
WEq>Ucms-]]O'?[Lm<V')`O1OdM=?eMD$7[YnaKMgAnWeW4m8$1S*G`!Xjmjm]l`[:.S]X-?U'Z"
a\:FKVWfLk?.T<eSXAXrC&quB[4!\h;];P+#!#4G]-*RP]+&r<+#)CWI9Pl@PKIcrT6,R$]?;jjN
M`r<1IkZ+g+DNDk(/%Hl#msQa[T/gfbOnrb+1LN%sUU](M4u1#fn/N/K4t[R![L?2TPIT,2h)(/_
2JnWTK;"6=F!=V3MgoOO6JhYcDXI4!+Q.TaV$X4j?G'q.\YQB#rq[6qURG?(OAJ;*qK4&'%omtW5
&pSV#\jgm5I!!NV<]++'?l7\3WCCTSFUK/c[_ce5QB)L")*FPml2j0<;H/g<`05^[XV_0AkLm0>u
NRV=2Zp+:LJ_qEt]OpW3fdlU4:1^Og1!8koe#]cPV::0:\<ZtPYrY\rQKe.52``$C?_(aQX?c`iO
D@jtk,,Sr]VqC5an)ZSNk9>k-n[42PjsQGj#+;I``tn'X*rm>)sqF5gCjLXP8M"ej!lHWQ@.`eF/
AJl=t)JUaU15a'G8+OR`^%"EcI9Z`prj^n+jb@mPK?q!@IfJ/d^uJ:WJDe#nHH=fi+j6a!LA9c;'
$M]>r_L)0!t('JJq8)AKQJ`4Hr<@>qf_h18#k\1t_q)uF(W)-ji7OsC8?k9U<]>W%58)tSdf`P8J
:1u^sQ^RU/tFk_)t^T2L%`C-k-I,]j^AC01Kr8uf9:*`iOH.djcg9KJn;-9oBn`LZVlni,</[k$+
0tL)-o&)5f+$Qr&\`aH9,&<AXL0N9uM@G5SLYa/92BNfsgY8V:C#+.5\J,?f]5(GU-,)gm.,.hK,
RRa^H-6dsmKNP^`\lT)iX0JH:u]0+Z@NXjkJfM$/uct7,:a*`hE6.Jn!Au@@Nj[E38n4^E?K,(lf
U+f,>2@K#4@<I!;>ph't6#2FM%($+H&ps'T,1<5+8efg[8j&./ougnX#6D[,OLYI]D]_g,IfW@M+
T0d*a;CJR]8Id,hi9mVR+8&<P6B*\sW/f!gQ2:kc=5oV[YL.PWDS]#]!NWVZU7OdNhcY&W+)3"TB
!`]$:/$Yt4]?rV7!;A=[gHt<4*R@U1Vj[gJb:%<SOX?u:^nn'.sjF+>-iAUFEbDRhd[t\-)n<[ZF
[E6UiDfsb!(o<Qn4/IF:3TQ./Dk-*T\=31J5F3GgoBRYNi7r,s9G^S,>NZD=e)C(@aqZS+B_-pQD
fp6pJ[>mDX>2B9BCG"!5>70j]Oi=ZVF$SgNO[Y;C&HnH:+\F%@\lN(V&9/c*nZ'o/f@:R)Y2rqU5
+pe2n/cBGflp\8mdOLjp+#(o95Mh4"S0@?`7/h]Z?fZDuTg7#^=-\W4X.h[B`"VqAS,ZA[>+M:AI
=r^ZRF_Ptmqu7PUEsD0s=6oYU?J,[4Ocf*\GB@?c;SkPE/Y;C//4>/[:BmAb/,*LCj%rKuK/*h9j
2Q7PS[of*tj1jHZ`5dp!-0EH4OeanjV;P_iI=J&H7]L\J[@j>UVQFo0]1CkPrpAGo[Au%oOU-79(
Pj'H2lMgk08fF;FESAJ9(i)_VH]rg5,R_=Yo6fXE*FEmZRP7<IZ,`m/6[SFn.WkR&9_F7mrrD]4q
c=VrS,VTkN3SWb!ke!s`9Kk%S_e)(ao1.1q1mIWAo)3[fTdB42htpXh.@Qa9CEDR!$uWOa/$DH,e
`I:OriAX`g^GVX5_?^qFud\nDqP+f4t%j;h"D!1>s"]mh5SEoBR6FY0h^?KTrIEo/\@!H`InKI.O
fkR#R'+9$,GYlF33\@Wo]I_26uX9cQ,kP<*P'NXIGp?Li19/`7jO*o]WmiDDN)NSIfF@Lk3ENc@X
;GEnG?J\/MQnQV]4o2q0"qMjHKgKinDP4>UEKBZ0Y]MYYLXt'S.7]Td(UTdpe$(-@A7eMPHm0*a7
@YnC[i&>Rme@:Ho9#9(rXP3'i41q)_rrCioMHb21?8GlW0A@^#,3r$pcVZ=*aK)W%<Tc.e85`[P+
KJJtHd*N#q=0d`6K;j_,AsPaP3p"Jn3sdLS`C5/dN-F#J<t@SCX-YA(.Rgc`/WBg:s(U@[s9tm1s
HdsG+Tg09b`lAj"dEmb7X_mS8m]p.DBej];e&8?nep:X5bUC-N%kIMC3'PjhdgQ3kng."He(<mhX
dW4Q'SHq)i8e](MC'>#MnPc2pVgD"il:h=6);NM.3p5J@3Hh+k9V9]+$^1A[qpH*#.$(@kn6P3kb
shJT+E[tWC&A)=rmR(/i6?ArD!13[i.V*EaUIE-+tHRg.%<)5Il>C@hT@;auW1amgJc^ZeZfrm"m
Wm\E)Xl):;U0*BOM[jae1+N$9Yn29<$"*:V\K12*(V*fk(^\\hX0$!,W!19!.k8qHo925"4PuJaj
pZ'US4jG//sAcoFWSq#qoQ,WNc@P]D.8,(G;g_,^0216"$'8Uc(*l=,=$`=2-kVeg"pL&8fg8,fk
1E'+#Sc^S%tK,E5e_1*rL9'-0n[%`X9j*QL*Lg@>iEElJE[#CqJ7t0N_\p,($sh<C(J6-bHCkW4&
eBftb1_W\H$6gV6ir8,Qd)`Ds=Yd@&0an0Ff:mCZQ@d%N4pYh?@,mi=T_+1=sVpnLq_X3RtH@-J>
LpEl4?[To9&BK`:O>aVp<kf`"2>/!W"Am_H'K2k31mfA&JCd$4FZ$^"IFsG9`j^thKkiqb&IqP:#
l_(<(=O*\dGh,i[Kre`('tdLD?cFo0IeYF$?MlcZl,d6#@sYEWk3YLFbtEB7V\f_$MWqu@_VenPT
dm97,;8<%6/leMY6YTp4C0G5K3n)Df!$tU(AIt7[gt;TFac#-d5^:hc*5bfZu#\r-oW2`on^VhQE
LAu,PUep*oX7]^dPqr?nU7VMuf]EVQV%8'HbUO6:[QoP3L?O0[qc,#tquZ[GZ*`!BoB0[XcpW:qE
CiGsDu;_;4bP_rD^HH,_7#5C[rO%c0_P`/K/G9!=b1%]YWHi#&gbT30MIn)O;adQ^Mk]OAKlG2i[
Hd0o&`AYKVrRElVTq4fMEmB.$"E`bqLLN?=rgFEQWE\=jW$eQ.tRU'"(SpfU?='U'Z^=&&Z3k]9-
qs-7[ijZB<11Yf.rFNlHO'*`fX^:e?*i*Xa[;R=Z*Zp=dNug=(Nehh<bIY?9cFS@0S:hP_=3H=rf
BW>h3]b$34,/:pkFZ%bkO",U8U8isY2k_+O'*`6?9M=iV\i18gI#k"GHEN@r,C\#`HM1@`77HIQC
H,ZKO*&gAJ#rlMpshHrABbU>nH;L[o^jhcMM<5j#NIkp[5?+;bA\EAH7tiTZ<9p4$r\e9h%<rCu%
=Sh0Qpe$%4[Cn1ber==TcrZS-/(?bSs,-ZjRP'7-JD-7h:i1$ZuW]dBHa-E@R9"g=KVhQ6TON<1P
-/c#!k06Bce#%+,!$I>:o,rbq*(Gm>OU\i7nBglp8(XK`uWHPHBiC$MoXAl*Hr8JmVV$#n>850:n
qop@t;;UklN--F<HHeiT3d4U8h[(NVYDeqlgRf5*k2KEXa(/4L50&_S&aM`pIZqg9(/aLY?%q`l,
m_N6)fIL[d$6ZICO\=<aTKSBQD)1DIi5"6aZH3>&G&a6DkD7\OUI^topLn5G1h_O,)(BFdjd*lhp
Rjp32HX+X]=;DMQ3td.GB%Aat0+FGY6(M*PY/k9ZKsQc;4%3EamloK#Xi](qm6Gl^Tc'dpm,oc7:
rP.a#^soa)F0S[g[P-*-NHB.3JRlf/\A>ZpMB.6F_+M)B!0EZ"+_$EM+1O4$sW]umjYgr@TK[3<4
mr*@^OUKB`EF#:!FXnp:*eN@Gr?`H&H$:"9_"@734c=p,@K>oFk?u<Z9irS>@fQtHK!QO\MP^HT9
gthFgfq[+t1MaeWf(!<2nPt&uQ3(99Uf_oE?=p8X(9DY<EXf'^N95SR[uVk4S=7VJQE&#5*-oIM5
WP4@I,FU\TBWf`1?*#>.jo1<XNN%#\oTse=]LY#:7]Nh7H?2'7=nIX88Ed(?5TkWg*>c4'qUD^Z#
=rh0g;G5>?GBJEUh<TmQo[gd(R0fS$l:meRHJ*q+mN3\qii#2c];ZM=7qNj!cDk$>U8p#!1+A^oa
+@AsD#\#L58(W)8%ldYBTQGefI-2E[C;Bgtjl#C"X'0_a%.rrD'j1rXlkS4.e*b3"r>OKWndGA1+
m)+c1R>kuAdWit/'-2"nAis\mPhB>QC-)-B.ocL+M=.E[pOEHH:h;.tM6"Z?5TB@7hCeM5L.DFXm
GK[K+/'\6"gne:RH'q[129&B[,`a"OKeFjMDT2VH-j@A%Q[7<<r+:pl8EYqT"JrM+_)LVaT6SfU=
2k_qYITSu))!%O^$&BsT1s_d:.NeDZ>pIR=2NS/#^&Q9TYPt:D95G1(#p;EL2!mC:0i%#.sVZM7[
@Cth++j"CPo9NFal)nrNV4QX$tX1jb0,:eTf)HgnD/(/dH;jW%$Tuc&h/^"T9J*c-%g_@d)&SHZ(
1Poo0bkEsf[+ai6G*A0^fj7+M,-RRN!6V`q@W4J1ILptae46SG'WYE,1^NmQf@>l8l_TX"%M1_u4
m=15`E2Tr7Pp_=O#h0"\Z[*uY*MfDL7XPLmphi)6HF*(d%9p4OPc88`e]=Gg=N-h=Cl,Ecr:OC*A
UOT?#i^`6o+3&3a%n?gb'pPf]X:0\ElQtDMVoJ41?Vl@8Zo@-g7^Q[F<gOehXR']li?7R[PA9(&Q
WL8^qJ2YXjSOJ[B[%Tr_lrM]MoKUgm^!p$]:2_A:?1,<,ZpP>N1)d;U"gIDYf]CU^,)AJVN76h]5
?D(5C;-gqr!Ea%dgo<&a\WXZ%tg$N)TIBno%Z68.-L\F-mu5PS#?gN!+U'nVu>=9%@H_Upsc6H]l
p[E`+Iq&_jrH9%5UX^3batN+fj^Q?g_a<B=$_g@#B2=3`@-&[Atmb,E;a06dErm=C)Jc":Nj??t*
P0O!AZ0m%3B+No^Gd:dI]$sT'8o7gF=q8s'2h-N*5p[s,"3Pr0o3LWE@5?=i)1hH:kk9PlL=NlmY
4ljVILXGbXqr;Am>K"s)37h26`j3Njq[NGIHX#>9Z0k?;%I(RdO_/c>='R^HK:W#B`_@#RKAS]$h
ECl=Gu7U9:QJW6dO@:s'"?-BpOV`Ce/-D$.iEP#SU%<Tj^/Mspq8b2a1AoldDO#ED%F"YG%'fK<M
tf32dGn9)F8ZRS&s%$mGB(CUZZkJ-\'\uY.(7CVbC_I,2:+JOA])i]M=unh/VpmQ-43MS&bYtn`(
HY(*]XH;jF'S@Wss+XJ2=TQ.X]YBZcNM)%q;(hlh:Kmt73P(pPdR5i+0Vmea)tac;aR=COZKS7oP
]Aa\FmjI,Bg^h5Y`Zo/.-aqEZ\(_r4@J0&'#TR##P-8G[,].g?ak3RM,_-$@IJD'n_\oJA@pRK16
ML"^MK[JZ?fp,i!QS,?%!ALR0Nc-7-g0P1jb_D-r1Q(CdU),CFN/Y?`YtBE)B#lo&j*Y7aXmU1;W
j>.\XD.>\VCDnqOce@CVq*Dq-7i?5$F9M)?aG:J1drKqA=<!=`\OhV'j.3N<`3hm>\<Ca?Q\,iQl
s.kGH''Y@L1fIXJ!6ob+F'Z:2ILj00+?BN0gt[Hbm%.BlOcsAL8_SYT71-C`ce.obY,BCF:EHc]0
d1MbC!(I.%lN40_O8#,D1ugW@P57d7d^]$8os.p%,qaN7QXpbA1n<Gl`0<:=q.DJ=5XIl3=6Kr^k
$=b1qsqom]tGte5<gs$n=hQ&Z353m>kNRO,)1S-9]$#OTc$]<.SIPOfnqeog4=l!=aXGWoqA<=,F
47+W"o9kg^4M5.*O0=KbpUe"!^GoXHE[rFr)QidS$$On)@d<N&IVXn-CA2tdh8P7rX=WHGoUk.dA
jfHc""UM0@$ag%q8""LC]=B^:o(R"QbG2DYGF17.ZaR2d?.]&!(^XA^K$cgN-@Sdim51m*SU+8S@
H9n]F2N59)-ArhdPI%ac`1.%iMPM!.:^qgJRD<Yf:'AVk&iC12T*i%WG^+jje_"d\"%2L%E:g5`r
T8F+^:q7Rcq;182JO*K,kp[I_7CP/^:>Gh)T1`^-UEEP"kRit+=aVgRE4*tie1=sNm<;lj<$Ll8Y
uXJdoNADM0Z%otj@*n-((Zf\[Y..mqI._*.9<2-/d<a%8bNF_'?Vt#^<6tt%n#%(26l1r:;f4po$
I\Rf`nQ\(<SGT-XMSb?'gJr,FrDGYWdX5&gM$JC._?X0-H\g53`aQ)nj3+i!^UA%1%@cmH$I_.TE
K55%YZF7Q.WLZRW2-%4pPX3qk6?OeS0ZS&+)K*(.PI;IHV`A%96B0hX5bUAl]l7E=hdrlmEHEn5P
DoI"le7Q1#5iRlWChpCQ7:#iH[A`:82%O#4S9[g-;uQOt1P*K=pTF0kB'g>QW"-Qc"Wg;6>1b<!k
-oSC)4:CBbPaXklf@WfpN0=@%Dic?W'/Y&c*/l)f./!#X=0/aJpc-Il2dN60V>_W.;J,sS6r<LX4
]HCq64B$L,C4HW,REdi.iIi4J5>2E5JHiC;_k3H.43MBn]CQ'H$FQ-26SA72!j@-.qCGph>`/-t]
(sOH+Ab_%o;hc$0=,C`R9#Cs\0J&clK6RcIEHd;`jMC;mn`Ktel8;pbd?DFMOgfJn)e=]l=;E3kg
0Y\f:"<T5@Z*kKL-(C_ZA3!fqFW9p+Y2a,CY'CKQ-T0KW5`N_X/r*3@jXF[79%W8][2tu?!qhJ6?
p!"&20B?N$OqJmd\u,0,R!D&493#$!Q2oA?8)iNc3<i<<O6upWCCPf7uB^C9&18U!Z;?7msbndkL
_ro9)/*^T1pWHc3/=?o'+ca+l#pVWcS+c.mK75&`7EXlt1SQU,1rD^cO)7nn6M4eYY8Z+>?e2+*?
5KaM-GnC?B:bGi,2m?EB&f:(_<@Z7:'eQ-?ZRb<TJ>8o:hNpXKRbknq1onsrQb;GKJNn1o'b$2Kt
XGTJs<NbYX0qHjpjImh27ek!-Q#-*N?5!BWcsaZbY?;%.D)fbcklI^=9bq(W%lbMi1$$INg`*>>:
--'r<YJjCj.(SqUQ8Os$VO"H)+Y+<)p"qj3_pU?N*BG5Z;Hs2P21E.)+O<#mkg&dhXBUtKINki=W
#K0hi)5C/E\@I0E2$Niua`;IZ.EX;Yso!5mFS+U4DRU374TdcAYlF;L:'hmkZ&_Uf$2o?N\8j#<:
)`DREQ/)_Zlb-"h5ka[g,6]B5/7/u/(CD"?U9NpH["4]f&>nZRn/XBm"7(VSXc3@IPc[48S=9tGM
9;!a3JhMb/aZLrt#P.8_R.F6`qWI'Z^0,qE\'I/JjA)Mgh9h:r1P>#gJ>3CMQniOTg=Y]SV)4p:M
cERS<R"pX%>Lp)Fb.b'WRrFpoMR/es,3Mhjfi@2?Fo>Or`Ycg^H2Hn4:d[`5Q"XGSGk7>eWc%p=k
-:/RIb;G-41hNDDWQb8i/EO<?85qBg?fF6(ho`uWI4X&)06_1*"cRuc7:O93kP\6\2c_K`G>'*V\
W2a:B!"gRZFXs6V3:*=+g)Xq6"]`0/ooeH,+Ut"`qsF(KB+*]A=FHmG#54-hR^*Hra!_7+:GkK!U
Os]<!aQPRWnJns9)3f!@e-Y%nWsq2>7@h]/!1B(sSk8Z>Nt@l0!YXC/(a;4#hr(&Hscm?3Z.NP#(
"l)G3DBa;(DqRE_R)+s.sS&AjB?_9?W3AlAMG?-m]0<K.^%]Y065Po!6]dMRc9?P5_/L>>\,1_MJ
fK4gWgS`=BA_bc#1+=Z]ID8KdG9$?(Z8%b6k%"n7X?bUcP)K.Qdb<5raab56?d*;83CRJLZomd\X
LZn5QV$KD-5%Too8ubPZ.3oVgih/EhF:p5&hLd8UEH],nmJQaiqg;c5[*KkJmCZYek`6HOaim-"*
]s8]8%:CNG@pPN2-P_]m7SP>*',e/gI*2A5o@RfTPj#^A4NIWbWAOGI@4]U"#.J<T;'"#p$#*=-B
)B:-\HkKS_m"q9'"2ek&W.CC^R_gV\jS,(oSD-.S@Eh&>HZ.('!SmD\\fGf5YI^@#/g132\DTD6m
A&=Zlr^$@E-95D/Oei[Op'iiOXWOZ^nk89DtW3"eQ1?n_SZkEB=j!Q3`X^,]ee`qQW*#XSg3S%gm
WFbnM0Tk15iqNBp9pkYm2o[l\V^V<*f97WU"Wm&<)(TktA&^?BNP=s5Q[Ea]p@#hHZU*[BP^';W3
YQn![Dfn[(&1I?lQ+ejb!_+@KA.?M0"g/&O)1k,OR?Y7F^pSN4>/^s72q;O&X+>")5en'\\Hl7kj
FId18#j<m_;kN];-.U3,IVJ=dZU@=JRcdVpUI\CV8i=-D#8lULi=B#O5m3j"tAj44Rf_D6\naSPH
5.HiF(*%`jPA"M5Opn"+M\e6*hMGdH^8APWNSW2-CrO^Hs_e[g"Ho]L5:_I?\oaf<dIXi0ge>YDi
'0r0_>9iY+]71Q'FZ-IV\GI>+%H*'207rH;s;B)Lar-[YpHg&l%icP!C\?7Q]g[:Sk3[*!\Xd@%D
jrH:]q/^1E*6*8ET;o/R2iAI23$dt[mkQAh<`0Z$AA!o`3d+YVdV-:,r8sc^egsBmfp*<B^NGK4+
t@-BNN%B57Tt;e&H).#EMD=fAF"P$a#19>[KA8X@1.jU\uZ]lDIHNM?F#N)rSBJ1cu/h;;ma!MYX
oN][HeralAim^@41^@YDeYgNt)@gQS463[$6Ui`3Ct?/S4Tjd>K<Rl`]%eW>TOJfoHa4omSYN<_n
[M-)TBH!jZa@?>j^+<V9][bpaFWgg\PsMfi@3.gIB^3caWF`cDg*Xkl0i?]`*]rbJP>mspRe)u#8
?;^[('<1(Y7[Gq8ZrhFi;jXjtbo]2dt2K0XO.]0TA``+GY&e?rL.;hY:Gg^r>kr@iN.SB.+echfq
h<M1hM's^io?L..cMT`[;:KMroiY*;-ggGYlYH/QWS).g!'$e6^"oaT.(du-81%e_fe*ai<`-#=l
3f,Sai^1q7$b`9V(\dFT6V1>F?JThQ(!&J?oZ[bF6>A-EiQS[)Rjg(mu>eX5-*^#MNVNriimV^CY
!1B\u],U(637RWYG2!D8mi'5(543f@A..n$dNVMfISA'p9Te?/R't-JT4<-#6c`j:)tgDkHJ:1;,
$EaW8<pg_X)_PG5$:IH;:56b'G-Oe2$$B0FXK0^13q<K<5P0@\Nh_F\RpgkSXHUFq!NM&tD9<HF!
hYd$YQ8pJV1lDgn*iSHrn.R;0Z,;6!rDp;kCF>CC,'or!X2lb_&l5T0_N_EnEgQ31=I])t_or:;R
4PPXkK$h4#j!3'85?Ep7,+8^hDcrK,kF/J=7:3gq1n@,9!.)d;asjp5_:*TpDk*GLeY<O`fCmc-7
XFojS7>#i&.RVt-KrYg(ZPRRhLCJ0=5WD%'o;kiS=q4T4CA1Uf+4BrdriYZ2G"Uo`ZmeJjl%&G"l
5HGAqsNQ9PV)SdiRhei64&n@m^cH<o%/#"V,T*Q9=+jrZ+S9Z@QVeIe^#s&"crM0%S<'C]<P/IHG
5ZmN=$sF-E(C;(B9;AZg:<,cFei'RQ=+]lN8"M?<)<_;9B!heY_mFeRhF.*1FNAs"Temd^e!B\"5
<Xut_J"Cg/f`0TKPl\q(k6P+1r1g*`/Y%RO\3[?"pf@q[`"0q[G2K#Vl_'p\f!LdS.q@>d_N#*.W
]k%#?-^MI.PhY?+^sk*3[u[%"?D7(eN7)AfaNAKJEaGq$6cu@/<RkD#.bFUXSP?5#_O6\_\%akXo
ck=jk6.dA!L%3(H1R2RJ))NeFc5'&BKR37<a8&$n46-ijfWYcY(E*jGg&]"EcJ"*eXA@5\94)\jj
ZAL^RE(V<8O0XAdj`*<_bDsSj(FK(6j?M^$.7ae9I.+c6TF!.f1&m,"5Vk`^LQ>;I/V.$-WAgEI-
\:'(4Wl=Wbla[mIjL=B@n+9DdZ8XV0729Ci.h^RaAiU6ueXW2h,3:UU<i.EC0V]F?'Q3HZqR!&tu
rOi!PD@IEN"^Y:b0I:)lJf046GT64hJnX"E5R,(=ZMfi,_BgjpflHDLCnk?2ed9[7*X,a&]404"r
c6mTpb5VJ%&UpSin[q1^A\H.+QI/F3<m2/UgL:)#25FDL8MUXY^+s:=]-Fj:K,:r_Lj@(#RpY=t^
U`ELKFmO9>ISk53CY7->`?7Kp%QFl;NDcBXlG[h10H("_[.?mS0[RG)/bu0'mD1H*F\Ue+]j_S8-
9_nF-H3Oa];/NFhX'ICZcJ6VC?.(qSOqf?IA;5Vc5SOnVo/pg1u)m6rBJAV2d)R9=WWb2FRdj]j_
&bX]EIj]_c1ZitA:\[roCF@bb_oo8Z9`VaH^fI7qtihAF=IX2uM2IuFN0<-bR43p[r&RT%RRk;VP
3et7?Lrai]^<F/R_29D'N5u"4Niqp:'Ue)N_`^NhIP0(iWIrVe"G*fQnRPMDI0,[n!8W?kG`St`X
QZZET0c'Gfnoe3<=5t]Jm2G!c.M6\tl2po:o1E@qNgu(Hl6cH,Aic,M26<=21?!\$CW?>3:GlRmD
>l7-0@P$Ui_"#djP[iH2?%WVNh;ijOhbZ0Z5kq?k3mQam;1I$'jd#4)WU!=016;GYF>#oc,h[[,_
A(,Hm<K=,K45%proc;Z(UPt,_(cGfl^;L<:TD_[7Q8p9'6]%`F:!f[74SCa6eNXT6s22AX$4QaB&
67`e1tcafU>ZGh1SA#cAX8k"se1X!b/FCCM<"moWGR;nDe;$aYdm1pXT\AZeBC@CWSuil7H(m>Kc
fm]_ZbHbaZF3cI`3ds)$IH!TCs!m8S!'"*!u?YrG]1%&r"*7g+#P[k%Pc*`Gu)"mOpfqXL0Uj;RE
MOdXK).&iVF_tYnP1d;_*:H1o`%.i9_O1,9E?RXG@M0<._p0+AMRr]F&9Z`O=R=L@qr/(ZEVn_.*
m[tV&d[tDXJ<NHlNF:>>P"",X+-'<Lm2i^?LYmtU0W,e=WqS?Z/Vd<X,O?ODPlfNCe2k8mFT?ZX*
oo,c=pV_CA!-ndao$&N@F1QA0Q4&Oln55A$bmPVAt\BhjeC(gZJKUIFcbPXP34QnWQ>I(:EH?l;3
pUIe`Hhf\LP(gSW6njj_6o0`=-"*18<KH*-k`_'MMRB(#&K?bPO$&(IIu?$T*Y6X0-<.-g8>N-3P
iIHB6&1DI0][>[_p/+H_t?di%f(AL_I.GaI4#g=[.`>#8Kc5=Bu[9U5!P-lcm[Iob0<jG\eG%8K@
i`\3A;oQ!2nX)qMQ[MD<X2\^f#8doQSLn'c3AE#ZaIo!_pR'HZ0J<O<.n@pu@H"9=]Td`<9%Ek%m
LspDA//`B<Fi\I,nl9JZC$Vek43R]rrCdtSh^fl@n(!]e`;$bBU;1bnR[)Z'Rs%gI!_Vn)SCe!1&
P$c`hJ02D#Wtakh%mN;%*!7)!>g$YRI<h06ke=HG#34D,*4U-BVKH^-U-bA"r;P%s2>rX5b?@`..
J84,6,I#>]JVh.Dg,XaFYNXAA:!Rr:Fb_QV`PD>g^<@pc9Co9nmn>2tTNN5#;l+$!d[b2%JIJ<HE
FC[[):LR6R/VQ8PXl:g$e>KO2iU">6;<39q^aV7@fc!s,ED!ue_RF43h0c;3&1uRXolCKb'$tMt6
.sP\LF=p.JUq$3mM`'AaMa[sg:%+%gOhYKp1>OkVn4Td;%dsbh]<dQBS#IdBhqoRL/khd2c+,.^e
63Mbh;5Q#hYX,Ea@&a&'n"t>/T1RgI+1cqVg`G;F-@h_XG$G0PVk*RXH)>clWHBXFB^4RU>0(Nf[
VQuaS3nam.Gs99CSQ6GRsVZ,+C*AB[5i8k:G1>/&@8Q+b!+R\?)e9Ib:2dUkrieK/\3Fm`"3%8RN
K2`=>[q;CLQ53g3cd+l"5LPU'OgjjOX_o'OnpXZWdHPZY*F0NZ]Fp0_+'fu7!Bacc$ZeN@/1I7TJ
Q^u*.)o0j\6lQWAYa\2J4o%k`lSJPPf,c9D;q8Hq$&QbT-FN#,N_Y_#X4-!\23FnncfD*A``f^3I
,D"4_nQY1%)gg[jGrWkd1i^0XbE_gGD;AZY6EdsXAm+Ciij<6(<Aa.`WA"0uWlg;]GIrk5ZP.m4F
?_cF[^?g_AuL<*HuY@E/^m@`!oe?+)WgYP"DTY^X7s_S\ju!,SoCaJXn5NXit7H>kMm,Hk3j35gJ
q=uPn+'k?K0n,:Wqo6>t."Cao64JWn%>!lJfQ<PNi/&r+*q&/(o:eP%B'SJ&VY%0&uV<bB2hdMql
msj=@^J(]7mS<W<'dJ^HJ'me'Y9&j(NE7Pc+SWBFW:ESf07Qn3]]gQq[tC6O.mE^d;$1%%BXRLBS
P?E0N5S>4lJQ1*RNa/e[!<N57X/j>mJJkOC'9=FGUUXcNDdM[<&Fj"*shoZ^_7[4@$7;I;A)5d.s
na@5S]oE8CGK(-OX-Fa,&k\g_?=?W:)$@.#^Fg#]>?*D)l3cO"D9JAH:.8>AP&b=R5GH$`-Q2XCP
VC:ZOUAnPQjm?352i1n@$KpSGO^jb,:._U.'.s!hs\Q[9)&p\#Ng4gF7YA?<sc!;8L,R1YNH%bod
!K'NRXlknnVA75>pr=36%]+f,q'XIPPo/d!*Ba/u9UVR`'t[lP.sPWp&k[<JFGM0mM__c?RE>Qa;
&!$tn\6FiL]sW-)`%-cX%^9;AmR@D(:Ba/#k9^O2_S>E_"IUj1a(?7cX!onkS,@^KZ7UE]N80)(Q
d:=CYdo!1>-c/?:;B!t-Lp:rLbA.>d=a";-`,b(tsn<_U(Bp;_rk?4Ei4C>a`/(YVL0=Mf1b_K5u
MX<Bu"SiL7hh)M0Ch1^kh8SAOn7fr-F0_iK%Jk!i`JO4$@9mn$(Q2c$^#1TLjo(]o?)q=K82d--N
[Jc0a<"<W;2*He]Csf,;/]RY;GV2r!cgA>>unA%1&h6aIeQt(+fbG"EuX977Sni8ms&=9$F_WBVt
@"%CFA#N^ZtoY\*`2sla=CC)!ed4hD!)2"B1dE6H?#rfu&_#_+K/q&FGb+O1lPOWGC?e@E;`B@-s
(_4QM,u&_>1-)dCa#!1OZkDq6RRmAg'A?`L:%j"dtR1Fb:'^b%OaPgjLJE^H,#pT,I(m4?C1$Y<\
[gAuXYC[X2V;3mQ-)iWHIc7fl0MaK"\'U'6uh[]-CS[[bPOeYU.b(J,(?sL'`_L5=*HtEa\f01uE
$:fPa!3rh7SoMB]Lg.RN,Md"IG^E!,PDfSV]7f?+DQ#`T$WY0Ub\ABY=5s:(99!^,:ts%\J>@#W^
n`>ZIE*4J"ju]GDOp7Z-FFK*Z]&7lZ.^!oO]=e%Y5RBb6SDb+aW3nDe"\>#)S/9C=g)ZMIdr,NDf
s&&W2]O1qPBrnIeX>mAZGR=GE!c.NSM;GGk]r6kKYn%Uj%$>e<H8LY'#>E>3!UPKWb)&N6$I,[Db
.(6+Z5^(s"&q`Y)S$<8l(/c;BO\gS<pLP]ak(HeD0,?5nq_p;/BsV9_."o^R24b!OL$Stfo1i)lP
RCL.uUUlOj@Wt^hu<LR=GDgl]')Eg\@CYh:([sV^QSXf][GZiAuT/b50.eqVS;&q-s-gGfKP`4:$
CZN[lLLQL1/']t\;X`HqH0U>DDWtr-A.@@sPk]5^=[-_*ak[H4;Hp<e+Je-0NT(IhP1"ZqmT[!p?
3qSmZ7+S`$i.a/WG*s%3[n/!OWNAU%G"Z<[+0Nj:/A<2WDH#H3>lrb4qd<"F#aTCdcu9dY5\KO=Z
Q-bj,<H?:8H_UhG2'^-ZtS&XDY-&GrB#2*NJ0JeP7F_nTgfg\s$d=bp0Z&=NWj$^+sbobAUJN-Q*
LGm4i4<oSU1m17K./F')qKI2BXY)&FJI=[HIq"m[kBG;[WQPb0>L33o%0o2WG+obI(ee8NPJMsi&
fSXrd*F'V0`^.b_*=uDLCjAmaOM01#Err<o7lQ'$>A9sD8=.FiRUYQ^10,Vq1;gSGLT=V2C%oL%W
7KoBiA#*CX32DdU\N(_iXnh-9e4g=0'pQQE]?$O0&dnTX-SD766[/pl(G(.!eXD"bH,Ic'k7bC3:
il?MLoR%M_RE.>&_lIT,dHrlmPjIWBq\eZKKKX%T!t725/pSO:Uf*=3Q<67G>.13ghus7h8!fR@Y
\<l/OPQR1KHZWPMsJ+r^uo/GT<,%#P@b1IS>5'Qg<d.HL[F;PYMbL&Y)SG2crWc;I/^pY*^rM\fk
YG];=;A6MJM[rD(?EWtTu1NV<t6]<1kF5\A!tcgcrel(N%Y"QE:OoCp4<O#q=rV^es9;t#o@p9`[
e9sGW.POF<miSIEt:-DV:8lB9m>SsR:o_mNc';nLlGf=%#LD&WCc=kj/XHgnTmD$AQq6.u]YW=%-
*SrJloj2@Nr6eAYc"$CT]@IR/rA>X9hCcSAZ8%IS^-(Gld'P1%LfmL<1"20hgim")4Cm16MJm_Sm
^UlM#A"ap]jo&T)J^MP9W[[54ChLYIW9nTBs2DW]!AKKc=+Cn**tYU.nqC]!BE=BB]h)9D6KI)Up
kNU?06h5`NE%cfhMN+KU/2)^^]c:A-CEjp8emga\!Q&=F.c8\EiUK$A+)\jQEB>!0!\G*\h!!U;I
F.iNG72+q6JV1V$Yr)gK1_UO5;rKkX1XD(ik=O/NepG#E+.c;1p/rr<pG[o?IM(p"?/(ArTb]un!
c*\h6]MH+(bMT-sbT/jK'mcl^0=-2(`lRC7m9].E#!;O@E:1%7cMnS8i?[@\\`ba.V)?uN:3gE$B
fm:OX&"LKig^W;_.ee3"C[^kZJr_fu'sYgI/'[6:[q64TS&("&e`KqsGdNT1X<'$=GYBhBC$-Q17
\#EUofhohA9kt.p,h2,iU0nFbq4Y#CFDq3L*'g\S/o)V]>EhqEPLB9ESG!9JaTqCYXPII`9r[=YD
qPdDajN<kZqMbMVV3ShJ7Tl.fojUIr$TrC0,PCNH<A-1h9LHrpl2BfPuJh8A6PKc.ka3Du%Kj5bF
Q0kAHL?6rl(lrc8_>r\B!%k;9#1:u11(q8?HbQm@)\3B5-MgIA\n/#FjamRk#m9r.sCo9(jP<:9P
$QVm)tdtG\8Tj].G1f@9Zc/]2pX56dOe6jSejat1j=)[P\CSoR)P*]C`H^;ol$p'^Jq=B;PPIo(L
`Md^DXgpWrE!).j,9_RbS-P:A2>V\AS.V*)P5]l=?[,nj=1i89<.DnV=rg*!!-bM!>"85HH?,,D#
G2jaMI#AUAOULs9C[ss\KRql_2Zb\hGeT=UbSYPpPQ!o5PW3KC2nt`\MsA#,unGpkBa_U4-O9_9;
j-X9MY(p<beih)I<.n[^;lNP(Ktt[!f1XB@;XCN!]h^oCg:1hMg_!!#^[sH[K=HMHCtjMoZk/>hu
p>K1p^QaT8<Z+3_$Nh.-WP,IZ4e*Y+7k+k<nuCQ'rFNr#L9&QbKX\q[Im0s+H&adkaKC9@k9]#(K
t)!^h&-#8CX3i,mljEnp?g/SV6HA^ugnW5QV,bIi8H)3C`4el@DQVqS+RoQp)IebaH58g@J'%Pn8
`9mk9h;%p<.lsYKh?3LrFU.[ra0pl:a`NqCINgkr:2QuLT,3*\8LSA!X*7F_P>dS9KU=^D.N6_ug
FqDJE4gYVnQ?9DEKc]rDRXm=._k+HWRFtNrrChUgEnKLq9;+$L\t_bXU4eOf+&kj[<g+r@Jf"Ee@
'-B]1/ROmhT22(MITCjd05TgZ+6b;J`Ds"09/dhu$s,WD`9+F2CAp\*]%%l;4<q>Oe+]V:9jIrrA
!/pM,,_RSS$N^#I1W9Kq5&J^Z6V,V4=')9jMA`*@umUMV?mN0u''=Po\/(qf^NQAu`I(0s1@I!N^
"8_Z\CX@]kjJ+J2%e_^$FN!tk+`ig0"DmV/>o#X`(q8/D`FA7@t]:I^1H8gi8J^]]X7h/.3(Kq8h
PZY9#Ufr#_DnB=uXILY<Ff_BbZ*)HW35dB&A%Ja]\*DtUnRj#2U\s(n*QJ*Z`^.Ao_='0&Bm&s&e
<IU6MO9Qte%.YL<c7P0nW#)C"e8n>a/J`<?_!fQ5.SfQ;i9/3Ep=i<Y?2h?IBDT,Z\]m9j6%W-d[
^,n^"-oj'4&l@AQR'3e<)@-<u$![rBB0kURc<A>t*H0,7,WU:O!c(`S]DOP('lcU;erIXH\m6<Uf
,hOD4ZtA#;QH6[nMo:k2f[Gb$:$6K0qAki2)hc$QJ`J'M4KH_D35<rl':&n*Pql^Y%heW3T`"Ra_
gk6;@+KVT3F?F>bQ*VD)8C[0!e\j&r?rc$.6&i-F+gOg$lpE:8e_h(dNBN4S0Z9:d-<$;"]ZaDFo
Mke+4r9!8;hlH9PAB?_mgcf\qH^Q2YS?XT>Vps*0K,;$ZR-rcIJ_L^7h^0n!"]LqIUnqke`j0TV*
cKr+d]-Xl374I$*jV.r9>9;R:8D&Ur+!.hgI^L38XrGoO.8WU`4*8Tc/;p]hb=VV#eF$dn8X;5ZF
"B4S>+&D'%5A'mkpu<Cc@4"P[eJG`8SDkA(1akDZ$#8Tg`J'm`NS[Ng^1RW&(]m"!oRqY3DX]$d)
/,M[^NOP%QUaZZkRQSbT,/`Y/k#b$/=b9;s$j7/hCd4ir'j`U=92IGcn;\N?@:*rSnGK(<uNh5ja
Pro;]O2=lN-Q^qGu95<ou.*Nr;f:?DD>-c?s9D)Su`F4^_jLeMj)RkZqeWEV?*c[(FaA7br=3=U&
aML?"9A83&1Fj6Fb)'FSprQ-M-Zrd_%&GCP7\B4@bZX[IeIu%f(7HkR*HA_),Pcm.&$N#G(rruaI
:DK2%=E('!@PKl*cSEs"V5K2g5jTYEk*KWDXG.iXJ!R2,-<41-b>o+I$]mCYsIXK;<%eqPP4O#&\
+=D@eT&gCC))_16300)$C>t?0`n:2'9"#ridJG;O4&AA&^[hW4cs:':BOY3n?UGdh52S7Vn#EZ=]
ZPfC>=h+;/mcFQ-)$@s$6S=*]@;lK[o7B:M9396Hp'_CR!Q@@At=eSc#e.YAXTgORN?**MG400e#
3e87>DnX5pE7B43RIYtuo-";)3T;?Pf;@e$TY5U+)>!'Z%CNs>4G3e/oUq2C&;H@MMReIQoi$I-O
h;\NGFCQs9i%fJ\o66,Lh`PShXM7<`pOfu$Is_lS+57KV9>@p>(f'3M[uV<6ZdiT]E_#e6@$b'2I
b5A&4*amG21*%.OHK3NpZj7G1&/I9+te>6rc`p@q8neZ5$(0b,p-GDIde7J5?5I7P]ZYEa#]P>f!
sQfBJ,MC^r\$L,Hm&S?Yi`;cETHnnahEWSC<!rp=!s%Pa0YggWPEeFqjoFk9n!ghi%^C0AgY*gL+
<s7(0rT=5Oc+(JbnQPkO=bacgpVk)j@_]?M0_K\PN8dCC#^Wu-_l($`#I&HRl0"o,P-`<>1jEC-(
"pK84\mDb*E#-sr`=#\DZ23h2d!P?SjZ+`CU3]<]KK\+"I/$rO"QV_/J^,a=W2r1m:GcZqnjPS6#
8j0Y1XtinPPClm(+DGt:GA;KA9BohgT2%W&XEY_fPAUQW'QJB4R!.\hq8osq[nZI_c.1SIi)0r1T
*>U_"[;Fq:o<K+hJdZSSE.R[[K?F!7.P013f$c<mhR\M=u)$/,Y4b+HGkLWRLG__We%?(LqIKCk9
.<COd&qrk7KnNQ!NBnfig1,;PLS!0FXdG[AWoH(Rp`I$HY_`;U@O;3cnHPGZ2][H`lVLht7FYo91
Rna-G<9YZmHb*jQ.-);T)"0X=ds-f&=YUslQ<M"hAel9D<7>[T[)2DPV^Gf7n]^[F2CeD'0+5F@)
s0nQXp$H?=2=I/]*:hLtCgV4TPFC("f7p@s4pC1[0]jFsaZFOB2]X0nCUR07MWGh6`.1D)`C]iF1
K)tkE95$L/'C%uFK>3q=+YDcpAhg=Ah;\r@l76Z0\\sZ]H06R17K$6VM)gUjGnk#8iYJH)elLn=j
#R,>2%A2G3g!I@GSVX5-uX@:<H",Rc7a31(G:.lSP,A<BRZI$?+Y%>GelIo7a(q&*KFjGX53mAZs
SpbP`3Gh/?aO`$<RphiU8^=lWHZfPE_`Z@XL<!j)i\`:/+l.AM;de$<'BJLB'cA$dY'^.JCiLGM`
PhYnJJ*N#8_d!n$:;H9V,ODg&VJc<=-3=\m9.._goZ\qk;i3B>1SIb??W;nmf!N5\NB;\j5Ip\2o
m+Hbfl'C)lZ%hBP[C:UD_UjBCd,rmVTfODN]@P*'H4#PkkEOU0m47!%Jrpb3>Nr>A(BYRn:@rV&"
.]CFL(%ZZ71nC26aN\_0,tIcN!$rbG,X6[.j3pf)T&+60"B0]7=MImu6f0o<Wjs2h"3,stq7q1'M
NqK9&<eWBp;0Stg\gih(d"(q.!T``?+dO&k)>E@Wmt]-gUm?M6S,&L;@@`a30B[G7<uaB;OT?'H%
WXOTt`SjqLu@=qdF&E:'o3"mOsV;G^=B-hW[J/n@:n.9W0Uk`p)1L482l1>40.-@m6D,b!N!C(@H
8SY;4Omqr]uiXc6;]0_IJcXBP&H:,$c'/?*i<aiZjeeaB^g\Vt:P`Xg;QieV(JacrW?rN7MLW$Hd
2)7i_OWh5>:^+=T$d85YoI:Yku]3cQJJmGgo[u02KV;4n\5f.`b+b*#^!3"`h8.b/[G=c-gC[+;4
0&!GSh;l\OgBsP@D.[?WT\QB_q8+]ib#GXNI!EPA;D2`Qrn/?drr?p`b.B;GUXbSiJ.!FQUNcIJ_
/&As)V`"Dd*l>o<S&*QeLo<>lZa/'[\?X]Yu3B>/&bu<;5[cK*RB-QpNZkG?%D_5GcI9nql?er6/
Sn"`aVQ9&*GF1-B1]=,rR3c3'Oi'40.''"5=oI`XGeF,nN?rFQ@5IfQbK<.I@%035YJmm=A*#l0d
`[;CrNl??8F/WpA[qh=S.6k-0IUZ@%Lt2K@n+,_\uPh0"])\GDcR*cRuVQ+jY.),s+eS=Ep-'eOL
>/@F9?9Du"jQ'f2OIH&=\e#P0^?[_Y!YY4Z#(r#[U!@(t>[DeBj'oMF`:T5/0.*+Y5*\q(s]u+KT
(`%dUgKtr\pZ\_rp1H4IMSp%_AlOs:,8e?0GSZE]']e+CEJbBm8DRT7A47cpRWjZ-`J!<+1=d2O1
sS^OD[(#r_p]nPf'(AJJn\g/W;L2V*cDrAg+HeKMaQeLoj7/5qJq5\ecS?`:DhV*20^=W]YsQVa=
0$_)GENAir^eEm]h,mmG1eTDORD-j$bOuSE>]]^3jFB)`7KN?LCUJ%(f.&CKWje.aC8;U=_=u\*k
Hs[f6?I,`eu(A&['fif^'cmf8U"IrbSmZtff8gg!t',>l(q#`c%"Id>QOG8,:BJ5hu8`Y?gb@k+4
YY[^>,pRLd[Tga:Il(K[Nc4PCRQMs!YmOT>N+8Nfi)_S<<+.5X+5AXA9AR6>P8QA!&8G*/9[-97n
(#-_/D<:T/,\:Sg=<sHE9)]5Jm?'lGcb/'ge][5?@R`[1\E$3K.k:2L2fu\/6nGBdk5<..WBY3`\
rpF!2hb_5:#FhN[&%H<@?!]K,EhM)QW]ZdMqa1,(rjB4F+IZ`p!%#E.dL3TW5!6@U&44.](^e^6)
+'7an6_>BrjaM[HIG.@aVRVh@aqsrc0`^21NVXG1(Ap`Q9(]^A(iT2&bN37OU[5ms<!7lFP!&kjL
&Sbo,>RUELn-J%R&,b5UsHA#&WJ4"7l*mr#J#I9spIis^7'd&Z1Zo!mnRqeiH[-?"#gf,6UA(0#)
0"MI_/]I#0)\'%e;)AB8lrr?kgL@9Z;*hrJdi-%S)m\PQG$#Rt.>$`r'kp5Tl_B#)tm*4F3?F[og
mGekk<P/Z_<4M#^KKIM\?F!l-4YM_4]ifO5>KE^p(s5qFXh_!;<r?Kbf5nJe3747m/Qhpd(og6=A
5Y(e%so.1h\sW03rM_JV_)hVbYIm]I.%'T?/7t=M68XY2gr>J7GalK_8"NJ^W$?g8Jbu#e@>50)p
WWXm[j\"_ogmkZ'lJQ]X%%_A1k!Lkjb%1]626Y/1T)5]YA:1OB1[P3@IEe<rCmim4FZX?b29`n`(
;S4SIc7"mEVAPSF)e@"/+UVXMeC;+h;SS;bj0jRnEk,?q5um<*l*HH):M<KWNRDVFiR%@,bqQc:/
DgK-]VpV`;qf\_W#,,_1@<`]GbHt6(BPc.4shZ6,G7h2bHc;h*N40$)<lb*0l"Zjf,6:'%d5CP+$
/=A6])UX\&dnT0fHcf![I5T/s3(@Mp88,pe*D>qk1cgH(5$qDRKsRA!AIfR1'g6f[rr@,rIr)%I:
p$(cO&^)9oj$Mi%#Me>QMtjC?*f</[-nbn;<OPfKUAL^?/\iMFTtB5GRcNu;H["5el1gpf:81*25
9og,,l-R[D.[LU*AL@`X#>fiZO^oE6:#[Oc-Pk7->uJUd&TJI[>:?UE<J3_#BpLPWa4QQ'tBV%nY
prQG6#k!&u>Q>s?VWi%Y*[Bpq@))&A\sIF9=NosZ_4]J\ci-r7@Np9/bRAcC\mN1X!lGPYW>,fnC
m7T">)'L)*sidEGOpNjp<IC;V;Gesu,+c3eOLEsiu:"*f]iInh'iaf+1X@U3nHA#aGIWBOPhW[Y9
8#$.AlP_>3'oY]hDf,9k"QW$Xqsr^8#t7fdnU]-2U\_W6V:5Z[><LuCeDWVfS6sL*#AO(',Suu_&
#d#nZSWG'YA\9&o2r[];PjIFpQ^]s5poP'(``KNA=L3=XPGmANNMX7=ONgS;)sl%lN_WL:j[lWek
g\o8oNQ(4i%u53$P0*hYlGWWjZu!l/gt'k#gI/pi:D;B8SfZ[Aa?-@8nL1^,16"mcVUi!,!QU;C-
2M=tlm]][^CHUiqocC7UWQQW2=/HJaZoW?O*QKspP-K.bhL%#^jcU;W=,ro;+dn!E9MHLmmli=<r
t\6"Qo1tbKDC>Un"7/DL87d0jP3EX-O;Bup6O]SWF9I+HYf/q,I!-_#*Nf2<Z_+uJr,Wek5HoTq!
p2r%Q287\4Mq"\fC@:Ef!1\-$]->C+IXGZp<\b?u)Z$ru-S^_%]$*`W&FGjV36=p&pHFuX3)eZ>G
H#egNYon,IoTNjc&A6=jRniUZrFEq$7/?6kO!OT<A)F>]VO#aP=`eWA>KV=d%40pTr/MLraidu\N
>p2Q(dLn:+Ua(>[Ai#c_qJ^/QFb>2H!n[PDh2)bc<t0O)kV7nRme9l3C8hLjQp`44q#iNb1*8e>N
i-#gQ@/\Ll0<FG<"R[_$a6]u^=KD<$7%C2Xne4$gK\nVTJ%n`__5b"(]?]HpJ_Z++]ANhCeOXfJM
(QVM@'BQ_:^7:)_,41i\AprUW=MXV.gb'"n4HV2?qq6ET-esq;;[g<R[>O(m==qrDZ>uii+VtE4/
f9?.0&mK:<a6f*RqN-koNp_k^/uoa1%_57u4-AXsC9hoF8YfJoGH;:FoqtYr?bKJc\t-;E^$2V\3
7LE%^DP@gH_^54auC2k`haP-,]^Q[`J84Z,fHf$Rl(50;9hC$PIY^T<j8\ebR<'lc)&f^jQ_X$q8
OW++g3bS;/@93mNX'hakAhX`hkgO.a*$C<XEs*0)T8:Td^c5k*d8k*8iQkot5W#qP_\Wke0$lQ-4
9H8ei<B\.1`'pLUAb@E(Vk*cCuXXE_jGb:$+=,sujOC#/<;gA7p`A\N'u;L=Dh!8l5":\hqTMBG-
l/EQm3D-C6r'l9[9=]Fm]8fgFmQ#:-F4T15'8tmLW]PmpMQ2ST&/6,7ENB]DN6!KN:on]l7<.\!:
idPYl?T7M+cmA:ZpQUTFGHI:V(.hmO@SJ[\Kum7'X2[]EkB3&KUYA=`HX%RX.P\;)j8q+5ptq#"L
,8[h^A06?q4m*k5>Q#K=ZbJ:d[>#lXtd[LTGMa-g]59ne?LPN@Xmit>7a00^+SpM5OphU6`G:2BZ
P?5lPQ\)859$Y`pS'/`6qqCr8JAL+6Oh$gbGkc@M7MV8aV<Heb[NloD\eN;[qQ')qV/$_+Nk\Efm
$XN#h1?EUP52BWC@nrc+gij.0Z87Qac0(cm)&I3n2/)A-cb!9,rq0s>sT??0[Z!*'QAgOaI4lKpn
3OB]<7rS$jnOP[A(pu7)+_&B-L!"K%?E-ohKD=`EB9%/`!<?(U]e:u73icdM/Is>VUZEGY25-FX5
<KVFL.-k3F7&!ed#YsHQ[:+'"aXNQ9k>U[.N9?/?:HW-EBh`!SaeUJbon0bmj!)_A:W@;>%<'N*%
66?t=pu3Hht)ZGV<ra,P+$W2(e(/R9:'\f#C6W3?LsZX,\fI;bfQmG^#d5R9@Ok16^3Pu\-%6#Et
<>3Q*..j^KQ8bN.sA@3tb-='R#kYoc8g&A7!#m!_K/1@fBs/:Yu@;@k^\3Mo4F#`;A^^F_^9`,YC
/V_1mp'[GF#Vp8qZsooYcO^DPepQe*)`FGY1gYNI7S2\Gg?IHa6#ftTLl<j6&hb9GugDIntAY&Rr
))]TN?R<)J(@U7D58)=PF9s)"O(JAsBS42k.7_RbK-r9:DQXm/:'qmegjEd-XO1laL3hi[](/t0%
Nd8_[&(>I5Ma#4igqIg/ZKk&BLUt%?p#d5&HcW6/$[78g!gq]j9?&Yf<>8[:Z0T$JroS9bXTeF1h
GKb>OP9!5@g+q<:<W?kCEk!<IXPZt9rd%5eg-[+ICsm=%Zf-i$NhR%&s)k/,XL=.T"O(-mJn+"*$
>:7,=1c$('leknWW7[\m@O/Ze5$;m)t:NPh;H'X!]f:.C]10I9rSbqX<c_r5hK4:`Nq0jSQIq^KA
_F+N,5^+l<s1hRo6GSXkrB-a+B&HE1!I]ND`tKlD0H$enX#AdHPuM/N]NrQDXQ1u1"GS`H&F,^f-
7^/9s@lt_Qi)nHGrRdl<1-Ym)\Le<c9rpch^2(tdXnPt-D]']Z,>&k:CD9:u$)Gu;p+dY=JYp@=^
0'JSlc3B>.laDrI=PZ`W/ueY]6ohCKq%u\o!U%;u0[<V]8FDpRNa)cLN1"BP]?3FLdqO8O"Q^uPr
B>8KOS!"f^Hu=BS[&l.]pGhK4#Bi@l.Leg,p0/!YMY:08mSb$,T"SB#(!TH+c,ZVPbc80!ri!U_4
9u*MUkelb`ta!!.hpH6:-ClgQ0hKP):tnk)4!tbNf4/6(mn+ZD?$6^\5/IF]ljCmKj&hI!8BK1fl
c\87`XDWsXkG6tW?,gEV.%IIJ9a.0tepnUbMrP4;(baJIF(UD3"3QZ![k\rd;'j';,+1.mGEh7:+
U3Hj3^jEf$sW/4<:h??ZU-HLk"o0m9Ubj;BKc1oBPV>5PkEUJmK3;qhlpZ1jq:X+]78MtYN#>.mV
+mG/Lo-1FL`YCL)CIhe%@jW5DgF(d\.PFCMZj14;J^d73h4G0,pVc$WZc38#ITjgUlF"832]Aj#:
IiE&>sckd:>7OP14T<(Qg0^0/,\DEWA0SL_i8Ye8Td6k>4(J<1cjI+<5/,\nm_PoO%t*ud.hT.!n
kV6O^`RNUYB-N[od&l\&"5`c36j]f%h1j-,R]+.DVa.YG+^;?.;U1?j&(LP(0t$3FGLm.+a*Akp%
WH4ic,gLB5#A!'T0O=XHi=M$TR!*\"qTQW(JqT,$N*W6&%q="2DuhVojU*Y%3,S\7*Ub$R+,EYI`
F7S83e3BtjACe'BUJU[9:-[S4>1F"..^0YhC*i*3#W."8roi%hq`3&&#`\EApIsc.(T%gr*^H'V#
M*ko-pr;ah%dr'6.-TD&M+c$OhSlY6rr>$?lTg`WqJ1c*FZS?qP$Q]W%FgY):SmT_CUQ/=#<Q,2b
GY!hWSZFa@]69HX:ES#S>"Y[m/"=uek;piAe::u8>"1$P_(:46p,Bp3Ieqh4!i[)'WlQs&IAkGQF
p5fMc+%O`a#?Lc>kt:8P10Yp>D23Ci1M]o^S25&XBP>nVPf^9'V=ZEMY%j(8#25l,h[(qeuZa*rc
5?+8ZOkFMBBkIUnr!lT$-;M'89KX03B0l7Ac<2>=oPl#\0"=VF?WXNcsn>bM_#r@MX+5rl>]`Y)H
*(0[("_$`%#p9%FZ6A'Obr8U%C?YlI3F5?fd:2&t\8Va20Z\)C0b3-Hb^a+s2;R>C"nIdrW:9\or
M`CGXmdd'-N-h>@rn1VArr@qAc:iaMQ!(39%c9LZ'kpC&T,@=VNofFa<f^SbjeFW?r0[d[9j#*ap
sr2L]&VX#:3X=RN'0c22RL,^JLPJ&>4UU9Lb[_H5<5Fe--O:4:5e^j_!#DR32nt1Jiu\KbXr@b6p
"u6(%_C_l,jej1EQ!e%7N^arr>'S0%X89']UtogI,RZ)&^QoeZ'9*U<S_,?LmXO8GsmlF,jNoU@(
-T$-8s0dW'heS=8,LgG:lLceB`eW]]1EnID?qPo"Z[P`Gb*4l7+'m^$/DE_r5A;Ye_"%]ae6C$#?
phZ;IIl.tY@?VG1/gES'uV;2LY!/srN1S+'<i'20SN1[<m`k"\<A>)rWPq0\(Ol'.S-h2c9Ufukc
I-.^`h07pRk'FA_*?I!LPdIJDH,^5t2o38p$$rAVI:kZlnrnMPn[j[g[!hEcW)h:f=@sdf2j<'G>
.e]ZT8Wk)XSZIAX89ho_\FBK8U2ql:+99!X*'9TgFG7eGSZJ[-[IG'17^<YM%&PH!4O#6=U0!=Gb
+32Cr#B`7<;^BN,G,err<q_*hoR%UR\I6IG$$E\s3;in@?0OMP@ar5>S:s/>_+^iu..3:jM'6A2-
ZS]MXM7?^cF!-W@V")tA[Mh1!XGjFU0Mn'3a;AVk:khX)O-BFS:NYUq.riOgn.Q8Em._5Vi^Qb'h
LYlqk<5>p&MfE7:oqgWF,Wa!N&Lau7<MT+u9RV8EWo8(`YVY??DQJGSU9WMK)=i=(e?YikPK0eZ$
iuiRHrpnL(Cm#Atc2#XA!.MUC^#SkB"B4!<2caHi1#U%=qOJ/L`g!5DHYZmGaGd`ph,WN94.E1&?
TCS^.p3.NVYJhL@:p.`;[Bb5Ha5#h.Yg"oS@*^QhAJu#LoYCqpJW$f\]d^e)T:pX5^*t.S0WQeWM
K9>Dp0DB[4u7BOr(<l6g?aUB:I,NrrC7-k(=ePHC]XjCOhiqB:i<jleK$,IUm4_-h3-(oj]DUm8'
71e?,2jSD?r73dutVPDe69n@;pW.8EaBiK/+hB(#g>1'omE=0/lGX!:L%k9rJA>iN>PEc1EZ@ijg
2&Du@kCg2tKHQVPa;HAZXi>dBiP$1AEUCU/bl+Or<X_),PW)B.NNm:OL$:M2t?ti5jG@!'p;,YrZ
"pL0b+hh29P[touR<VLe)<s_QW+2nBl856Ik3k0O4#jA!9k.kmHGU1krnnQCVQpU;Psn?pjN`?9.
es#53j=EapKPmhb,FtRJKnJT&rH7ZhBu[B6QP#Z\.b-QauX";QdJIgn>)eL`h/cc)O1r=N./.U4i
J3'a:&+<A4s8KpPmU:AKIm"dj*TtZFm\o&f,Z>^%6d17W2,UGZqo,O4,l,rpWLXNg9-uRa,DDqL%
4:6u2$HhH-`^0oa86g>@jME`3fO)n\bW:?I>h=[`r8@oEE?^GtmBV?@M^qR&e6^'3qR%iFhMN'J-
@ecX*i(Yho-QK$%-_66i"H:,+rP:4'\-#jsSHoV7**[.KGp:6JrQan7)M[T`skc<\4N(QT9ZG:kL
!@9^B,ZS.nZB-`oiu<jL7XFdfopV>/?^&JqH1oFW=!kZ70!'U3[dO5>m>YCuk7NNUl;J70!TR%1V
;4q;4Z-H>`7&k,7VZ@%CV1%s4(8fnBo:Q,MM/MqEU+D7]6r0Ib'g`uQlTIMX4-;\Cn!<AZ]])%?>
IhK]!`VIB%IF'S[2>f>/2Z\LM^)^D!OY1kVI?]P7[/c=Yc[+imWJ.Fa@j8jO]_3Z_t$iBm8Sbd@=
`1h8///@S11p'$UH.BkWGL&pnGqbZ@lB36j9V[>mqCVt1BK:@[@(do%efa2MErS>?]@_GsU])X>I
]Ls#2R=Ik's9BM5j>P$_#Vqk1sg)#6dUNJW%PfV%1^)TZ-PbLB1S^U_5:-Y'sKil=Mk9:Rj=pK(Q
K=52*8#b6_2*lGVhjUXY]Ct>-UCNfD<CF!4qdf^,&<i>idnhmrasl11XB:;KiPXLEJP51d-ddOIU
lhhqV^^*Z7?hM$m?amVddK61d6!te7U2?6TeFV7?7QSW+SS<h8Mqj0]NOoSLpBOdq>94XnIpC/S7
Xk)b##KM)WOR1rrCj>0rofp*IQoo)8j"j7i-R21GK]TGJ>Iql:^eT;Q9K0)+mq"oZ7<D6h:)%=VL
45`>+Ur:3#Pf=c7K(M(A%ZU4R2TYAANj\9C@8gVug9/uN@48+V7WCuI9M"KdUN)@`<<:p@T7OK7X
s&:!s8(k/(?n8=Qu4\;k.\<26KA+\W:>A#(BVmhU*-0q_'ja#IjcUpEkrr@4Gn8X.E4/FmHTPP8n
3Cgin\4?>,P6o*+#te.%LFL\6DtA9l]F3m.'CP+MgLe0l+3lASOHb>3`H+L-9KJ(D#!Qr=W-7p>l
:@/KVbTZSPXD^4.c$WW<Lm?##E1Hq2C1)'.FpthQau%UTZ(Y=7f9f^FNBR6=A\q0kurOtaG!7HEU
`m6FadO#D'ofcE^*+BoNsQ&4?AV)d_r1Nc;"bV[827U3dMF"'9/^.k/qd,I4hDolFiUl:e#r4gO_
IM0irXW;!(Vg\h]CMpEn0@"J.7SIdfB@rE&+VaWU6f?bQSG,F\L6H-.]fG4AgJ)5N-<pR"F7]6O"
#U\\olS-jJRZAi]PF%<2CY[>+:aH+kIc:ODIXngNQ7$*sFSfFhGf0:.sc6qnjq8J[uMn:_pZGDX6
Y8K65ff_MFk9+p9Hl0%imu(\Z3#rE[X$h6in99&UQefO>Ni@-jJa*ZGOB,64^\3:Ur*mX(D.4$/j
`45`M(TV&`Lg,?6C2jB6.ETr'3"LblFiBZRW%YQ.<Nk&a^b5>L4)ghCI/;UfF2n*qA\16Il<C"gb
/R^OV-t>IU4ahXhY$KJ^`077`3'Kk6]YJ6;1MX4djD@pOR+^r+#BX7r'nNQ])<C;K53>5AWQO:1%
8)XGd/U@Z>k>O&Dg#'*Fsl8:_BCmi"dM&(J71c>l@_^Ab*V4QFFt9=7P4(-7,)K3g>i!DV19k;p@
[n*Q'*N.KZ5=/d/,@dfA:A#]=*`eDN+AOO,Abc5\B]bhb$LrT+)3lV>gZ&d5V6ahMU)ZK6C;M8kN
D=_u^FW(`SHtTC_(6%A?aR-9?Eq0/0%F[Ao89C][O&"La!J/WIp7LLqeEpiDlVr7@ol$bYH7[/_6
Pd`*(/cM*P^?c0N0$7onm8"2YAk6"G2l1l3g)OIa27.F*n#g6SopO:=%aIpI,u@U7K"D#\C'Si1n
n<7c=oGi7[)[n:=>MoA5i5-;MeaAX?CjY"JVDU+<^%.3/i]5DHu+'?^(9jjh46.J)+bdqdeY6ZEj
8c0L^>+P.YCb9Y6'b?3MUsq^V8`@l\_>D"1ZQ5,FDlf@KKgI:KKC;7aV*42-@%j45'6iA*R[C6*V
eR4]MmQkcs:794<kNTk]0G7W+,gC]MXG.5*oLreP)q8JV+mu#=>NL]!lQSlhXJOO7`13/\`M]<D$
lF^Zje>eLZm!6CJG-^P02sVM?$J<)`Itei]6tC]*O9?q.nmhfi8-4j9.3=*pV@el5^?d-ESOUq#b
JocX(sAkkUDcMF%83`8TQ\;bEja1f1[k1SQ!-%J:T(H`*oAS;\g,PkJQ=RcG'W6>O(:tcj2GPhUL
M1lp[:%^..18DU86*;,VLr3<CpN&/5'NNn4&$,\&#3)<dd:GMgod^lSFfko(k(O`_.0bLH#p@![j
KCHDNH(i]VUI!ug>Y.LU"p5n518]jOc6e4R(9P4:sk^IAMaQg)ugi7q6(?sXjFFS)"2\&L7NcBRK
]prM4cQ)E`d<VL:njA-U--Q,@Lo\%M@qaj7L.H*!Q8%2F'mH57WCh.I33#+Nsb$].Ymi+oJodAcE
-Y#Bln@Pbb4\r>JO!Y$A0)LX[V\jr4&S/Q6ata@ok)\!:YDqT[mJrSnHsl6`O)<WVl\m[pa*YWV?
'6,.\PSj@9:+1VmgjC%\>PC4PFfBa19WQ8Jl1nsV+CUd\FG8CYG@,qOiI&s2Z#-+""i)NUdT.Wo9
/R4L+)Xg\B5nken"1PesFZ5>FSMdK)mBaTH86'AGU_ZDF7cI$["MDYNf!VQfq44)JT%:c7\PG:(,
WR/QBuB?rsM,jn-rIh`3Wo1I77:miV'?8BMNS$GL_6[R=PG75&m+0U,A446Z4Jj>r,WC+O/PTA(<
[(2uih'_i-ufJ^_sf"YrR#RV:?0+!Wm,H7XcHhEg]/op5KZ*)Q"R?U-)9lOD@SOMe<Fa?^saX"@e
PD%@IZL2LP>[p4"a\_6;':9f547HkI8dLjsqlVNn;L;s=YMV1X_E6J9lN)i/0l<%C'p=o?"c'`sq
i]d,#L6K:b_WGr,[*Ac#@FW:$F2X"gJtV.h\_Jb9:4D;4QCW^?1%3(J8.qCV/e].Z6^pKASV,4Pl
C`^$.#FZWH51PIl7O2qF7kt?tM/ha%4nTN<qE*KYMn;0k<1Kn3u_-Dg2[RZCr',]I.LDf2e\`4)]
Bk>E%B1X,X[$<qU,OHlTX**haH9;,u]6=/S#RX0Z"LF.9/XRFM)-mrm`5As9aC)WLkXP3PV"KOQd
pmgiKuFJM<cER!)7!1`6?=Z+Elj.+P"`$B7#!-81U28>s!D2i,`YA=RNr+?knqXp:+gh<IA,G:R$
Z<NO3]$!I?0`8iQS=;BZB-\i;hW!E!`];P=6AA>rAd;a&;#JE_mE(mu!1bO`6tIFYAan^*XU`ghV
baS"jIg^F?)_*)0KO[o3?:,&8L_u#C!@NZP\&F0fUD3tOSm*plj^r4+f@'G@Q%LN^E78!o]@6OO6
'.$gEW6RKT=/IXV:A7deM*CE0jRerr?r[f(.2c\U?]2)lL<n(""=j0V'6!:NO1*eERQdjMoof9??
e''((XDT*H2VeESgWlB7Zd0I\HU(=HR@7W$,.$#auI=NlsTKSXV[+]=8N.!T%K!^W#iW`7p5<dc*
WBlRo^c>`E3'#(4:`:>uRpq9f=9'T"t[(t&R^P-/5;47FCY]YF,FaC@YO@[rThK(3JIse>bpJ_EU
IsDW.G0cl7^Ji21,$c3MO42A>/^QOh'kYQNCH7ENW=;;=3e%Gq`4PlF2ULVW6_sB9Jk2!o4u8"c#
B35rZ`ETO;F(@(EgRkel,],Z?W"o(9Uo=LqKJoF(rgcJh4s&_o7Nu3`+:1>D*,f=V&cqTMW!dLOQ
s`gpCj'h;4hu/4J$\a;Ok=@0ULZsKc[%G6I+7Unjb[]d%3D\ot&PbdMn6!o8n9hc&kZAMVld7V`C
UHY8gS=^8s.40M5f2^$**97L3.#GfJ>.(R$Hb?shBoc./('HfGIuN;GlR*'?6d3`D1FeJ5POC&kk
jc$\T?DXiQ9AJH30i#`i'pV:fK?T9>rW7C_,-d$]0`+<..U7KVLe[T2?W!:is.I-p5S57D9<Z[2,
]iF6F[LW)%mAD/L4>kFSDYCa(<P-eZlM8[b<u_/>)'83P0mhSeik=8R>tSrlo9r7e0E,u#UhSipr
0MtQbi?8%Ge\LDeir\t^<-oDr6Ej^qnVEN;PW&c[+8lj="(77k;6`tLC`MK=g*boLT;d!j(ktUeC
Q3uR1IL[W5!N]H,VmRFAK#e3fTWJKS\1UaIPZ+<bJHKq>9)m`_kNNHE;=i=*#Z*DY&`K<7RrB_>/
l2@!XTo(*ciB($b0_DgF&-N9\n'X]=%V5;_eU&+tM.(1[rRp.PUTflmP?`TJo;+m32mStbb)c3Hj
[oa("^4Q@)b)1DEq0BOan]%(<2Wme]3_QOHHp?hjaakB^InXNBURu4'te7'O\;;OoY17%4077->_
[X%%:oFh0(N-Fe(m=(IOf3&>hQ&Y*JX@@q=%e]@^VrGP>\DQrc$]-!3iDjH3g9K9"PUAgI$J-rG[
W9->CA7$/lP`X3hXET3\QcBCUrh]cq=0Pr&HkSWY8K:&#&LnK>RoD>T[I,sg<4SL\T5a9E(BF=51
,Y:^+5.PnM0J.)(.OdJ__rP^N],o.<`DQP=LSnIiHj)m?"`g=7K*b`kHCQiqH(56\]WT<_7S^IOD
&<"@n"*2#dOje"@t]m(fSgT%.\jG_K(1anO(fqO#5e]/`Q6Qf=dhPG3K^b:kt@WM\=+_i8'0gjC!
[0?(e.\Q.H&rrD_S8+dq9i*5e7d[i=@Uo\ciI$rQ=WU?!P>EANHlG0]?;sOO_/>Up_P[_i.pCbbi
qP/$UUZT,2r['$54E5pTcK3h>EEKJW^D4bo.g4IF7c/KdeNZcmQlU/6l.UPO`cOT\P[]*4MdOJL>
qW)mlL*hf/M="Xb:eJMW]b&c7ES;1#C3/tA(@a8,6$9=)>952;35Y-=/dUh7=2!M`^/BQW'.KWqF
tlEX5DHS[M<.@Oe:YqrgEE$M*cI/bOJ8;ER\dCmfT5m!DRQG$$Bh[V>BeRd/6OAh&]9NV:U[Kmd:
#SS>+>.kO2aJOd,aN`G7IYf7@RMY;g%eJ3<\g"/8N3P;TZ#_\r1O6&`A"8TgW_Hu\<hr(RL&G\[D
YGATWJmd`T2`=Ck108rEk8Rj2NT2%'&oVR50ne`AOcJg36pr8QPO<7MErrCOfM'oaE*4WbU$%M$Y
pe2p,+,Ec9(f%%n7cV,3EqMKeQYs<*1!dei;3m+#0WqS<EbT#:o7gONW\7W6Z^:[)5*/-mV_Rt,R
RnYYk4m"YenYhXDmoJ^l!JS^:G3UKJ(AAOf?3WJrrB0S!nZSE_"8r(M\)5'icQUB"kTj;5L^o,rf
p4pc>!8%&sB"[!+#nOe0J1uf(SXK.-D"g(Xp+O=B,PirF\_VMV-M]`Wo"+rrE'!8D]>A"S?&B#K`
L^io(1aUp\IY10nh<SP<(tjq3'<XeY0,]$/8KcImKW.lfgW,4#!pNd`0-K,Jl*rT+`S9P(Xhiq6j
VFr'Ohk;Q&;dHO-4<VDfDlU.YM`4)!i@9]Q^Pt7$5q9Fg"HC^f;#7bc;83u@jZQ8s0,2P<>N"1iM
WZkKf*M#W'NC:*5cc!l$6&i#$k7:M]`8\XlG.#<lo]F,,B)_j#FNNqVJ>)Fu3.L2oAJn#q.!(>r,
QVV,Pumk@>XPG%KsE.G#>glXK9K;sl2E.cFK.u[h\^#$[dGB>?*4E.7';[%N=90m`$]t.kjCBkr:
t)$6bD)tZC3!T$b]jO7+M,b$bkDJM;)/dba%T;4lH7ZpMH"3p;#ma=1>(E`c4a8I=D"'o8ccQ2]A
eYbk31*Mr0j/eVP%uHtrc@Vc1+gS?^WdW(q+M.WFQ,+>1]Z,LnGs;uL,FF"Ek;=JBQN-Qj$Fot6m
%;X5VamQJ;U4aW/jVk62f!Er,GYbfN]Cjk;<;.JG<:Qm+s)t)=Z$dQaG,U*0MK+&Qqe/A3U=#F,)
?Z&!g4%enrpqZ0Vrgsf[S6er<CYHuS2IrDUNk]ALr6g#,dQX:^#+^3TKVIo)/gtJrkJN64)t%&DM
T&\Tr:r=fdBqUl(pYai7'gJ38.$I<E/?!o[H@S)c(5H9*Rnf(`Ode\UCVGrL0*[AJ<B:/A_@gN9>
KO'FoMG$`s"]FJ0f]HZSZ5kr>&J)bM)ZqrOcNrE?-)Qn/;;9,lmGpo[%<:P.We<&_-g6]Zd*GRiD
f.<F(c[E-\NdQK2DCOXYL?fAAIA]3G\2\$rr#g>`]>[(!VD4W,Zj#5bmL+3]QlG?^uj'j5\B,@*k
jm!^`dWVWcS+jZ>s3Q6??/Nu=,ggl9#rNM0&pIla>rc%NJIUr?0OfoWXAX3#&4jfM)ece)i2es#2
lQe$FFrYhJQA^[)N(-)TVH/]Zgni"pU">o/9REVila8O?EtFRIrF2<ambMd`fo$]gKp!]geCX5Pl
Lg1.YahH!%WdY[M7rmkHJRYnSQSR5dl64[_!HOi3CGp"Gc`tdb`>GfUrlIBnXT+6Od5XuQVj7Q._
T@$b$2Q^*pto`A>:a4m)HWrlN*>C49#:@U>1Y\>]7`=1H!MIJS@a<.-PS+UdtUE7$U=?j*+<t*jU
g@_t>5pe<6hjg!A+40bTKAc'jk;;Ye1LK(S#^?r6BIr\IX(2K-SB?8A9NY@"_0d8b\11Wt,b'cif
U)l\2SFc"rQ@m<#\c6T5Q-/JWg!Ug(dd@5OTk6+iRrc/Fa\TooC^Pp*h'u,^M?+<afb>q]f=<mld
qdrC3@W&-[DgP2gPNdO'CXQA]P5bO%9J,/lR(Ls0%:T7lk2/3RbGWjJgF1-=%L]9Wi#CO24^pPK_
f_usYP>mF2o+#1h@&BH)6\P0AM=*7`aUjTXF'2[/Tj;V><f1rGR]&YBf.M]k-IW:n'2Ho^#.^+2.
msa)Z5^.gGG4]kP3#,&c-q<-$27W1gOI(GSF/N>[&h=$mg)e]B@`GmQK1@@>it-(0&^iZ9C3tl5J
R;&t^D:$_mVVD-g0&^bN%KisSk`][$Vf([OB`i&i2T54rM*KOgXlX+eNF:XoW*ip-+'il\5!(SMZ
&D>sU+Jp'5Ql^o"'off>MVRaon(edC]E14oC/d6Q[4/rjV:]CDnDCKo<[>HETHGZp2<Ss=>ER()u
Ut)oGitYo5f8MWf!4_o63@;Dl`b<JASG@/-GY]F,FS1+u%]<&IH6"D+Mdm_56r3_Vq83ilmd9hOR
Z%1Gdg7.`da`ZV#+oa*[R3b"D`NJG%sM6r;HNNRnaAV*o[!1/LZ3dm/uMdR`a.(aek9PDXII\bAL
':e&mHr!H+ME$Y-1d$9h<E`%\1m,N4CRRjTM,HaG]9,B6)E&_p2eS]PHP05OLVH^]!!k(&mTec88
\L/P%sN-TOBdP)N_FDr6L*!j]-$)DnqZ<Q1GJIHpQ(Hsr!,U=1'lh*X]=GSj;a?F)/oCjE4//k&?
^IBd+QE>#rE'I%_c.XCpLF<.eI^#4ZGMm33]N21d5;!`m!c1-#Yl-H6HEfY);!4amhT1OJi@$c4.
jIsB%@>qtgW)7ro3(E@=p@0*9d)cOA\Ttd/9)GB^eS?Pfdd#Tp7O%p]B,f'"#,Wj.nWU>_i:huf2
X::]"!Cml/cPN%#@gKs>/a&u:-;pAPOd9,OcoGo$hK'F2NP\3FX7BW,>IO;h+:8?eX';'V6i(.`^
"+kbHX'7(hCc-r5e6_hQh'DC>2PSR?qu8Q(E2Mf%'L/CQ4]-Ws!\;c=JHlFc$r8"X2La;C#&ln4$
IQ<D1Ue:Pu+GCC)U:0$pd,S1;Dh_4Pe_1J`=S(T+lY3\bJ(oqr]s1u=4K`W_6Vodm^Y"'?>9G9Hd
rjSo36m%I_EZfSL-c;@??e60B?3/)2>Is<L.grR41rbrmelFCi',GI?Ll-tlee=L\8O%tF"GK1Q-
G5d>"]<a;A]>(n<2g*%e/(YW:a,n&bN%2p:AUa`>85;q;T6WaA#rfP#KRhcJY75Cg>t+#O]hKI`/
@o`oqd=)9.ZssT=gd)^OS`b3=@W#3].kTgdJ(MM[Xp-di56pm!531Gil%Bn;"!Wb;gXpC%el1`cX
Eii.&Q$NG1NuhCoU19I.;lVIJ:FYO.;sWPdpIm`Kr<m$X1\pT&-;DNkbD/7O`ZT>e'j,dP#ZAFa+
DE4RDTm1Cs2gm5!K>_L<l9%fW*ZqNHKSK15u7NEuPaS3<2d[Q=<JlB:C'KT2T@5;T;7Ep:P<RdQ4
r"7h(BP8_tQS4!ZFrNd,079>IYG4d*'[(giTm;%p?JmH\5FlB;l]5?ur`bqS6?Iu.mm+q_P;/_ft
_&/u+X^QNr^Sb`%/bjh4ot&Ts&d?;+;Q;DnMQ*\QH6$$,j[9R4=PYXISOq%T&I:]LpIkJ2UEs59p
W_V3cUr6T^hQ6lF8TN7)^RL6Snje2e;mD9lm:D4o7.`tP.:+O7NlMPn`LYTVO"pUG5#d,[eRKE!1
V`u$=<MJl,GNb?!-`V]p`t8AK%bd"`E3]n1W;UoU+N3X#\?$;3@)1c%ZVVE+V)DCB1IaP-oN8ij$
&![kLgW899"\VtEZW9.^q12clfmX!l/i9f\G]3/O>k'(C(l8<$pX4$]/#i-KV2UB%cN@guMgDr0s
b!q5'kpp@6rK)S:&ZZ666^a3"2Hd"lh?,ucB4I=^4k-,MKgNS4\BX>,(!PG6DWPj):P+:mhn!jFN
]&s'r[4Qc6*U0XJ(:SF\YbLY#YW(&BO.=F=[]]gnLfU[a]cMr5nPbKEUXPIip^@I#Ybs7a867fYq
JN*?5Knq,[+p9H]I*+a<c5Y)W8=_&rJq7ZH.FOH+_U=q,V/(@@dpda,_&L+mu(MV%i=C8k)-utX#
HM8IP;\jU$$:nDp%"Oh<?Vkc2e$1gb.0;@NK;58`'IIS-'4b)6onV]^P6W/.V92$s[9bTPF("LQ'
Wtn!p=K<lYep-Ynt("-Q[_ZAg0t,:0A@,o-s"-+k2dj96GQdtb!Rkp0#r5j?Zhh0/\?-qGS/'$n!
F=/F8sU.;<X8f9?JUC)OmF;m6Aa^k[_83<7\Hh2m$ck\k;J1I*>p;+pAZ'lV%1UfgRn]h`22Ft;Y
MWq4tT/[&4HhW23abk&JNcfmNX!?NV2eXf5nV_+ujrNj)1:=Km(hg&O!,aG:m*%t)Ii&XK+D^elY
=2q>.B:0&#lal7!&(sSC>PGr9RN7^*Lj3N0S-7U]Nr&!lhk?2K[p.YG&&f^.Z_:FT6[J?"3:&7X>
Fe^/'l3Xe^9&87CR^r-t@S:I$qIBQ%$q%`?_&obXf?<0FX\SdD]LE1;PTHS>k"rqDr;F]H'2>C=E
T/GZ<^<W#:re\q0$hbiKN2(Q0^Jm[kE7*YB[RoXNl!HNq)j0:9E3gms`.?os[W^MGB^H?"`(^[m"
c7&\ge@LC:oNcs74IIN1g[-_B-hILNkUVYsq%s2SW/oUWM]'aL*_O;R3gGsCh5!FQ^?^E)cCnoig
rfcjoP]]`n:kF6IBQ;qHoq3s:WsI):Q+7#('rQP/DXpa="P;NIK_h$:`Y5<oG8`VcKJ>Hna^'5<T
>UXirgFpnFG.8;k;_elC+8Urn7*m%U:&6*ZW:=8nNW65XSR#,3ea;EEFB/ErrDh&Y9u=j-ZtPeji
Vdi-.B.tdGrd[7$AL40"-&C%[8=EOWAd`nn^T;oeu[H!#T%?<kE'pfg-f2"t-W/c_7"jA%k\XT6J
^=K;[;$$uWeY0oroXa:1a=mYlL+YD3*N]s]chlNo"k*pX:h1ch,qIE\CFP)40oO.8_9Z>YirNqKu
GLlE5]UYX`@QY>%;+ON:9,9,c'<`>n<W^u?q%&GO-*XpZM0=Hf%m=EbV4hL'YK?c%,(ds+II:Qt>
P20'F\T:Mm\er&ZE"&Jk0ohOZ%s/XX_ObL)icaSVI%[>f<_*Y"F]]AS_ndto^`(4O3Wsea*csLC1
S-&d0<OOH.<fM/H4%G38dQ+XF6<s?&q<kh+OqBM*j(usKh=O=M_:gII9401"$`+f\'E;:bi?,i)V
&nk?S$T](s.jJWe6h7aV0kr0PLu#n7I>b3VIRSafbnp8,=mu5O##rlZG6o%bA6nDQh&QE-^pr.Pu
OM1D@4W)a'CV&r6=!ph-q8plE`X7StW1c>K'M!,M7urgdF'YBeb]i2fGGf@Ao?d>YWIXn.T&G?eI
A09[eZcm?&HiU&#lm1atDB1A)Eg!4:ej<M^HLp,d)FF'i`[N/,poZN$YSJ1>g*ZYl+9;g;;7s2q?
ktaU82muZQ^H0q'!-]\aFhD1L;MnZi=%2UdHR7_Gb4VOYmf*9oac#NG-XMO_S0LLtZrm7bW/#Ajm
tN*`RVKAgSN`=[3g5e9o@dE"0o'63CYO`m\kJ;oJrHN%E;\%Z;N,O"guA<-mooP(KuQI,?L'Q!%L
,;lRH^9O;PP)n;oNNr[E8duKqrhSo&/uKB2nDHNek^P_gGKs&b>r$!;g:9.JHn`'Be4JW^a2%k,T
Do?VAME(P6"m];kG[%[oi;X39ZI!B/U[>1&$Ke:\45ZZsj]oYOnSF7-7Yh^L38B3[[Mrr<kb,Lea
p5YEZ%rc<1r&KenZH@7P-rpFounNW$2=6!gSDS2lO/m&ePoY2!cl@)CRR+DXeRJ_!_iR!YEhQ5\b
2qf^+,$dt:\QS(D.k//&FQ:k=NNnL.URlZi>:?u@#Ag?&o]:nJO<s]5N2QW29)%(hp0bod74L3fp
u^72#5I@B(k*Mo]4BNOEdpeTmlu!tF"+!\[hc0pnWZ*LHg6S*c>L[Fgg8fF$=%.=8KV]l3Hs)#`u
D*uWk1mS.i7mh>Q@E^R'`+Q`eB9Z5(,DoEElkeKT+`.HoMYc,i<%Ae8n@Lb#]s+DY_EJ8DJ$1H#O
1Co\`f]/jE';S`"(cK^i0$/#9)4o&d2!Hs5Gf_%P2/@-rS8f=T/[UT:^&lEC'"^Mm*i:p8Ube<Yu
?i\-Et4_>D/irZOF.+u_EilEX=h^OK>gs^TBK$$aJ^i-$NCZOY#5pN`$q8DO@>P`p2$S1,_lR:7>
40Jl+?n@9Wpb?8'Fin]+C/nQYHt3[r)XT/cil[(M%G'DR@T.;`n/)HAl$n'jNi1o`rEbW6Pj>.YU
.-;rS[E>oFni"VeN>)$i%,C]?T:J3*L2OC>cWfH1?5NU"'@,shA_Q7<:Pi/Xd++JKN!<Kq:5Qo><
bpY_]G^lIZC5`n]iUu/cIPLIl2e1,`A*;^J!_mK'VIR1hu>bJGY_C@Io;u\>&N*XX!`5PESJDm]`
T"V5hLGFf,GIL6#KG-LcY)ELO#]330(90c,f[]>/*Q/UbgXrCai_&A+dXAab+b"5U.>L,H\)`//=
XZ)^i*m=AAR9c:3GcL'2!WiM+!%le']'Ip9nH8%gp=q4LapoF.FQR"9.ms(#Voos3dO67:d7l\p\
X4`D'-`dL?nc!NT]bulnXN=fK!1\ocQK`BcZZsATiilb^]"f98?>WVpPBt2e2R?FlI:;'KirGXXc
MV.<^C-OU3+AIt7^+c@;u:0N=8_4=kF<?i!pdi;3*Q'/_jGA^YTL>f,'.[mlhi",Slti[QWL_?,$
L&Rc>hT-/UB,,"oIHbU[]?X_ag_t$*rgjB<0-s\n=fq7WFdo1:kGfl#,<`AOt=G6cJ*d"3&aL'a&
UWk.2IRJDRj4auAdB,,35=<[kE#<7?A."MU4-O1_s\3hM2logs<&?^i-=7Roo9P8Y$jp@l=Ja?$o
9i8k@2F?6Yk0BsL5nWA;K[X\":=3l*gAAV*R^@^P,.Y[VIh1ghIrU@ZI,mWJU,hnReG&pfuA/aFL
Vk_dpmu2RCd;H.8A3hG>h@Q^#Ib7)95jadtYQ"S="FF)>LYilfrN#2UJkV.r/_1Y-aoLiXGk[.[0
RLBZhU1@.oH[D9X)a?hnoo;-3!K]gfYg.$LSlI`<CIroiP<L3q+XDhj#+TOX^5q)JG$M%OMRHM4@
WPY[S@:SW]K<Fd[32rOk9K#;0O4emb]u>-Ams%NV;mn@kOV!?SrITR&h;B2lp]4j!KssC;A8DI.j
JNS1D1/,?q7c]3?piL[aJV<Ok.DGSH30MQ.85h54?rUj'fBGVVq,VMrq#+6'nu0Tp[/iDtp+l8q`
+2s&;NE?akkDPuGX):O7WU"j^`nh]UW,t\SMifaGK$>mB)(qj7UFR4s;i-aMtNH[cPS(B88:/K^;
S:"`Zl1IrX[_8R/O?-Z5@8Xe21R0?/=M0[!\guKR(UgihXGEAKX2L76;NUl[27dr)a8/+jAQI@l]
Opg(6T;rn.ggNEUq:),9#DNs8#KP>OW[,mGKKH@N+Fu].H/,\#2)TmS7;.(]2%8HS.A@^A1SbZH\
`RirT+5N!P'0**]%p#Z]*^.q!j:!8huIW0?t81Z7Ep[n<cn4Yr>BX(sVLF:B(<2EKQc%r@FbB'](
c>SKqa,&).,Q7=uGp9E+8_B55R;O,!^?8@!SD98h;7X7j>HVb'l,HsTF5;u))=nmhF>9Y\k#G+S!
tZAT`(j1/ng_ZM)dONCg?-R)*mOJfiRM(;E3ej3G2*O19l2.>6q&FI"=]k1j153?.QH/:4:`**l*
"p;oP^I<P,4iJ5O2h"FkrZECRKKb_%!)DG/-Rj74KWfc2<Tp_iZb$(Rb]F`=3uFkJ,eAq,KW)rha
[p@,/U>.$d[+f&/f6m>%p0;^D^d-b8Se>+Cftm1-dr^bGedc"OkUcH(nSt-]bgk+ZL6L77j)05nF
^E@V&F$D0HKFK;*L+)@`h^[GX(c;-*9i6_(!^=0BR:U#47Sg;.f,4cL*!u\&J[e3cZ*FZefB<ar@
8Rap<jLc+D!na\E-ShY5;FG(kO^bQLqLkV<BR7-u6Q4E6"^;rUt=hn5V2?,u3=f((S?Ac2G@J')7
3l,N%2<K'RS\m/%Tltd6JI,K=/)j;&A1BGm#OlVfa%E**0]jsd7b)`(2Q0u45NuGg1#LuMK_rEZg
M\*Jjo8<\!lIRVt'=M2un5]RVe?$9Ao7lS/93^N"8R:3p0quNSc2I<cd8l)r<,eE_/7g(*m",0:o
\bb2YM=)fg(dBf!f1"r?8E,BI>5;`*oF!6FhQcq=p)bH&l%p](rXT]4?dd_/ZVn2N+V;Ai8X6\S1
U!4B'X%hA:D&0/,RGhT^FJr=:(k.T.U<-(/`ZWrrA<pO4C4Il^_X>rB(+ojW[?r"eE)sKjrJR$@$
P9GjQ4'*+CRD=:UN9;.7qVg:ogF!@dA6QhRROl<:Z><R(eSl"`e\EJqFY\`X'<`MAE)(Q8*/P-h-
n.RG<#;;8^q5\QBmCBJbPd`A@/3$,u7*N8r,Str_&%P4X3]Y@eB2+CrIL#fmL*.h!s2_%%09G-<%
CYf8S#urC7K8_2kjbeFtMHA<sh?bY<J7?L+7BcBR\fnS@5():2(<nKlObkO:J!UG9-'bkTQop.Ek
3>p@BHM%Uh(H?J18q)S)DI<u,q6h<.+Y@\i(Jn%ne_[ZqtsqoWS15bnX?fdm&'u]*]q-#S-Su5]L
4ncKq*M%@hJ:g7_G2,']4+q%L&,_a8in[@P#^9r+)WJoVmFt7G"_&#FuIuR$%CKL#SHWist0UdJ[
HN.%DmOFXjIlm0hbWC0,UcYQWeJ1X.mN0p*V1fnZOU8\$al^2Odd^-#G2ho9E7N,>iO*j@QV^ZRI
heZW,.<rV2hN=+R%!8dOr(%15l2/Zla(_l5Sf&Rp#fqE<fB$.,`7i0H-E/G`R$a\Cr_"1gi$MN9\
mus#8XE<mAe!";GAI?0V"LhJ(]ZIk<#)n1Pr]Ok(Mqd^GD@/=7/=(Wen[o;Rm7pm)#d4!@Fn`keo
7[na>SsNN1\1P?5c4U%RCD:tTI=_\&%2PfF,-8Ec\sfc+8rD[SU/.C_U`VljtCR2"Q.8eNDJHV\?
WEF>VM4>B5D'*m=6*j=q0Lndo%pZpjTmBr4,u!*^>8S7#LoNrHZM^3UOEs&"Fat[EbA1qKLqrdgj
aqj^nF\/suD+@S>'S2D/&5#5i*#92G/NQV79<1K%LMJ]$q*e,KD;B.>t=q)HB]niZF(P..?7mWAI
/dNd5"RK!9DQ\+N^:]3AMG<VE:@c[6@lL?)::\;VlFEE!7f:O2sQ-X^Q*)1O:eq;7!^B6HsV&QQ$
am=>dM1;DI:n2+0h,![H`5#^.[EiG>,om::LFL\5RjXula0OcNL\b@\fYW9e(pa/A:?oX,++l_;b
50K<J#>$Kgh:C\O(0EPQf._V5X2sgIM_MpqN#jBT$3nW[,Ca67U7uJ%cMfse(*K.dWOWYd:det19
]U`q/MaK7uB+Zl1C=?E-o:[SqIEf];A#?en4^q+$Uk>rhQ='U`p&5hB&TkB3q@eFR$QjHh%"eFRC
_0_n;=B*D!Rd=]mnfjEQ1H*E]\QF]-ETrr<AXM\rA"cYrfTo6P3*#g'*t0156R:*!Zo7/4g-Olu]
_nGO4nr5ciQ7Wn_2q#664e!+<7HnjY5dce@VlqZ0hFEBU#dZaLKDr>TXrB:Hb\e'ib36fqM/.,:]
\#^?\D[24a-#PILk@S@E/'i\MdeqUU:,;Xk+7Y"qprW5["b$",qY`Ao%^/sC[\3(QC*a:g.kY`ti
H6Vil&gkt7S7sPT_oTuETm%,8k$dU4<C;R;W^>qX2]Q@iR/_fN-iFQAO!V142Ssd1,4;6qNeKrP^
5,mb_0>[7i^p"WSaep!c_[TSo]e!^CgI-pg0kBDakq#imMHMi]kdgUe@W8A`/K<f6!E*'6m)R@N+
Z[opVIQ-;+(d^Yo+l&^m&[2S%D_k(D>7\F&ujI@"%)o-)j2C#mQ(`-rWSH=VoE9>'LRl)EtDekt,
X#6">5(4rC#lLWm4!:`:-R:eu8QD7SbJ"^s:0k&B,fmVk+^2T/]FX;N_Td5Lnh`:>Vrr>Lc/5NIG
)4*hA2d?:^:]CEmaZCg"l0Ie_rr>T9f?=nE\DQ0te?29RL4(TP!2VmiQOh&.pPpO!!;*AQE]W9p%
Yl.jih!RLP"%l*=8'QY+tRjYH0t:>oZ,ue5:nkn<E[Anrr>LcG>/RF9=QihGFJ5ll=Fru9t+3-h^
R%>`;Mt@IbAuG!2k9&)Z2_HmD_n.JG_\LW:[!jCs0'*K_gJd#>R/u3NSipSH/Q0p<#<8Ug!8-J;(
oZ!do'?/!H\((Q0Tq[/=J\7G/j_Y^TcYc7`a;!UbB8J0ZB9F!]"6qeHI.>ID6Jg("&^`Y'rdnb-l
T8KBp3=M*r@f,S0TDK6(gk+&,BM!D0R\+[p$F:[_J8k3jKTuCCspMn6Hd@)h#5^\mWb"Sr6X2/%N
%G(47c/<660&Ckl$C3H)S8jH:%>X^PI[Dee`'rU&$JjEs(Al(4eTR(:o7d$PjZI`f"fVmT<$#VRh
iCkr5?/4(Y-9]f?N!n/Ea$6Q^qX>0A-*R*5+5!P9<5,V-?Tr%9D?FeddGl_E9i]3SQ$V8qm-Q^et
k#3fibI(Y_j0$a.*a+Tlm0ZQk[:",$JmLYrhC#-1cS)I<8^p26Vs-.p@@HE#G!f;MWjCo9g\s@7Q
Kk`T1UbDAAEA)r*d^]hY(bl'8p;fR<a>68N/)m.$:'g(mh;V[BTe6$tWs]0\o^iiisC/6ZCci/`$
S\ia,S8jXh#SDX]s.O;4<N++J0?J7:/o7D#ZT%Z-DUGk`Sp9O(dk6T:l`+@bn<qb?T=I?@9<keZg
.>`PZPt^i3(6fKMk5l#;?_#*c7U5PuiD67MKrZs;6U\Q3e'8?gPHeG40aNfs*jE_M@ZhDHY&G.[,
@<Jrc7'Q5ZPmgeA"`D&3uEUH:,%Wh=O__J/hqdn97EN.anFSu<`Z@QL9ML`VMr`uVe,IX(m-`r4V
'b:2&)gn2I7,rg/0$RegZ$`GQ$VaK,TjaZcee#L"J%1GecX?<l+Cr"8ZjV>id#pY#q4GS5.^hoJb
fJ&-ZQSDOTa5W14CU.2Yp+?LUQ<Fmkh($Z;%$_Y06rFcR)n/c=Pn:P6XE9Mj.1]Q$<&pr*V0f^#%
QOQ0XP!g^hjY4$aK]=V^T>5HM$0;b;3Z!1D/bQ<r#X2?.Y:GmEkHIJ.7:3&W.omo$u#5.`R&S9Wr
L#kBu\"F4+F*=_aJ"/!9<Lq$]'jK]?C3IVq<q0L&mVm^N=_h+fBK=<,FbV3Wj*2NN.F:1h",RuVn
Y1-UmUZ:#2<)%Ai!BkHo7H6,!&WeIg</Lue*V:76_[Xo!,\F@nU\%0p;_\`21G-h#Y/J&]ELH*gG
ZUogDWJ<<s5HYkkBY!!5JG>hiHLlhPbD+S%p3fmkIG/KjpZ-o7kGgrU:2"Xo)V&DuTf2>G\592#d
R$"X#=n=g/m[1M3;_1bj!FOAOL3nUrR8$DR$Ch/`%KOB:n9QYg^&hJTMF<7VAhfMQdl`r)e`']]k
&.=V%bnZ;#2i<^7SgW$Qqn=>3h.=;/M`tD;4HVF\,G>I^Il:;O[Ea<H4&&V;`N+![.5Na6Nb:V@q
fXot-lqd,T+LBlDRDm_@KA4l>b!$d7K,i8k=15B/dr"%?k>]2*.49X)XE3lNjd"]?#!!`e'^fLYa
8Z.7l:e^7o;qR/dYDA%mONQXN#9Jca50Nt=-jH5oi:Ij=2!]+_aZ^omu6nY5GiF9n!3*(nBtNtY;
l6pooD:""HckJUW.DnUXG[A'RcCseED?&-ZHQY0m[dkNpqJZF`AgOjV_$IC\+)cKs:mC`_N!P<G2
IU4niB-aSX(mB_S]K`_KlqS<d/nr]/BQ82me2>iX+X/u53r7.Kf2E,p'3;e'AVdgh[p^,G)jVZZ1
KQX"Nha/1dc2b$G.SspPiboVkdo;esoSt)&Vh_@bqj!KHD_?G@1=P]:TYD^8N<$m(h$7/6@2h6r^
V^`:]0*B[[+?GZYrr@&Y=\A:LO6`H+QN:[GLs*))F2lFUZ&EQ+rrASTT2$'fA39?Dbm;<3^K>')j
CJ="/tK_oN(-$eP=o6]F0U[:2n$U%mhU<qaZ_lJ6M)=g%`lOo'.c)9VpU>BM$pIHel#B0:5Z5tk9
c4EnD.2gT!e4Cd<-YWkF?@fG;`gJ:_ZH[*q6t-%;T9tj@$a5ldEAlQMe7-1:qarl-e$$rU=u3;/J
GW&=6(;>2QD8,8hnDCFYDX[==^jIZq2So=W>K'C&)>M*?(%BZaCS#m^-mrpk8VhkB/UIMhC0HDM-
6XHo7F5Lt>63a\.*_fV2CCV%%%&<u3k.3708;uZjU575!d=lS4^J"lQJ]0PA91K4:E4u%J&I&B?.
/uMP1&S/uk4D39dn<%/5U-u,q_:&>`9ohmI'\IT)G%Don5-.p>4L\HUr;F2P`+SjOVb/H@;N9nC1
O`D'.!ta7iZ2,h-#^Z^i?%cQUJcIKjuo%PA(*]&3mB(tQD1@uH?&19"iSIKl'5/-gF'7e.SjD>\m
_,l=/m?>*Z`gX]gre$9J3B0O1'-&Ih)2rh&t>RAX@u26h)oW^R#l0R_88R5DeUsG#"hQ90'T`!/-
LkO%g((i-u6TKVo*Ak-e70W-!B\`F@FV]VN:DOSYeCH+bV#OU=0qe'<'+8X8#T+1hIISM>'Z?Cgf
4G1G'+57\VS^+=0-&3gaT4e7'%!*=n<:W12-at(Y4]ZmIcrB,F^V@2uHC&T#KD<[=7*1oa%e0L)I
`P5PHr\LoQ6`=peAi!DAi&]0D+HGn,WW&D['e"9AK3@$Xl?B0t.s/<"35U):<-$u6_WlWIoQl[t#
/-GHec,Wuhh<X*0=d]Rq/E`g&".2G:IaIG!.1SG76,P;WPs74^\%O+iM6S5#5P?ZrrD]X8Y;@#Q#
,5o51G?WEKjr]bmnR/r4g.bqW/"YgQm?<bOTu"_b^:49E$,U\924dfVTnn!2b7q"f%t<4?Vf_o]Q
!0>Y4S`f("7T^BFr;raj9uZt97;?RVMNq>C6mcp-%'^VfC*D^RdAJ&h3?Er7K2iA[C(6`p,9]Y9u
E4p9_!^B@@)Q?Q[rYP+liIV8Z\j$+7l\poJ"`:b&trYqmRSrCqbSEK/mrr>f`o_-;.N)$;I:JFKA
qo_4iIgLL8jnWP*_&FoBR\`&mY(!$nV#D7(CqgM0ieoT_r?q*tD13S'!7J,-rrBDh+..4UcYgDTp
*0%/kNQp]8)^RVP\*Qgg)2NPqBt3te!,@GH8f$AEW4GVT!-Gq4g(uW0/;,spIS`l!2f3KgT1>SVg
eunmY>69)n^r5rr>^k)uYoaCQtB%^XpEsYE),9IG6k*0hNIFO&`jhfmi7HJJb]j7#?4KR/[08Itj
"_Ep@cq[)`%<T!e0^KDYsUn^mS")+pK8KD3]An:1:<#=Op<HD<=q1&h5kHJL3?;EZ\Ukc[_kLdEt
DZhuB-G=eC6o^eTc^[#f4:DiZ6rrDnFV:4Z`-RcI!3I7=oG%=`B8fS//`4lE+Pf8e2\R-F=gX?!^
H>Ir+e6DZlrr>hJb500#Faqg=!#%eC]=(X[-?4JfnC$PoI5X1MOO(La5>#sBF,R0q55@#ObpKc<r
g2q0`]c5A(s&t[=6DF:P%34LmW"imj+#I&1Q*Y78m^f'4Lb8:nV+`C3<mO\FIlC:Q8k:$KQ"]>^c
Doe/p?me!7H?5e=B+)3]S^LX8\+RaIcEuR^(m+Znf;HYE5*1Qm0[.rrC7!r+Du^oNBtMnIUm+;0/
=YJEoKrhWFQ--^6Bs,Sn(bH1d%0/cPet??)KG8$tK4,0J')]sX+*!o'dgdPnq"*RArof,S;)QYqe
%ec;H%([/NZ07uaCN5_b;i[_(m'pT5mSp]I-;S[L1OINQG:)YZIGDu[Md;RnJ=(\F\E#XjZc"e1%
`Rha.,B;b'i">?SO5._W9gg(8!rI=+?2jn,jQ=POh@s`u5R+8%Bct)j+$%m#r$UJbj8T*k8!9XKV
uHaUr)qc]Kq,6qQR9U#c`)*QcOG0;%Y%&WnXeXs^\%!SHA?uUH=^H%SET3eT_*PS0!(I6_&"nWl/
Ot8lWs`Ro`"nK=KEiJ;WjM35([/+r@X?&*B(R6N)T4#lLbOB&i[IM_q!]"0-AaNd/?:r=mbm3dg/
b6NBD,(5P1+[^?#%t^@=@%ac(P4qL5<P`ic2\1nH>+e*;FYF^kZa+!Kr1jo1\CR)"rtD*Rga=ad#
UECp6P*.6c%7p#DcXSQ44]&?.%R4b>%0a>C_TRXI%pb7MEJ+8(9\e?s/?(lL/NQ(=d4ugS"qK@5(
K#B:c,Y=IfqnfTe2qoo;+c8nBmbcch5F1T7nNs]4O%o[#f_7Mb1o3nOnV4Z1)?Y3F_0>P-*uIIDP
=pI2":pIVQf7cPkS#8ueFOWZc>")f^ZPq,l9[JI)ll!BFSZ%jdsePLdpX"fNtm=!qM/oX?h9I-]3
IElVles0o2lM;p&>$;AQ8)6nQ<%*3;!htDo#NW8QNLi"h.MsHp,7\(qf:]H8#Q,DE(fteh(Z^d+9
@U!U%H9lFk^aJPgJmp:LFtNu8WP]D$P1759G0KH(5LDsa#!:,Z_-"`&AGh_"j.7l^L[.I^j>i<$k
\c:qP3N4_TsfdD3Yg"L2d2gF9hdG=DR0uFO=FZP96e)"%&O1%Goa8=0%h#j((RVVW:Xs)9l+%rO)
+QYXkL5RrY"lJ&ON2R*ZF%,CG8AS%:X#$f/f8AiND_<;]4l<<$#HX`ZP$#W0Y%<.Kplm?'AC\Ra"
??Ick5c6Se)MOE'7bllLW2IQ0Sm2FCmq:VfBLnWPLDA^3RZ%a!jlq+ad>I(Hl1<DlL^kWB!oVq8*
MWMjan-&;3:V)%\HP.5>7C@VHjN(G:SVEnn$e,,lo5prNgUa2a,u:N:f5cG'+6V->,gJ!2PT-3PQ
fGTag!SYLJOEP\'B1DNt/$0\:)Hgi=$5/LPG@qG?OH743j>dj.GAb/)JF^QU"o.0r?qPKi+fb2XL
)\Yg29+$MCP,Po..j"dI':/o!:Z$HQ5,WTqDVC`HCX.%(a9G@_:?qolX-G`TGZ[PkOb7B@h%ZaL[
OUD*gb;7Q2S<5br-CEsW+LPqTL[r&`$Oj]EI:Bg!2t;eCIP;WoMKiBp\;C!+DG&p.(W2dOf8OeTq
XK_5l3T0iFC)A70qP#caJ]"mEe`H=D_9E%r\9Pk&gQM>ZhRN^=d#+"@n_1=kOhkN;3M[%IF[XH"7
SJAeFd',UX>#)!=nN+HXVt"VY)N9Ref_+lT@<e7uP[QU<UF&nOA-Or#lau6mMW>p+aO-r3fsumf^
4mXj_JthKfFc@c,V!ZNRie;PP)qrZ#V4k^F^K7m^,LPEu-gU"o^OC!MP6([7cs3Qh>#HP/&H/+8H
PbL2e9=NWMB(RI3KU2mSOCFR]r_md2@VPlJeG^oW`@TeUcZ&@AjBR&,GdHCUOQl0J!$9Dbu!;`o?
0!e>q!9=O3ZsMB#7UhsoA,(;iX%tihcBhJ%Uug!`hOj36lR*Sse*&IsP9\P0kgMpJGD:NPXm/<(C
VWPHA@-EHiAj#p6mS6b*\h.^H`CFMX^fPWrrBJX7SiWem(`N+8"T(m7FT@(lsnuL`f*I?e&TaSbH
9AqH?&=7hJ%+YZ@9NXA3L5B]&OWepMp)m,7/h-/'l^X3S'F(T9RB$@!]Luho`FJq9(B:j]P*O#c)
+e1mDKo6^hAPk3C*4?Z0q[k+%K2XjnVFh,&N>NA8u6l#bS"p>j;<8EkZ.*cmcu::-kP*j9GX23/R
D"Nbioi#Zr(<_$e>;*kI?jj-_PdpJuX*Edh:mHp$bFd.4$@X!!o;BBk<e-M<FK"ghul2>Y[UNhA_
/=rBX<p*9;i=-5c:\0DV8^ip`kaV/%O&^%`9P`CR;<^OtlDKlJS:;$lX@cc*V#LF/+7CghM6RT8R
j5D?"n6fTg%`"!oOeHJfK(gs&KqmtQVCL]Ek'%KM=s0Z3okRmkk+_&@m9<taM7!<2l]ET<!IZH6I
?uu]>\1WB"L@l=YY'8&&Ima:fn)"]L!Vam#V%YYMVJbag/H@p`n.AOd"ENXC\9Aq#2$l1?)MAe`#
fm&=>&-i+lKZSN(G<h5\TCUFt3d]?0A??;LMFUXYfiVeQ!CPM"tX;3WBae""?pD`QS'kt$cd99[3
;6`,u:p4\!@T7M60rr<\MqX8e*SPf;WJ_T[)iR!8S@3L['NaM@O8g&-Qm'ROd\oS1,QPZ\Dkkuut
GFP]uS0CF:IR>='JCUoLbk%iIUg9dWh\D2j5?!Uq]iLWm]OS"kmJ<Q)\/b/.3TeWl$1udab/\CAo
]64M\%94TFa>b_pK!`uPWE>,H:19TS.^Y!;mi]NnNA,)n*M/@k60tqlCGkg/,-+Nr`;]tDL[#a2c
HIa*Q)*n_f1gm*gpZe*uI&08G(mTOWWI35WU?2h<i2HFN8AgJ%e^ZpqE40r55(Qj1S>fmu4#KIWJ
=\>rZ23gJuNdWu+;b&=:[aZ(`SAHT.)PFND%ll[Ehk)!@(nUe]hu[kH0n<-3^8]O&sQ4s-CWf?u=
MU^,)_8]^(]7(%hLk;iD=!/G"3N-e-#YJ,=J("$&ZP3PH1[+&H-=.9`;^!C[!k0H%JrY)'VlH/&Y
]_0@OBEbNp8,ee$2p[IpjW7e/!/p(>2ps)J00+@9k^=2qBp55Mc[=uj15W2UpVAJH*i6rAY&>u$S
I_$Qq@j625(5G9cdr*re"(M1\oW`IWG=]'5D/4@W?O!2-N8"_p_\,@6/=17c'bR2!+.#EO+us"!2
P'%&LI5s^\F[$5O#qAOWm#DRd.<[F78/a89Fj6T`m?%ANqf!?+n9:)6`_4lQ6S_rrCOJ!/j@^VFF
ALT>l*S)ufn2G(\A8001oTJH5V/b;jdU1!u%\W(E/558c]7b*a>1pi#d8_:^^ZOaj1LrrDgI__Wt
8Of-;#!1`fDdYrOI'u"^D80r,4BZArNnpgQBiHi)-8#HpnhG^3KNVH;db.9fkN3/pD?TI@C?/3P]
?t;ma9kJ=*KRXN\l#q_)U5R!R8&oE8LTMr2Wn3h_*PGQCV<UZNiG5gEp<bqW?C4sGPj%0cSFYsNZ
99g"?Wl2.N+4:PqX@mJ:*`st?f(AC3*AuCl#aLGaF@'Nk^FVT^"m"QGNQ6d+k)EBX$!nX7;105hD
QV'r's<;B4ehpo9]HHB5FHXB<0T$0o]K9^iYV&m>Om3Q;Qi<57;ES7R\<W6_cV/&5QVR]`L:Oo2r
\$UcNIY3[7i(F'fOAJ($c`/A1K((pEd[iVRioGkG;'7Ws6&jp+DSA^fJq(/aurSN1g;@Cj\[/m"I
krr>@_2"1FfT_qGBHJdpL3DNqFemhsKnXeq'@FU6F&l.fpk;^mq.AI)f\dIQnS"=fr;q9i&r?&3X
O%DCirr=Z.5LMb4gM]h!RC\_FoNaYNV67fCLH#/jid:ZbVUDU:Qen==o?mR<8*eGs%`7sMMuNe::
g$JRH*97#YinhpCo%E8R(ICp+2W1imn$Te0D\TljJDj@3[i^,hK0k'E?[W;KmA^hI2>]KI2Ea>pS
RP[UXJE!iuX`iTe=eA^%e!;*S15j7dJ&[cq^<[FMDYQl@'hSdSbK2`dW'&H.MmP8lC_s>^1Z7dF%
3p^?A$Rm*#Bn1-[ncraDt_;85(OO]6$J_n>pXIt?gF0;fDa:f_W9-h(p!Z_[df7j%'n?Wd(=1+kZ
G1h-1&/;g&HIK'<G2b$<dpjU>]/S*2Eo\i3LEI1GrDI<&7-`h\H@3@F/jZ(5B;<L.+nJr^hbG%M4
d/O*ac9AH.J)R1Frr?uDok:GAk^'s807htm`QHMZ5*PtunAO&F7t2R7-iTJ-7f0bmrrCL0]qTheF
FAT-Q">9V%XXaRoF@,Onp_4Crfk(Vl[PZ.-`"GOHD0Y"4ecH33T0_N/ZiqC0>RQ6(FQ%(,q)MM@%
K8WBaGF-A"QjihClkUh`nOO`kG6WZXJc]./FFI.>>mbrrCO&Z-N-_Nc(NV%=$0%gQ;=rZk\Yt;sC
[U^.bn5:&XXeLO34]1dVFU\DI0qLrG?2)u/8SjN*PJk6VKEV;n)3W^l$haHkqI*fJI3`c#Ylf>_r
TUS(rOo7gO+bg+TmN1rd0C0KMrN5#PFA(kq6XR,kKZlQ$Jo_N2rl*jj<d^W`\;2kV3H+l#npfObs
8`''I;3:7S4aG0])W7b<B_+G5Bc%86c856?X!u):CHH1XZJt$F>BTYZN^*Ab7L9(5LZ%,DU<:R/N
_Jl6"aS'PI]k`YkKj)hM_<]=UV*)k%]-QWYE^oolkSOXhC;<7)Y'/-nYQ5W,"pRdLH<"rK45IR]p
[AmaGAN1&(=R!P5/*Q4ZF34SN[Gm:VT`1*W5(o@5#eF08t@H:0+iJbELU6b796La?9:T?_dC!ao,
7Ha"JKco!qM[*p\(2)C>1f+4<sJ`FYe^a/%)jGk%#T,4G=$Mq!k&l>!V<dBlk(j>&AeW*QY16!6c
Q9,NGrak=(<J#nN6rhd.DO/&4Om]A2oS#7phQEI=^M#C5K?*Sfh$f;6hV*(U-N=NEm5`8t7!R^.$
G>TH+L2_s,qb`"fY?69B42g`Ke*>/<FJBc6N42D"`eRk?*PN;g^molj(19FiDi->67>q0G@PN-[k
:QhI\"#8=rbP/++frqj7g&OHWlXRgnW<\J]ru6Lf%tJql?"9,Eo+O`c.(rH'(503"n*_:?D$4<\(
\WjnF;J:,H<,-.O2$n^k)LcE%cE?jZ6ZqI\:493HYp`&bq#0V]@k&3I[T!X;&OZJ%f5CM%b!A]45
ZhK_sgMdQ61+<I<a<3n$jg&MO*[H8+TIOP*30a#Y6KUAjt)8R,%`;VE_tb&jCBRK_M`F0+W@ri"G
($dhdh'+1@-W"Re0oopU-C-36D=EMaBNU1]c!1$1D2eH:j7!bCP^*A+E3diL_[/U,Oe(Fs&U,q<M
4ekSt"!gj;PZPATh\Zik[PTVGSM>d9;)Rg$Icr2H7j_2UF[$9r<=UlhNeBJ\ro]fT=d^au?Y@AEn
!@(Oh(NKB`!1W%&h+RSEd^]H%\sj-Morf4.g)Wr/1on!\6c$%"<WOVP-2@@N9*!thWr,pr8c_,:U
@?(CY`W<pI[*AEa;*W;*=H<opWV?J&m9Vkri<Xj2\(P,3Y)F0R$5^X^IJcW6BV0VS;kH-D'-g\$o
Zs-U*bEj&&_o\#Rl3hqgBAQDH!jMt<,I*gbjJ%t'4oq8S)6>7s2Hopo$BCdpIkZ>g_6BCX8_l5h%
52+g&N.G_=)GfdcZE<Y5[dF+))f$=O3S[4YWQ"^,!TbO/<n!Bd0T5$;eX@)WRXHXE])Ns@@ArJWS
LYLi$-PZ")jQ/"_P%n4'b;6Ra_>9dl*q%,==5WD8o-GHNIU'IIORDaR:JtWUh7nS&kLKJt.7Om%`
Oo\S=tH?ZC<ms+A^rh3FJMAF^#V4or)aQr8/o!'X<7(3H^Q6X;-k5Hac!5J&(D]8Em<i=W(!4nGf
Ari,LX&t\R1lglG,T_4g"_a3*3*.lh$!qHE`ILjaT$@"qi5c)<5i8F6_Xn8O+eGXlj%enDBPP]qV
r&-'VK7.'b3I,KBQ#]j@U;;*#;9:0I'tA(0\!+p20C8o<dcrrC(>FQ:u0(TZ3b%iAk!YNTFAh4.u
<];UUUAV+DZe<+mZ4D!J3:/R8sh6E1+5-dR;d<f7nEHd[Z:H=h+`k6P9`DjScd-@i1fIWe$GCfZV
:2u5*9,!6,e>4=QReXkO?QIc)FCEU]-k[K5Fcu#%mq(3G]5$AKIh6sFJ?Kfqr73N35Q:_bW/SJ",
:Q1lf>]Vp,2?!YE'pJPm;Jq+['(<V!`<MR[MG1hLaCQ8'(/J>aSpQp`9GEQ=*8_QPklCL*B;A`!P
FKQ(mEDtC8fb=)EK&:fM4/(5ckfOH![+JPXs(b*2@+k*@XX]d2=!>3QcG.2@C"bDI./&Qp$7S.M]
J<DSD=S=8o8iiUoN;I-6pFr'sUp:T(?KhD]O(gF+fqVp/h!=/::0nFDbP`+tb%HJMrFo_ZA/j*YL
)*h#FU)gpL7EN),".)bX%k6=))YL)PX%]HW]0,?DOCi.K1l,I^iehcg6h.(PQT8;\g*k\:(f65lf
%rNa(m/HKIV:=d&_Me/intl-M`,n&B:,BTh`c[p>Cgok>36T0/`bKJ;GP(`ZpX@pgiL_N%E*&"[r
rDV\!!4tErrBut6`L'a2$9^EGIGs9cGt17PtR]%'q#A-UW0jqGd$.W(A4JI7RF[F7L/IuMN:dY3R
>7$Gh\4T4c&bW@oQ[X;<l<s)c(bAoR@MH8R,TSk>t0S)9%(ZCnMg,*j,229conO:glfL,-6-j^,F
s7MN&TsR8AR2X4p47>sGC"][e+J.!bQ<lI>+8?8G,doZ9L[M`Kr6V:1K$l0S$)N(N\5FX8J7Hl5J
uFaWXq?A!BIA:n4]J;ffSa3H5JGJ2C:Y6Hg5<E:c3+Zdk@[,jN2T1?*6n*Shtf=\3"%mI($kp-JU
k81&DLIrutPBE(/E2Q@0f3(7'&H;aM-N='nI\5Sl(XsM%O44+/c5Mf/'f)Qs%FoN`/RBT1d>QI8U
3$0@LStlZO2eK!H4kjT@8Wf.XNN^Om4XgbcQC'?SV0TbrF3t>h(j#W]=U+#Ma>*;WAQpuRc-8JIU
!2.J.=Gk5p_F&O',&j^?,30-F#YHXr_L.WY,.f:,8]S^2a"?*h!M38md$e\gL:Qh+dcb)LXg>N*Z
?mlqHYcaH'k-!JdrnM]*,8N(,PdbI1\AeX)&fe`^eS=+B'<B*%!q%p7_E#`1R`:&V?V7SpF<kJf?
)I)pNt.P'@QoUe()PJBeqUit2QEB9%g]D01#UEWLbb5L>CVD2n^hn3Yu$VZ_t1/mIs-sFDhFIo:q
>okbL?BNtd1:fDhh*f\^eapdM<5m5eg!1;;kse'`Df,ZN)q53X@M=kL[")<Mja''j;Q0s%!0"RFL
Is<"#4_R7$efZ/3-F@5r*o;rVnP!X@o2>aV^V8d`u-M5i/,M?J;42;i:d0NP53;AE?22Jiq5O`^C
/Vrqg3_0V`;8[qJmf%cD7ReHl;Ri67!a]Wl9a(k:&Vmj#&ED(%Y^,rAiDe]4'6&m\?78^SJf2+>p
gbgW&#]LOn.$97Ldq,#2<epu<kAf:0KoV%<KFNjRuH1]"3GQfuq^^#^I[im?9nN=M0Z=OpbO3-*U
>8ci'9I%^U@``6q:ZKTAGRb*M@Q3V?W_2S3i!0#_nb!JY6P6C;/WE<]2d?`):GOHt=3b3($O_&5D
/*BAuO)eKaW5O_X]=aogM62i6-?LFakOpW=cpV.AkBR0PCL@dilR1T_I;kEdabG3m`OO'<rrD)\4
gKq!.'%fEjCijf82KYF51j?^VN@n@58+Ogm&G=p0os;oWBSs#"S-,mQ$M8UrJ\spP]N"7j?`^r!n
Cl>_t;c/0)T-4(T)>WC%,QqoHOY8NpE1W8+?#)gIbP6*^>6$`3iZJ]P<+p*g^_dCZF89:T52WlKp
ItS?``Y*B;Nn9SZEFU"d:"XNl_?/e'l^@V)GOo:0aX3fog&jPOkPDiB.+<cn)Q7#uTNF=Jd"8gND
Va3G`m<]tAs&W\hYO%rB!O3bgNVPbJo:e6:c0u"tXfqn5(m?^S'hsl[E>:D@++qa(8Lq5+'*ZZ*)
.D?4tiuRL)mKkU1*kBJ2'ngqNW5\$nO&7.@IlJ_(aPVLV7-dI/Znn8T`?jDI`s4URFV^S>4qb>bp
\#\*mb0H`4ZJ=Qqg$/4"T!eS1Qtr>3VBtF@PZQ13XFM:doH>)LJm((a0mA!3AN4sHC[%g,q*mZ.8
IdOAsYkZ`%+./nXY/#fc%Lo$M]`d3!LIZHBs#._BEJk=/YhBO&"@"fB%*^"pss.pr[\g3HWG+@rF
!DF*\?FQQF3>^9M7rPmQ2HHcl!?OT,<k+d9lRO%(P@Hfo9L/5:cEkuT9rGWk.KNHXUQ'I6'U<_KC
8X-`@BUOFHZ+MEu'=3in<5DX%MT)Se^0.#/5V+E06i"#TopK$bG-FJl*QX5*\k!Gkaa*CVROl7AT
PNcf,BHtXkj%H<+e@>0_'=Os]?<LST:8^uoKH&d<6`E,3o[e<Zo3!^FfcqNf?5XJX1d:^gC:nk`3
jEn'bWhR;(.u-`3`5/udT>=N,,O_(`T5i[H>cXY<9U3LS9g5BO8#0J>BJ0ZhE6,!:g!S,YnTbS`U
efg]4Rg\S#=a\De*2Z;I-W:O6krh:C'E4[k#Sn/%MEWbe$1/!:c5Ei];&_m4W>gW?05mIf;,q)iV
CbdBUu,nU1nC(8`@bXIi_q6;d"q7M1_(Wc!c1j3d,:pMt&e`a>Xuh-I%NkF9s#!)W[qW/oKW,)5=
QL-LWe2'/=_+3mZZ4Z5$dS,"dg-mP[R^MPta_u:io(sAep%U>g-3<@qQ"3kE7b$n63o]4(WTAPU'
3a\2$E]9b?aN(4SSPXX.iq)1\q%;l3\$sYS>]8U=:9W[9dtgKSrrDmh?DlfE.O!6u]5mK+>dQOTi
kc]s8'AQ8Y7<b*8+Ag"KcO\E<$1QZdJP@e*Nr:Ma+_gK^ju-PB8fuabn@MG/;Tp"Wbs%W`^FJ8M5
h$_k3GC8FVY1lIH2MnUCQP]#rgll.dE@"l'3]#PFph^-q55C^\F=.+l(7Za,G1FBEOF.aq9;_=&I
SES<;9R#6(@j]5YiUn%,[73>qJ#ZB6!G,O6,-g!Pa)e*`*dii,"$YOA!Z8oI<bdueSGEFAG^G1sE
+S7L%2RKV<-m^ZQ"eoO?\eV+H'%Pjj%<=8Q@iF7J<>-En6DDrg_G_B\m)h'KaTe>=c4gNu<JXU.U
pAAFC*D()37+WrMq/M(FM8*1qbJsq79bF!bX8-CMX\,O&rrD:n(8TR2SP`_Edj*G+=5Viin]"o36
t5e^*Xt/8''f[W7$O%-i5bitVs2BF4.Fd_T1O,)F7YSFnE$"s-/?L'4,WFuOm3SW7`[X*e])g/5i
TN1^%&>an!DPU/:RuSWp103[UW++"'2!W1SsDZU;m(H/BIKs/u'N3O]cZ])1EFeoa)6JqG&cApip
"@X501bnDgB=>b3__#gk0qnWYD3Ab;*(h;F>`$[`c1<KCeX;rGpMl5/XP=fBS;UYPeWphp;0][<l
:ahl?N?a^80R$*_&Hue?88$2SmVfbT\:'_!07B[eL?c7$u-.CX@=q<UK.8!HK+cN$*h.AE"V0Z!l
]AAd@D]>J%kcN;QG,@Bs(.l1Zh@!p#G.uVIeKl:#*Nf6"%_"b-aZ\D_(hI):oVOS018a_EQ?]:"F
?=fl"c;s-Q`.Sa-,i!iGWJc73<mCQ3'*4*%-(3uCG//HUXGK'ZlS2?*g"88A:30TRL6f&``/1^b+
BPr&sPk$d_E4g*Q[sTpq1<3Nm6K-bIR?4#Yq1ddV,M9`9KZfKgoT[FhFM8KgtM+bq>P<Upfl_j&_
C/M:*Mkjmq?nbH!''O%u\.nX"QAmB]i&"X2<GU&(S\bhJ+B3c$+u+d:n/>/Sj^9;M\q8$E@oirfk
,=/dWEWEld+c5^3f;H&77g=n4Jm>84Rrr>mn4#H3WTC\B5p.SiA<S_3c!$eeV1@g-;*[I%K\oUfj
n<tXYW:1^GP&,tB=WZV#8Zq4N_^(Htf!rt6C:6SH&FR6_$$AQ`k4>+6bDOjfZ':Ynm`U_NXA,*i9
''R-OjI-"5l+r*p;LYp.WnMfM^$>M":j`BYEM!CSM$G!o--d^nQI[I!'54Fo&So^T6uD?aLok\0S
B_K,nC$8`f"/0"3)c=%DK6*^:3#ZO<hDVAJn)^UQ:p:FJ<n>]-2YF15qS(nVD=7P@d7$Y>p*P=Fo
PaDMM3V`25HaNBX#U[GS0Hi0a._d?Hf\1(QuaV7X(RF3gh*IIgI(LKD&GrpCB:@(]=u\48BNG0*C
!AGYl%NLpAe;KY'V<,n(31^Rs02jG@g?Df14b(!1M_9n<AEj?r_hfa1c?*J'Pp<I@tKk8":F@nrI
BM(d8D99VdDfeLtelCq+S4G4kHt@c,,Jf^SS7O=kmdn>]H<n[VDN)>'\#E5dAtMg:";`EfR>N)hT
9hT";ijkYHd7)D:U!`+qJor1ZMe#bQ&%LJe0X2-N$!h(WqX2;H;qYX\)(F+"P\DC*MXRe?bG%m"f
Y_go4NO\EJj<i;48/[h;m/U`\]_Lrp8[)3]]q1Pin"1^2ol/2aucIVT(d;S`Iai5e@EC-R;'j9C"
N79m4qQ%o3"KO#XtP/u3[nSo5Mf?d(^/[J@jr=s'=kT9`$Waa#U_ND:okpoS)gdYq;U]=hiV3Wqb
:;=WrmkYunaFcK3:SPHN>4/3!hE?&;OiuoNjG0X<i``r#*g[<OIeGL:gN18\\18<.\G7?T/%V7jW
PK02]du@;+nmuA:0NM8fMWUF#_m3+"@H+fAe*-Ep7PH^aj!ujk+#kDV0($EDS?"l0l,X/t-M]c43
YCSK8p#Q.I\B#uF.tfJk`]SjB]FWNh6SaOa*q4K#t4.o:#SZmS#c*k6cpO2)DPpnZhSQ4e@YY3<U
l<0p00&hL9udrg@i#E=H\^u;C&q'XK)QnomheLgK?;N6g9A)o*I"=38G(UpNF^UIKqGGN+*L)G;k
R(QiZf<0hAbj^2&[6decqG$P;R./rYI$UtbaH/?+MMK=Hc%lB!(Ph_$`q/_bNg$&F;rm;^J#T6YA
llE"`'5';TW+U1cmY-0;'Z>jGH`nYABSCl1-`JfO$Q6,U[SYW1jEK#"OQasGk`gDuOPZeU(d?K=0
`&57_c8#1'i:2ql'%FUAafC!Ch9oU+`9?k(Qpr?tGB@U':?bWK2tNIm8;eeaJPZ-FX48F35ssLNF
S`=QXV-)!c8]omUdcV>J*uopr19b.\#WKlX^pT)$qaFr?!;!C^V#NO`'U4/XQD-nrr=ZX#>iq&ZK
BB3e#9`%]Q9f1)"krlS_&4T`pguGWBF"QHF#[/2"^o=Q%u=E+%fi:A]%i$q2CZO1FnJ=#=k`\;C#
Tfo:>h^8#=p:-ek_m/mTL)>u5)#TY2tZ!kZW2P*eZ,A6tcD2fFYFUYR!=rc79?X8)+L9)EP+/$3f
JplG(L%#Dg44Pt9ko_NBJeO`8<86/^f5=5@@PXE^a-Yh[k*gFeU2/eSnJ#W'L[l$=2dIEb@r#\Zi
mHk9V%$\L7dcR1$d?7t<%^/KS>Z6%-N+8h.dR?>Vml5Xn!/Kh2<Gk;e\hffoN=hKIn2d/=;6'$WH
cd@9`.,49?TC?f7Hsq8gJ2l\0ccS2gp/6*c_W)5/>TSPit(7498O_hEhq"O0TArCo8P&W:=);uQQ
g4*5DE`%Vf]"b%iJ]?qL;*PoW<J`<CLCelT.(#]+4+5Gf^#/\fEjB^ocWmV5fJ&^KAmG-E=<C:m]
V=F,Ddl,88q*ETPXaqG]Xu2.Bo/Qg%1<?G5JW=)a6c06D292U8a$S,(+#4B:KVUc"U_3DFS7mb'F
D`HlN!"M8fAPTiS&9%[02IGIYtdtF)kGeR":6KopR@A^)S0kE7+5e#jrbH!)CpjHXG1nOE.miqH$
;-oi^IqCAi'pDKZHqqc&?Vfu`eDg`YNl8o,_?S0SnUaZ'[Jo@76@Ba<T2P%+@k++\@pRm&QoiUY/
PY&4H;KP1g:AdNW?0`Kdusf8at!k_LE#f^.jUE%I(1C;[OW][W)9sBL9]2h^9.QPO)`MHHfG588i
>EjlR18^BO,/_1t>BugO&IT4-=FLhJi,e!&f7/P_'pg*9U*q>4!u'EdH9dL-==MF]XXX(G0W2F,W
h/Fa7koECZTVINB<Jj#2QOE0QsgK!iTr+D@T/X@TI7Z)3P$L,%pa%?CF.jojs4Xn+/4+S@*-qJRh
4Y3Vjnc;ns2QfD3H#O4sf/_>.nIE/cGbMHo%6r)f]D=a$R9T5_plLs<ZdJQ,(pQ=@&o:7Rl#aGt-
K38(A?6-1hr\O'5lZE-?:/b7b<"A)#29%2^8O.XJ:W*/5@GMI5oq5U^H!*Y>%\t6?$-#E6;iDJ9T
ms-@Ak$Tn-D,Mtmd)cOS#Og%[&\T:GN.Y*8G%i)o^5sb:20jK)WJ9LZ2Xg3Vs:F*DrIhd)L6C;R2
JAhOG90-nV^;'(qNN5Qg$h&.[ri,1qI^lRms#Q.\kIFe7<AX\O30H7^^4s3cVU1NH87#d?Vc<nUt
u/$tCb=aiWq+0sB=bq/a.8"`DOtoI&0\9<$t&:)2%%OWL@<jP\2+9W3mf*e#lLr"NT=mGBqNrY/?
"FjJ?Kk7%@j>b*98;>O(HXtd[@7rpST6qui_jcP2ZnXH?_mZCacK95<B;Z0h%R;SMr/7MrNZF22j
=KlJFDg._rj^EWeMOLk-dKd?P[D.YOklr?85.-0HI4h&OGH@FU/;Utn`9h7"(WL?T.J-3bZhqc=U
tsLgot/@[OR6X;=ec$@XiSHGF[%KdS]nKWN(&pAQ$^]m:<0CgV'uWnoYl),nV?WT=EG\-M^'D?I_
5BHMpsqZQ"8])nEIuign21LNRa7[#qgU<0B+%\e]3EL3e)mSH:pH&.@OWBYc`e),K[\_@60r>eDi
O0]$m24iMG.ZTMMjB[-'%90>_hG$Not[\!?Pe`*Q9'/4%jrc4d(l2^rXSa7V)j?jTFT\!&;-;5tl
g:$M"]q2SgemaK(75*_21Z6m)mh8OR$ji!5B$Vru)f_"3/*+!?tVuUAI7QiMNTh[C,%[jaLf/8ei
1+BDLJi8s?&>)4XM)QL4.F)W!#YBt9dSnIi2>^^+$sgHaPWkDL%(J8?%_m!62m.:hS7G&j2BFU9d
\ngnn.P/Fq_Db(V[DfM,Q@bljLUG"]Lfd#h\9nFT1k2QLAq!q7lO<\@fku9DP'B5HVjC+?][l<2e
Sa"Vb)hK<Q7(!%\-t.i;Rb:)Edu'Nh"["`f+f-8aM_?)*nOu?3UeB:BAa$dlN1_Dk##"r%+Nbr7N
do(Nkq+F7m+Mg6L@(h*.&p<MARQk?H>MOH'.Xr.3V\UDo*MLJ3ZQVdDGT,c!7-Vp81O$Uk=*)"^e
5ZHAPAUh_kNq>mL4]\2&VnBI$of`~>
endstream
endobj
6 0 obj
<<
/Filter [ /ASCII85Decode /FlateDecode ]
/Length 45
>>
stream
Gap9`0JkUF;(qu:F&N/pSXcaRn!Rsai0(8T#QSGa#/g~>
endstream
endobj
xref
0 7
0000000000 65535 f 
0000000031 00000 n 
0000000051 00000 n 
0000000110 00000 n 
0000000159 00000 n 
0000000296 00000 n 
0000109317 00000 n 
trailer
<<
/Size 7
/Root 3 0 R
>>
startxref
109452
%%EOF
\end{filecontents*}
% END STANDALONE DEPENDENCIES

\documentclass[12pt]{iopart}
\usepackage[OT1]{fontenc}
\renewcommand{\rmdefault}{cmr}
\renewcommand{\sfdefault}{cmss}
\renewcommand{\ttdefault}{cmtt}
\usepackage{graphicx,array,booktabs,placeins,needspace,xcolor,eso-pic,hyperref}
\newcolumntype{P}[1]{>{\raggedright\arraybackslash}p{#1}}
\Urlmuskip=0mu plus 2mu
\hypersetup{hidelinks,pdftitle={The Saint Who Left Was Never Free: Religion, Suffering and the Body},pdfauthor={Sohel Georgiadis},pdfsubject={Philosophy}}
\begin{document}
\AddToShipoutPictureFG*{\AtPageLowerLeft{\put(17,110){\rotatebox{90}{\parbox{510pt}{\color[gray]{0.50}\fontsize{13}{15}\selectfont ``The beauty of a thing does not depend on whether it lasts.''\quad Sohel Georgiadis\quad 25 Apr 2025}}}}}
\title[The Saint Who Left Was Never Free]{The Saint Who Left Was Never Free: Religion, Suffering and the Body}
\author{Sohel Georgiadis}
\address{Student, The Reality Public School\\Kokrajhar, Assam, India}
\ead{\href{mailto:sohelgeorgiadis@gmail.com}{sohelgeorgiadis@gmail.com}}
\vspace{10pt}
\begin{indented}\item[]25 April 2025\end{indented}
\begin{abstract}
This essay examines the prestige of renunciation and the association of spiritual authority with suffering. It asks how visible sacrifice comes to stand for inner freedom, and how bodily life, joy, and sexual desire become objects of suspicion. Through comparisons drawn from religious and philosophical traditions, it develops a distinction between freedom within experience and the performance of escape from it. The argument treats the body's impermanence as a condition of life and questions the move from impermanence to contempt.
\end{abstract}
\vspace{-6pt}
\begin{indented}
\item[]\centering\includegraphics[height=133pt,width=.90\linewidth,keepaspectratio]{The_Saint_Who_Left_Was_Never_Free-cover.pdf}\par
\vspace{4pt}
{\footnotesize Caspar David Friedrich, The Monk by the Sea. Alte Nationalgalerie, Berlin. Public-domain reproduction via Wikimedia Commons.}
\end{indented}
\clearpage
\section{The Saint Who Left Was Never Free}

\subsection{Renunciation and the promise of another world}

Why have almost all prophets, godmen, saints and religious teachers spoken against this world and in favour of another world nobody has ever seen? Why humanity respected those who rejected life more than those celebrated it?

Wait. Let me say answer this more simply before the answer becomes too decorated by philosophy and saturated by advanced vocabulary. I may be simplifying the history of religious institutions here, and I don't want to be ignorantly dishonest about that. Maybe I am attacking the institution because I have not yet understood devotion from the inside of how a religiously devoted person understands. But I know how it sounds, and I am still saying it. I remember hearing people in the school corridors use the word mok\d{s}a like it was the only respectable desire, while the same people lowered their voices when sex or the body came into the conversation. That lowering of voice is also a kind of teaching, just a hidden one that people didn't believe as hidden anymore.

The state of being sad is not an accident or some form of enlightenment, and I want to begin precisely there, with that single recognition: this is not some sort of accident that happened unwantedly. When a tendency appears across thousands of years, across cultures that have never met each other, across geographies separated by oceans and enormous silences of unmapped history, across people who share not a single word of the same language and not a single image of divine, something that unanimous, something that consistent, cannot be dismissed as a collective error in thinking, as a mistake that everyone everywhere made independently and simultaneously. It has a logic. It has a very deep and very human logic, even when that logic is ultimately rooted not in wisdom but in something far simpler and far sadder than wisdom.

But to understand it fully, we have to be willing to take it seriously first. We have to understand it from inside, the way a man who has grown up inside a tradition understands it, not from the comfortable distance of a critic who dismisses what he has never genuinely entered. To dismiss it is easy. To understand it is the only thing worth doing.

So let us begin at the very beginning. Let us begin where this tendency or urge begins, which is in the plain, unavoidable, actual quality of ordinary human experience.

The ordinary human being, and I mean the vast, overwhelming majority of men and women who have ever lived on this earth , experiences life primarily as a difficult thing affair. And I do not say this with any sort of bitterness or cynicism; I say it simply as a fact that anyone who looks honestly at the human situation cannot avoid. He is born with enormous longing. This is the first thing. The newborn infant is already reaching. He reaches for warmth, for nourishment, for the sound of a familiar voice, for the specific flavour of milk, and pressure of his mother's body. And this reaching, this fundamental and constant reaching, does not diminish or fades with age; it only grows more elaborate and more demanding. The infant grows into a man, and the longing grows with him. He wants joy. He wants security that does not depend on anything outside himself and yet can only apparently be found outside himself. He wants to be loved in a way that is total and unconditional, and yet every love he actually encounters is partial and conditional. He wants to matter, to leave some mark of his presence that persists after he is gone, to have been somehow necessary to the story of the things. And life, as he actually encounters it from day to day, delivers all of life, as he actually encounters it from day to day, delivers all of these things in small doses, briefly, unreliably, always mixed with their opposites, always moving just slightly beyond the reach of his hand.

The joy comes but it does not stay. The love arrives but it also hurts, and it also changes, and it also ends. The security he builds with such effort is always one illness or one accident or one conversation away from collapse. And the beauty he finds in the world is inseparable from the fact of impermanence, from the knowledge, always present somewhere in the back of his awareness, that this beautiful thing will one day not be here, that even the one who is right now noticing the beauty will one day not be here to notice anything at all.

Death stands at the end of every corridor. And this is not a metaphor. It is the single most consistent fact of human experience, present from the very beginning of self-awareness, irreducible, non-negotiable, the one thing about which there is not and has never been any exception in the entire history of biological life on this planet.

This is the ground condition of ordinary human existence. And from this ground condition, a very natural psychological response arises, which is the:

``Desire for Something Permanent, something that cannot be taken away, something that does not follow the law of loss and decay that governs absolutely everything in this world.''

This desire for permanence is not philosophical in its origin. It is not the product of careful reasoning either. It arises directly from pain. It arises from the experience of having loved something and having lost it, of having wanted something deeply and have found that obtaining it did not satisfy the wanting, of having built something carefully and having watched it come apart. And from this lived, bodily experience of impermanence arises the oldest and most persistent spiritual impulse in human history: the reach toward the eternal.

Now the prophets and the saints and the religious teachers did not invent this impulse. They found it already fully formed in the people they addressed. They found human beings who were already wounded by life, already exhausted by the gap between what they had hoped for and what existence had actually delivered to them. And the prophets, being sensitive men, being men who had in many cases lived through some version of this same wounding themselves, addressed it directly. Not with philosophy, not with medicine, not with political change, but with a complete reframing of the human situation. And the reframing went like this: the reason you are suffering is that you are looking for the permanent in entirely the wrong place. The permanent is not here. It is there, beyond the reach of the senses, beyond the reach of ordinary experience, in a world that this world is only a poor imitation of. Turn away from this. Turn toward that. Give up what cannot satisfy you and receive what cannot be taken from you.

And the people who heard this felt something like profound relief. Because they had been blaming themselves. In the privacy of their hearts they had been thinking that perhaps the problem was personal, that perhaps they were not clever enough or strong enough or pure enough to extract from life the lasting happiness they were seeking. They had been carrying a secret shame about their own restlessness, their own incapacity to simply be content with what was in front of them. And then a man appears and says, with the full weight of spiritual authority: the problem is not you. The world itself is the problem. The very ground you are standing on is inherently unstable and inherently incapable of delivering what your deepest nature is actually seeking. And there is a very immediate and very powerful relief in hearing this, because it removes the burden of personal failure. If the world is maya, if it is fundamentally not the place where real life happens, then my inability to find genuine happiness in it is not a failure of my intelligence or my courage. It is simply evidence that I am sensitive enough to have noticed what the world actually is.

This is the deep psychological mechanism at the core of world-rejection. And I want to say this clearly: the mechanism is compassionate in its origin. The prophets who taught it were not, most of them, frauds or manipulators, though the traditions built around them have often attracted both in great numbers. The sincere ones were genuinely trying to address a genuine wound. They were doing what any decent doctor does when a patient arrives in severe pain: they were offering the best remedy they knew. The problem is that the remedy was not a cure. It was a painkiller. A painkiller so powerful and so immediate in its effect that generation after generation reached for it, and the reaching became a tradition, and the tradition became a civilization, and the civilization taught its children that the painkiller was the medicine and that the medicine was the highest possible achievement of the human being.

But let us go deeper than psychology, because the psychological explanation, though real and important, is not the whole story. Alongside it there is also a philosophical dimension to this world-rejection, and it has its own internally coherent logic, its own beauty even, even as I say it is ultimately built on a misreading of the available evidence.

Plato observed something that no careful thinker can dispute: everything we encounter in this world is impermanent. Every object, every feeling, every relationship, every condition is a changing thing, a thing in process, a thing that is always in the act of becoming something other than what it is at this exact moment and that will eventually cease to be altogether. And Plato asked a question that sounds simple but has consequences that ripple through the entire subsequent history of Western thought: if everything we actually encounter in experience is always changing, then where does the idea of permanence itself come from? We have in our minds a clear concept of a perfect circle. We understand it with complete precision. But we have never seen a perfect circle in the physical world, not once in all of human history. Every circle drawn by any hand or any tool is slightly off, slightly asymmetrical, more approximation than pure form. So where did the idea of the perfect circle originate? It cannot have come from the imperfect circles of the physical world, because the lesser cannot generate the idea of the greater, the copy cannot contain what is only in the original. Therefore, Plato argued, the perfect must exist somewhere other than here, in a realm of eternal forms, of which everything we see in this world is a faint and flickering copy.

This is the foundation of the entire Western spiritual project of reaching toward another world. And notice what it does to the spiritual life: once you establish that reality is elsewhere, that the real is not here but there in a dimension beyond the senses, the whole structure of the religious life becomes the structure of movement toward elsewhere and away from here. The men who have turned away from the ordinary world and devoted their entire existence to the pursuit of the real become the heroes of this worldview. And the man who says that this here, this actual, immediate, impermanent, mortal, beautiful, terrible human life, is not only sufficient but is in fact the only place where the real can ever actually be found, that man becomes the man who is confused, or dangerous, or simply not yet serious enough to have understood his situation.

In India, a very similar philosophical move was made through entirely different vocabulary. The Vedantic tradition said that the world of multiplicity, the world of separate beings and changing relationships and desire and loss, is a superimposition upon the one underlying reality called Brahman. Brahman is without qualities, without location, without any characteristics whatsoever. It simply is. And the world of ordinary experience is like a wave on the surface of an ocean: it appears to have its own independent existence, its own shape and movement and direction, but it has no existence separate from the ocean beneath it. When the wave subsides, it does not go anywhere; it returns to what it always was, which is simply ocean. And the goal of the spiritual life, in this framework, is to recognize oneself as the ocean rather than as the wave.

This framework too, in the way it was most commonly taught, placed the real beyond the reach of ordinary life. The man who withdrew from the world was moving most decisively toward this recognition. The man who lived fully in the world was most thoroughly caught in the illusion. And the institutional religions that formed around these ideas reinforced this hierarchy at every turn, because the institutional religion needs a hierarchy. It needs a clear distinction between the spiritual and the worldly, the sacred and the profane. Without this distinction, there is no institution. There is no one to teach and no one to be taught, no one to grant forgiveness and no one who needs forgiving.

So you have, from both the Western and Eastern philosophical traditions, a powerful and internally rigorous justification for the rejection of this world. And when sophisticated philosophical frameworks align with and reinforce a psychological tendency that is already fully present in the suffering of ordinary human life, the result is a cultural current of extraordinary force that carries almost everything before it for centuries, sometimes for millennia.

But here is what I want you to see with great clarity: neither Plato nor the Vedantins were simply wrong. This is important. They were pointing at something genuinely real. The experience of impermanence is real. The observation that ordinary sensory experience does not deliver final satisfaction is a genuine and honest observation. The mistake is not in the observation. The mistake is in the conclusion drawn from the observation.

The conclusion drawn was a binary one. Either the permanent is here and we can find it through ordinary sensory engagement, or it is not here and we must abandon ordinary experience entirely in order to reach it. Either this world contains the real or another world does. But there is a third possibility, and this third possibility is what Krishna embodies. The third possibility is that the permanent is not somewhere else. The permanent is the very depth of this moment, the unchanging ground on which all changing experience rests. And the man who has genuinely touched this does not need to reject anything. He can accept the impermanence of everything precisely because he has found what does not change within the impermanence. He can love completely because he is no longer afraid of loss. He can engage with the world without reservation because the world is no longer a threat to him.

Now let us look carefully at why the crowd worships the rejector and distrusts the celebrant.

There is a story worth telling here, because it illustrates the point more clearly than any abstract explanation can. In a small mountain town in what is now Turkey, there lived, according to an old Sufi account, two men who were both considered holy by different portions of the population. The first was a dervish who had lived for twenty years in a small cave outside the town. He ate almost nothing, owned nothing, spoke rarely, and when he did speak it was in short, difficult utterances that his followers spent years interpreting. His face had the quality of weathered stone: still, gray, completely composed, without the flicker of any ordinary human emotion.

People came from great distances to sit for a few minutes near the mouth of his cave and feel, as they reported, a kind of extraordinary stillness descend upon them. They left quieter than they arrived, and they called this holiness.

The second man was a baker who had his shop in the center of the market. He woke before sunrise to light his ovens, worked through the entire day shaping and baking bread, gave a portion of everything he made to whoever in the town was hungry without distinction, argued vigorously with his neighbors about the price of wheat, laughed enormously at his own jokes, wept openly at funerals, fell sick like everyone else and recovered, loved his wife with the kind of steady daily love that has no drama in it and therefore goes mostly unnoticed, and went to sleep at night in a state of simple and complete exhaustion that was also somehow a state of complete peace. He had never meditated formally. He had never studied philosophy. He had never used any of the vocabulary of the spiritual life. And yet there were certain people in the town, not many, but a few, who said that an afternoon spent in the baker's shop, listening to him argue about wheat prices and laugh at himself and slide fresh loaves across the counter to whoever walked in, left them feeling more genuinely alive than any amount of time spent at the mouth of the cave.

Most of the town considered these people either confused or romantic. A baker is a baker. A holy man is a holy man. The distinctions are not difficult.

Now I am not telling you this story to resolve the question of which man was more awake. That is not the point at all. The point is in the nature of the crowd's certainty about the dervish and its uncertainty about the baker. The crowd was completely certain about the dervish because his holiness was marked. Every visible aspect of his life, the cave, the fasting, the stillness, the withdrawn eyes, the rare and cryptic speech, communicated a single unmistakable message: this man has given up what ordinary men cling to, and therefore he has reached what ordinary men cannot reach. The entire surface of his life was organized as a legible text, and the text said only one thing: I am not here. I am there. And the crowd, trained across centuries to read exactly this text, read it immediately and with great confidence.

The baker's life was not a legible text in this sense. Its surface said many things simultaneously, some contradictory, none of them organizing themselves into a clear spiritual claim. A man who argues about wheat prices does not look like a man who has transcended the world. A man who laughs at his own jokes does not look like a man who has seen through the illusion of self. The crowd could not read the baker because the baker was not performing holiness, and the crowd, which has been taught for generations to recognize holiness through its performance, is genuinely helpless before something that is real but not performing itself.

And this is the crux of it. The worship of the rejector is, at its most fundamental level, a response to performance, to the visible and legible performance of having given up the world. And this performance became necessary because the crowd needed external proof of an interior state it could not directly verify. Since you cannot look inside a man and see whether he is at peace or suffering or lusting, whether he has genuinely transcended desire or merely repressed it, whether his stillness is the stillness of understanding or the stillness of defeat, you need some exterior sign. And the tradition, across its long history, has settled on the signs of renunciation as the only reliable exterior signs of interior achievement.

But renunciation can be performed by anyone. And the performance of renunciation, maintained long enough and consistently enough, even begins to feel, to the man performing it, like the real thing. This is the deepest danger. Not that the crowd is fooled, but that the performer is fooled, that the man who has put on the garments of world-rejection begins to mistake the garment for the state it was supposed to indicate.

Kierkegaard, the Danish philosopher, wrote about what he called the three stages of existence: the aesthetic stage, where a man lives for pleasure; the ethical stage, where a man lives by principle and duty; and the religious stage, where a man lives in a direct relationship with the divine that transcends both pleasure and principle. And Kierkegaard observed something very important: a man can be at the ethical stage and look, from the outside, almost identical to a man at the religious stage, because both are living by something beyond mere personal pleasure. The ethical man has given up the life of mere sensation. The religious man has given up even the comfort of ethical certainty. But from the outside, both look like men who have given up something, both wear the marks of seriousness, both carry themselves with a certain gravity. The crowd cannot distinguish between them.

And Kierkegaard's point, though he did not put it quite in these terms, is that the crowd tends to honor the ethical man most, because the ethical man most clearly resembles the popular image of the holy man. He has given up pleasure. He follows rules. He sacrifices for others. He is serious. All of this is visible and legible and confirmable. Whereas the genuinely religious man, in Kierkegaard's sense, may sometimes do things that look to the ethical man like ethical violations, because his relationship to the divine is direct and personal and cannot always be translated into universal rules. And the crowd, unable to see the direct relationship with the divine, sees only the apparent ethical violation and is disturbed.

This is exactly the situation of Krishna on the battlefield, and exactly the situation of Jesus with the prostitutes and the wine, and exactly the situation of every genuinely awakened being who has ever lived inside a world that could only read the marks of awakening and not awakening itself.

The worship of life-rejection is also, at its most fundamental level, a form of collective self-hatred. Not recognized as such, not experienced consciously as such, but functioning as such with extraordinary efficiency. The ordinary man who prostrates himself before the saint who has abandoned everything is also expressing, through that prostration, a wish to be rid of himself, free of his own hungers and contradictions and failures to be what he has been told he should be. He does not know how to live peacefully inside his own skin. His desires trouble him. His anger frightens him. His sexuality shames him. His grief overwhelms him. His happiness feels unearned. He moves through his life in a constant low-level war with his own nature, and the man who appears to have won this war, who appears to have successfully transcended or extinguished everything that the ordinary man finds most difficult in himself, that man becomes the object of not just admiration but worship, because the worship is really the crowd projecting onto him its own most desperate fantasy: the fantasy of a self finally and completely free from the burden of being a self.

Krishna disturbs people, and has always disturbed people, because he refuses to offer them this fantasy. He offers something more demanding and more real, which is the possibility of being completely and without reservation at home in what you are. Not free from the burden of selfhood, but free within it. Not arrived at some other world, but fully arrived at this one.

And this, finally, is the answer to your question. Humanity has worshipped those who rejected life because humanity has, in the majority, found life to be an experience it does not know how to bear. The rejection of the world is the spiritual correlate of a man who, finding himself in a room he does not know how to inhabit, simply leaves it and calls leaving a form of wisdom. And the crowd, which also does not know how to inhabit the room, honors the man who left, because he has confirmed their deepest suspicion: that the room was always the problem and not themselves.

The man who says the room itself is fine, that the issue is only the quality of your presence in it, that the light in this room, if you have genuinely learned how to see, is the same light that shines in every world you could ever dream of, that man will wait a long time for his worshippers. And when they finally arrive, it will not be because he gave them what they wanted. It will be because he gave them what they actually needed, which is a completely different thing.

\subsection{The image of enlightenment}

If enlightenment brings joy, why do enlightened men in paintings and stories always appear serious, detached and sorrowful?

This question touches something very real, and it deserves an answer that runs in more than one direction spontaneously, because the truth here is layered in a way that a single direction cannot reach.

The first and most obvious layer of the answer is this: most of the paintings and stories you are referring to were not made by enlightened people. They were made by unenlightened people, often very talented and very sincere unenlightened people, trying to represent in visual or narrative form a state of consciousness that they had not personally experienced and could only imagine. And imagination, when it tries to represent something it has never encountered, invariably falls back on the categories and the associations it already has available. And the available associations, for the ordinary unenlightened mind trying to imagine the highest possible state of human consciousness, are the associations of seriousness, of gravity, of detachment, of the visible withdrawal from ordinary human engagement.

Why are these the available associations? Because the ordinary mind has been taught, through every cultural channel available to it, that depth is indicated by seriousness and shallowness is indicated by lightness. That understanding is accompanied by solemnity and ignorance is accompanied by laughter. That the serious matters of existence require a serious face, and that a man who has looked most deeply into the serious matters of existence will have a face that carries the weight of that looking. And so the painter who wants to represent a man who has understood everything paints a man who looks as though he is carrying a tremendous weight, because the painter, who has not understood everything himself, cannot imagine how someone who has understood everything could look like anything other than a man bowed under the weight of understanding.

This is the imagination of ignorance projecting its own categories onto a reality it cannot yet see clearly. And it has had consequences that extend far beyond the visual arts. Because these images have shaped, over centuries and across cultures, the expectations of millions of spiritual seekers. Young people who embark on a genuine inner journey look at the images they have been given of the destination and model themselves accordingly. They practice serious faces. They practice measured and deliberate speech. They cultivate an air of gravity and detachment. And after many years of practice, they have very thoroughly developed a very convincing performance of the state they imagined they were supposed to arrive at. And the culture around them, which can only read the performance, receives the performance as the genuine article and confirms them in it. And the performance, maintained long enough, begins to feel, even to the performer, like the real thing.

This is one of the deepest traps in the entire history of organized spiritual practice: the tendency of the seeker to mistake the costume of the destination for the destination itself, to become so thoroughly identified with the performance of holiness that the actual quality of aliveness that genuine inner development produces is never given the chance to emerge.

But let me now go to the second layer of the answer, because the first layer, though important, can suggest that all solemnity in the enlightened is simply performance, and this would be both unfair and inaccurate.

There is a genuine quality in certain states of deep awareness that can, from the outside, be read as seriousness or as a certain kind of gravity, though it is not what those words ordinarily mean. When a human being has genuinely arrived at a deep understanding of existence, when they have genuinely passed through the full weight of what it means to be alive in this world, when they have looked without flinching at impermanence and death and the suffering of sentient beings and the extraordinary fact that any of this exists at all rather than nothing, something in them does become very quiet. Very still. And this stillness is not the stillness of sorrow. It is the stillness of a lake so deep that the surface movements, which on a shallow pond would disturb everything, are contained at the surface without reaching the depths.

But from the outside, deep stillness and deep sorrow can look very similar. The grief-stricken man is still. The enlightened man is still. The man who has given up on life is still. The man who has fully arrived at life and found it to be, at some ultimate level, completely all right, is also still. And the observer who has not developed a very particular kind of perceptual sensitivity, the kind that reads the quality of the stillness rather than merely its presence, cannot distinguish between these.

There is a story told in the Zen tradition about a monk who spent forty years in a monastery practicing meditation with enormous seriousness and dedication. At the end of forty years, he went to his teacher and said: I have been practicing for forty years and I have achieved something. I no longer feel anger, I no longer feel desire, I no longer feel fear. What have I achieved? And the teacher was quiet for a moment and then said: you have achieved what a stone achieves without any practice at all. The monk was completely silent, and then he began to laugh. And the laughter went on for a long time and had a quality that the other monks in the monastery reported was unlike any laughter they had heard before.

Now what happened in that moment? The monk had spent forty years achieving, through enormous effort and discipline, a state of emotional flatness that he had mistaken for equanimity. He had suppressed anger, suppressed desire, suppressed fear. The suppression was real; the equanimity was not. Genuine equanimity is not the absence of emotion. It is the presence of an inner spaciousness large enough to contain all emotions without being overwhelmed or defined by any of them. The man who has genuine equanimity feels everything, feels it fully and cleanly, without the secondary reactions of resistance and clinging and fear that ordinary consciousness adds to every primary experience. But the feeling itself is not eliminated. If anything, it becomes more vivid, more immediate, more present, precisely because it is no longer being filtered through layers of protective distance.

And this is why genuinely enlightened people almost always have a quality of emotional aliveness, a quality of being genuinely affected by what they encounter, that the man who has achieved the religious performance of detachment does not have. Ramakrishna wept and laughed and fell into ecstasy. Anandamayi Ma, the Bengali mystic of the twentieth century, was described by everyone who met her as being in a state of constant, effortless joy that was utterly contagious, utterly natural, completely free of performance. Nisargadatta Maharaj, another twentieth-century figure, was known for getting genuinely angry when asked questions he considered dishonest, genuinely warm when someone arrived at his presence in genuine distress, genuinely playful when the moment called for it. He was not managing his responses. He was responding.

And this responsiveness, this quality of being genuinely and completely present to whatever is actually happening, is in fact the outer sign of the inner freedom that genuine enlightenment represents. The man who has genuinely arrived does not need to maintain any particular emotional posture, because he is not afraid of any emotional state. He is not afraid of anger because he knows it will pass. He is not afraid of grief because he knows it will pass. He is not afraid of joy because he is not clinging to it. He is simply present to whatever arises, and his presence has a quality that those around him can feel directly, even if they cannot always describe it accurately.

The paintings of the enlightened as sorrowful and detached are paintings of a concept of enlightenment rather than of enlightenment itself. They are paintings produced by minds that have understood enlightenment as the maximum possible distance from ordinary human experience rather than as the maximum possible depth of presence to it. And this understanding, this fundamental misreading of what the destination actually is, has shaped the spiritual imagination of most of humanity for most of its history.

But here is the thing I want you to hold with care, because it is both honest and important: the misreading is not without some basis in reality. There are moments, particularly in the early stages of genuine inner development, where a person does become quieter, more still, less reflexively engaged with the ordinary social world around them. This is real. The mind that is beginning to learn how to rest in itself rather than constantly reaching outward does, in that period of learning, sometimes appear to withdraw. The man who is genuinely engaged with a deep inner movement has less energy available for the performance of ordinary social engagement. He may sometimes appear detached, not because he has achieved detachment but because his attention is genuinely elsewhere.

And the tradition, observing this, has taken this intermediate state as the destination. It has looked at the sincere but still-developing seeker and said: this is what the finished man looks like. This quietness, this apparent withdrawal, this reduced engagement with the ordinary world, this is the mark of the fully realized. And having said this, it has frozen the tradition at the intermediate stage and produced generation after generation of sincere practitioners who achieve the intermediate stage and stop there, believing they have arrived, because everything around them confirms that the intermediate stage is the destination.

The genuinely arrived man, who has passed all the way through the intermediate stage and emerged on the other side into something that looks nothing like detachment and everything like an extraordinarily vivid aliveness, that man disrupts the tradition's image of its own highest achievement. He is inconvenient. He laughs when he is not supposed to laugh. He is moved by things that the man of proper spiritual gravity has learned not to be moved by. He is, in short, too alive. And the tradition, which has built its entire self-image around the performance of a particular kind of not-aliveness, does not know what to do with him.

So it either ignores him, or it edits his image after his death into something more acceptable, or, in the most extreme cases, it simply does not recognize him as having achieved what the tradition values, because what the tradition values is the performance of a state that he has moved so completely beyond that the performance itself has become unnecessary.

This is why the paintings show the sorrowful detachment: because the tradition that commissioned the paintings had long since lost contact with the joy it was ostensibly pointing toward, and was painting, with great sincerity and great skill, its own image of its own highest achievement, which was an image of a man who had successfully completed the performance.

\section{Suffering Has Become a Religion}

\subsection{Suffering and spiritual authority}

Why does sadness appear holy in almost every religion? Why are laughing saints distrusted while suffering saints are worshipped? Why do people feel closer to a crucified Jesus than to a dancing Krishna?

There is a very ancient equation in the human mind, and it operates so deeply and so automatically that most people are not even aware it is there, functioning with the same invisible reliability as the assumption that the sun will rise. The equation is this: that which costs nothing cannot be worth very much. And its corollary: that which gives pleasure must, by its very giving of pleasure, be suspect. From this equation, which is not argued for or arrived at through careful thinking but is simply absorbed from the cultural air one breathes from childhood, follows the entire spiritual hierarchy that places suffering above joy, gravity above laughter, withdrawal above engagement, and the sorrowful saint above the dancing one.

But let us understand where this equation itself comes from, because it too has a history, and the history is revealing.

Human beings learned very early, as individual animals and as social groups, that survival requires the willingness to defer pleasure. The man who eats all his stored food in the first week of winter will not survive to spring. The tribe that celebrates when it should be preparing will find itself unprepared when the environment turns hostile. Delay of gratification is not merely a moral virtue; it is a biological survival strategy, written into the architecture of the human nervous system through hundreds of thousands of years of selective pressure. The capacity to say no to an immediate pleasure in service of a future necessity is one of the most fundamental human capabilities, and it is, in a very real sense, the foundation on which civilization itself has been built.

Now this biological and social reality, this genuine and adaptive capacity for self-denial in service of long-term survival, was taken up by religion and spiritualized. The act of giving up a pleasure was translated from a practical survival strategy into a sign of spiritual seriousness. The man who fasted was not merely preserving food; he was demonstrating his mastery over the body's demands. The man who remained celibate was not merely avoiding the complications of sexual relationship; he was demonstrating his freedom from the most powerful of the biological drives. The man who gave away his possessions was not merely simplifying his practical life; he was demonstrating his indifference to the material world.

And in each of these cases, the demonstration was legible. It could be seen and verified and honored. And because it could be honored, it was used as the primary currency of the spiritual marketplace, the thing you presented in exchange for recognition as a serious person, a deep person, a person whose inner life had achieved something beyond the ordinary.

But something very important happened in this translation from practical necessity to spiritual currency: the suffering itself, which had originally been merely a side effect of the necessary self-denial, began to be valued for its own sake. If fasting proves spiritual seriousness, then the suffering of hunger that comes with fasting is evidence of that seriousness. And if the suffering is evidence of seriousness, then perhaps more suffering is evidence of more seriousness. And perhaps suffering that is not attached to any practical purpose, suffering that is purely gratuitous, purely chosen, purely an end in itself, is the highest evidence of all. From here it is only a very short step to the flagellants of medieval Europe, who beat their own bodies in public as a demonstration of devotion, or to the extreme ascetics of the Jaina tradition, who stood naked in all weather and refused to brush away the insects that bit them, or to the countless traditions of voluntary self-torture that have appeared in nearly every religious culture across the entire history of humanity.

Now here is the thing I want you to see about all of this: the laughter of a genuinely realized man sits completely outside this entire framework and therefore cannot be processed within it. The laughter of such a man is not the laughter of a man who has not suffered. It is not the laughter of innocence or ignorance or comfort. It is the laughter of a man who has gone through suffering completely, all the way through it, and has arrived at the other side of it at a place that is so completely and radiantly all right that the only adequate response is a kind of joy that has no name in the ordinary vocabulary.

This kind of laughter is, in its own way, the ultimate evidence of the deepest possible spiritual achievement. But it cannot be recognized as such by the crowd, because the crowd's vocabulary of spiritual achievement is entirely built around the marks of suffering. The crowd reads suffering as depth. It reads laughter as shallowness. And the man who laughs, therefore, appears to the crowd as a man who has not yet understood the seriousness of the situation.

Ramakrishna, the great nineteenth-century Bengal mystic, laughed constantly. He went into ecstasies while doing completely ordinary things: while eating, while talking with visitors, while watching a theatrical performance, while observing a flock of birds cross the sky. He wept with the same freedom and the same totality. He once became so absorbed in the devotional feeling that he climbed onto the altar and merged with the image of the goddess he was worshipping. He did not maintain the careful dignity that his position as a temple priest demanded. And for a long time, the people around him genuinely did not know whether he was deeply holy or simply insane. The two possibilities looked almost identical from the outside, because the marks of deep holiness and the marks of madness both include a disregard for conventional behavior and an inability to contain one's emotional responses within the socially expected limits.

And this is an extremely revealing observation: the crowd cannot easily distinguish between the man who has gone so far into existence that he has come out the other side into a freedom that looks like madness, and the man who has simply lost his grip on ordinary sanity. Both laugh when the situation does not call for laughter. Both weep when the situation does not call for weeping. Both are apparently indifferent to the opinions and expectations of the people around them. From the outside, both look like men who have stepped off the ordinary map of human behavior.

But the quality of the stepping off is utterly different. The man who is mad has lost his center. He is adrift. His unpredictability comes from the absence of any reliable ground beneath him. His laughter comes from the disconnection of his responses from any coherent inner reality. Whereas the man who is genuinely free has found a center that is so stable, so completely solid, that the ordinary social map of expected behavior becomes genuinely unnecessary. His unpredictability comes not from the absence of ground but from the presence of a ground so deep that it does not require the surface structures of conventional behavior to maintain itself.

These two are, experientially, as different as anything can be. But visually, from the perspective of the crowd, they are almost indistinguishable. And because they are almost indistinguishable, the crowd, which cannot afford the risk of being wrong, defaults to the safer option: it trusts the man who suffers and distrusts the man who laughs. Because the man who suffers, whatever else he is, is clearly taking something seriously. He is at least in the right emotional register for a serious engagement with the serious business of existence.

Now let us look directly at the question of Jesus on the cross versus Krishna at the dance, because this is the sharpest and most vivid form of the larger question, and it illuminates something that might otherwise remain somewhat abstract.

Jesus on the cross is an image of pure and voluntary suffering. Whatever one believes about the theology surrounding this image, the image itself is unambiguous: a man who could have avoided this, who had the power and the intelligence to step aside from this, has chosen to go through it. He is suffering maximally. His physical pain is extreme. His psychological pain is perhaps even more extreme, that cry in the darkness, the sense of having been abandoned by the very divinity he had trusted completely. And he is dying. All of this is visible in the image, written in the most legible possible terms on the body of the man hanging on the cross.

And the crowd that stands before this image recognizes every element of it. They know what it is to be in pain that they did not choose and cannot escape. They know what it is to feel abandoned by the universe in a moment of great need. They know what it is to face the approach of death. They recognize Jesus on the cross because he is doing, in the most extreme and most public possible form, what they themselves do every day in smaller and more private forms. He is suffering. And his suffering is not shameful. It is holy. And if his suffering is holy, then perhaps their suffering is also holy. Perhaps their daily pain, which they had experienced as evidence of failure or weakness or unworthiness, is in fact evidence of the same willingness to engage with existence that made Jesus the figure he is.

This is the gift that the crucified Jesus gives to the crowd: the redemption of their own suffering. And it is an enormous gift. One should not underestimate it. For a person living in genuine and unrelenting pain, the image of Jesus on the cross is not merely a religious symbol. It is the most powerful possible statement that their condition is not a mistake, not a punishment, not a sign of their unworthiness, but something that the divine itself has been willing to undergo. It is the universe saying to the suffering man: I know. I have been here too. Your pain is not separate from mine.

This is why people feel closer to the crucified Jesus. Not because suffering is actually higher than joy in any ultimate sense, but because suffering is where they actually live, and the crucified Jesus meets them exactly where they live. He does not ask them to come to a higher place before they can approach him. He comes to where they are. And where they are is on their own private cross, in their own private darkness, crying out their own private cry of abandonment.

Krishna at the dance does not meet you where you are, unless where you are happens to be the very unusual place of complete inner freedom.

Krishna at the dance is not an image of meeting you in your suffering. It is an image of what you could become if you went all the way through your suffering rather than stopping in the middle of it and calling the stopping a destination. It is an image of what exists on the other side of the entire territory of pain and loss and disappointment and the fear of death. It is an image of a human being who has passed through all of that and emerged into a joy that does not depend on any of it having been otherwise.

And this image asks something enormous of the person who stands before it. It asks them to accept that their suffering, real and genuine as it is, is not the final word about their situation. It asks them to believe that something is possible that they have not yet experienced and cannot yet quite imagine. It asks them to trust that the dance is real, that the music is real, that the freedom being expressed in that image is not a fairy tale or a theology or a consoling fiction but an actual possibility available to them in their own actual lives.

This is a far more demanding invitation than the invitation of the cross. The cross says: come to me as you are. The dance says: come to me by becoming what you are not yet. And most people, being where they are, find the first invitation far more accessible than the second.

But there is also this, and it is worth saying slowly: the man who finds the crucified Jesus and who genuinely meets him there, who genuinely allows himself to be met in his suffering by that image of divine suffering, and who lets that meeting go all the way through him rather than stopping at the comfort it provides, that man may eventually find himself at the dance. Because the meeting in suffering, when it is complete, when it is not merely a consolation but a genuine encounter, carries within it the seeds of the joy that the dancing Krishna embodies. The cross, fully entered, leads to the dance. The dance is what the cross is trying to reach. They are not opposites; they are two moments in the same movement.

What religion has done, in most of its institutional forms, is to stop at the cross and call the stopping enlightenment. It has offered the comfort of having one's suffering validated and has called this offer the full extent of what the spiritual life makes available. And this is the deepest possible injustice to the people the religion claims to be serving. Because if you stop at the cross, you never arrive at the dance. And the dance is the point.

The sadness that appears holy in almost every religion is real holiness frozen at its middle stage and presented as if it were the final destination. It is a man who has gone halfway and stopped. And the institution around him, which needs him to stay exactly where he is, which needs his remaining suffering to maintain its position as the authority over that suffering, tells him that where he has stopped is the highest possible point anyone can reach. That his sadness is the mark of his depth. That his sorrow is the sign of his understanding.

And the tragedy is that he believes it. And the crowd believes it. And the tradition builds itself around this belief, generation after generation, until the sadness becomes not merely a stopping point but a destination in itself, not merely a mark of depth but a proof of it, not merely an accompaniment to genuine spiritual seeking but the thing itself. And by that point, the dancing Krishna has become not a figure of the highest possible human achievement but a slightly suspicious character who probably has not understood the seriousness of the situation.

\subsection{Sacrifice and celebration}

Why do people worship sacrifice so deeply? Why does humanity respect a man who suffers more than a man who celebrates?

There is something in this question that goes very far down, much further than it appears on the surface, and I want to follow it carefully, because the answer requires us to look honestly at something that most people carry inside themselves without ever fully examining it.

The worship of sacrifice is not a simple thing. It is not merely a cultural habit or a religious teaching that could, in theory, be removed by more rational instruction. It has roots that reach into the very structure of how ordinary human beings experience themselves, and until those roots are understood, the symptom, which is the consistent elevation of the man who suffers over the man who celebrates, will continue to appear regardless of what teaching or ideology is placed over it.

So let us begin here. The person who celebrates, the person who lives in visible ease and visible delight, creates, simply by existing in that state, a problem for everyone around him who is not in that state. And the vast majority of people are not in that state. They are burdened. They are worried. They are carrying the weight of things undone, of relationships that have soured, of ambitions that have narrowed with the years, of a body that aches in the morning, of a mind that will not rest at night. This is the ordinary condition of ordinary life, and in this condition, the person who moves through the world with obvious joy, with lightness, with a kind of physical ease that suggests the body itself is in agreement with his existence, is not simply a pleasant sight. He is, at a level far below what the watching person would consciously admit, a reproach.

He is a reproach because his very happiness raises a question that is extraordinarily uncomfortable: if this person can be happy, why cannot I? And this question has two possible answers. The first answer is: because I have not done the inner work that this person has done, because I am still carrying weights that could be set down, because my unhappiness is at least partly the result of my own choices and my own avoidances and my own refusals to look honestly at certain things about myself. This is the more courageous answer. It is also the more painful one, because it puts the responsibility squarely inside the watching person and offers no comfortable alternative narrative.

The second answer is: because this happy person has something I do not have, some advantage, some circumstance that has been given to him and not to me. He is lucky, or he is superficial, or he does not feel things as deeply as I do, or he has not yet encountered the genuine losses and sufferings that life eventually brings to everyone. And if the genuineness of his happiness cannot be denied, then there is still one more option: his happiness itself is morally suspect. He has not paid enough for it. A person who has not paid the required price has no right to be happy. And therefore his happiness, however genuine it looks, is somehow not the real thing. It is shallow. It will not last. Or, most damning of all, it is selfish.

Now the person who suffers publicly, the person who carries his sacrifice visibly and whose pain is obvious and readable and confirmed by his physical bearing, creates the opposite effect. He does not raise the uncomfortable question. He confirms the watching person's existing experience. He says, through his very posture and his very suffering, that life is indeed as difficult as the watching person has always felt it to be. He validates the difficulty. He makes the watching person feel that his own unhappiness is not a personal failure but simply an honest response to a genuinely heavy world.

This is the first and most fundamental reason people worship sacrifice. It is not admiration in the genuine sense. It is recognition. The man who suffers is recognized, and recognition is one of the deepest needs of the human being, and when it is offered, the recipient feels toward the giver a warm gratitude that is easily mistaken for reverence.

But there is a second layer, and this one cuts deeper. It has to do with guilt.

The ordinary person does not feel that he deserves to be happy. This is not something he announces. He would deny it if asked directly. But it functions in him with great consistency. He has been told, by family, by school, by religion, by the general texture of the culture he grew up inside, that happiness must be earned, that comfort must be justified, that pleasure is acceptable only after work and only in proportion to the work, that the good life is not a gift but a reward and that a reward requires sacrifice. And so, at a very deep level, a level where it is never examined because it is simply the unquestioned air that everything else breathes, the ordinary person believes that his own moments of genuine happiness are somewhat fraudulent, that he has not yet done enough to deserve them, that they could be taken from him at any moment as the consequence of not having paid enough of the required price.

From this internalized guilt, the worship of sacrifice follows directly. The man who has obviously suffered, who has given up his comfort, his pleasures, his ordinary human securities, has paid the price that the watching person secretly fears he himself has not paid. The man of sacrifice is therefore not just admired but envied, in the precise sense that he has acquired something the watching person wants: the right to be taken seriously, the right to have his inner life treated as genuine, the right to exist without the vague and nagging suspicion that he has gotten away with something.

There is a theological structure that makes this psychological dynamic extremely visible, and it is the Christian doctrine of atonement, particularly in its substitutionary form. The doctrine says that the debt of human sin was so enormous that ordinary human beings could not pay it themselves, and that therefore a being of infinite worth took the debt upon himself and paid it through his suffering and his death. This is a structure of infinite sacrifice as the solution to infinite debt. And what it does, psychologically, is not primarily what it claims to do theologically. Theologically it is supposed to liberate the believer from guilt. Psychologically what it often produces is a permanent, unrepayable obligation to the one who sacrificed, accompanied by a deepening of the sense that one's own ordinary happiness is an affront to the magnitude of the price that was paid on one's behalf.

Augustine understood this dynamic from the inside, because he lived it. Before his conversion, Augustine was, by his own account, a man who pursued pleasure with considerable intelligence and considerable success. He had a mistress whom he loved genuinely, a son whom he loved, friends with whom he argued and laughed. And then the conversion happened, and what followed was not simply a shift in metaphysical belief but a comprehensive turning against everything that had given pleasure. He began to describe his earlier life, including the genuine loves in it, as sin, as evidence of his distance from the divine. And the tradition that grew from his thinking has produced, over fifteen centuries, a spiritual culture in which the man who has suffered most is consistently placed closest to the divine, and the man who has enjoyed most is consistently placed furthest.

I am not saying Augustine was wrong. I am saying he was telling the truth about his own experience, and his own experience happened to be the experience of a person for whom the turn from pleasure to what he called God was genuinely liberating. For him, personally, the earlier life had genuinely been a form of distraction from something more real. His mistake, or rather the mistake of the tradition built on his experience, was to universalize this path. To say that because the turn from pleasure to asceticism was liberation for Augustine, it must therefore be the path for everyone. And from this universalization has come an entire civilization's deep suspicion of joy and its equally deep reverence for suffering.

Now I want to bring in a figure from a completely different tradition, because when the same pattern appears across traditions that have had no contact with each other, it is worth very careful attention. In India, the tradition of tapas, of austerity, of the deliberate subjection of the body to conditions of heat or cold or hunger for the purpose of generating spiritual clarity, is extraordinarily ancient. The great seers of the Upanishadic period sat in meditation for periods of enormous duration, in postures of physical difficulty, in conditions of genuine deprivation. And the tradition that grew from their practice has consistently believed that the more severe the austerity, the more powerful the spiritual consequence. Indra himself trembles, says the tradition, when a human being performs tapas of sufficient intensity, because the concentrated spiritual fire of genuine austerity threatens even divine authority.

This belief, that suffering concentrated and willingly undertaken generates spiritual power, is genuinely different from the Augustinian suspicion of pleasure, but it leads to the same cultural result: the person who suffers most deliberately is the person whose inner life is most respected, and the person who is comfortable and at ease is assumed to be spiritually shallow.

But the great ones, the people who actually transformed in the traditions I am describing, knew something that the traditions built around them have consistently forgotten. They knew that the austerity was a means toward a condition of inner freedom that, once reached, no longer needed the austerity. The Buddha sat for years in increasingly severe asceticism. And then he stopped. He ate a bowl of rice offered to him by a young woman, and the ascetics who had been sitting with him were appalled and left, because they had confused the method with the goal and could not recognize that he had arrived somewhere that made the method obsolete.

And what he arrived at, the inner freedom that the tradition calls nirvana, was not a state of suffering. It was a state of extraordinary peace and extraordinary clarity that included, according to every account, an enormous capacity for joy.

Alongside these dimensions there is also a layer of simple social function that has always supported the worship of sacrifice, and it operates through the mechanism that the French anthropologist Ren'e Girard spent most of his intellectual life describing. Girard observed that human communities have a persistent tendency to manage internal tension and aggression by directing it toward a single figure, whose suffering becomes the means of the community's own cohesion. The person who sacrifices, whether involuntarily as the scapegoat or voluntarily as the ascetic, becomes the container for the community's collective need to see suffering enacted. His sacrifice is the community's own aggression against life turned visible and safely located outside the ordinary members of the community.

This is not cynical. Girard was not saying the community is being deliberately cruel. He was saying that communities in tension need release, and the figure who suffers provides it. And this need, once a community develops the habit of meeting it through the worship of sacrifice, becomes self-reinforcing. The worship itself creates the role, makes the role available and legible and respectable, so that future generations find it already waiting to be filled.

The great French philosopher and mystic Simone Weil understood this from the inside. She spent her adult life systematically refusing any comfort not shared by the people she felt most solidarity with. She worked in factories. She went to the front lines of the Spanish Civil War. She died at thirty-four, in part because she refused to eat more than the rations available to French people under German occupation even when she was already seriously ill and living in England. And the response to her life has been, consistently, a reverence that her extraordinary writing has never quite achieved on its own terms. People revere her suffering more than her thought. Her thought asks something of the reader. Her suffering only asks to be acknowledged.

This is the key distinction. The man who genuinely celebrates does not, merely by celebrating, make any demand on the people who watch him. The real demand is made by the man who thinks with genuine depth, who investigates, who challenges, who requires you to reconsider something you have always assumed. The thinker who also happens to live joyfully is the most threatening of all, because his joy removes the last defense: the defense of saying that genuine engagement with life's deep questions must necessarily produce sorrow.

Nietzsche, who was himself a man of considerable suffering and therefore spoke from some authority on the subject, called this consistent elevation of suffering above joy slave morality. He said that the systematic revaluation by which renunciation becomes virtue and pleasure becomes suspect was the great psychological achievement of those who could not obtain what they wanted and therefore reframed wanting as the problem rather than the inadequacy of their own pursuit. The insight at the center of this argument is accurate: when a culture consistently elevates the denier above the celebrant simply on the grounds of denial and not on any demonstrable grounds of genuine inner depth, it has confused the sacrifice with the destination. It has taken the medicine for the disease and the disease for health.

There is a story I find useful here. In a certain village there were two teachers. The first had fasted for many years. His body was thin to the point of transparency, his eyes had sunk deep into his face, and he moved with the deliberateness of a man who had long since surrendered whatever physical ease most people move through the world with. He spoke rarely and his voice was barely audible. The villagers brought him offerings, which he refused. They built him shelter, which he accepted reluctantly. He was revered throughout the entire district. People traveled long distances to sit in silence near him.

The second teacher was a large and lively man who cooked well and ate with obvious satisfaction and laughed at things he found genuinely funny and was interested in the lives of the people who came to him. He remembered details about their children and their work and offered advice that was practical and specific and often immediately useful. He had studied deeply, traveled widely, and arrived through a long interior journey at a settled quality of presence that the perceptive visitor could feel within a few minutes of being near him. But he was thought of in the surrounding villages as a pleasant and helpful man, not a holy one.

One day a traveler who had visited both teachers said to the second: is it not unfair? You have clearly gone further in your inner work. Your understanding is more genuine and more settled. And yet the first is worshipped and you are merely appreciated. The second teacher thought for a moment and then said: the first teacher gives people what they already know how to receive. I ask them to learn a new way of receiving, and that takes time, and most people do not have the time, and I cannot blame them for it.

This is the situation precisely. The worship of sacrifice is the worship of the familiar. Suffering is the most familiar thing in the world. Every human being knows it from the inside. It requires no new perceptual capacity, no willingness to question the existing maps. It simply confirms what everyone has always already known: that life is heavy, and those who feel that heaviness most intensely are those whose inner life is most real.

Krishna disturbs people, and has always disturbed people, because he refuses to offer them this confirmation. He offers something more demanding and more real, which is the possibility of being completely and without reservation at home in what you are. Not free from the burden of experience, but free within it. Not arrived at some other world but fully arrived at this one. And the crowd, which has been taught for generations that holiness looks like sacrifice and that joy is either shallow or unearned, simply does not know where to put him. So it calls him divine and thereby makes him safely distant, because what is divine cannot be demanded of ordinary people, and what cannot be demanded cannot disturb.

The worship of sacrifice, at its most fundamental, is the worship of the absence of demand. The man who has suffered publicly has already paid. He asks nothing more of you than to be acknowledged. But the man who is fully alive, whose joy is not performance but genuine presence, asks everything: he asks you to consider that your pain is not your most essential quality, that the life available to you is larger than the life you have been living, that the gap between where you are and where you could be is not a gap of earned unworthiness but of untried possibility. And this is the most disturbing thing a human being can ask of another: not sacrifice, not obedience, but the complete willingness to be present to your own life in its full and genuinely available aliveness.

\begin{table}[htbp]
\caption{A rough map of what people usually trust. The comparison reproduces the author's categories.}
\label{tab:trust}
\centering\small
\begin{tabular}{@{}P{.20\textwidth}P{.23\textwidth}P{.23\textwidth}P{.25\textwidth}@{}}
\toprule
Figure & Recognized by & Crowd reaction & Demand placed on us \\
\midrule
The suffering saint & Cross, cave, fasting, ash & Reverence & Find meaning in pain \\
The dancing god & Flute, rhythm, celebration & Beauty mixed with suspicion & Trust joy without proof \\
The world-rejector & Distance from ordinary life & Worship from a safe distance & Admire, but stay unchanged \\
The world-transformer & Battlefield, street, city, unfinished work & Discomfort and resistance & Become different \\
\bottomrule
\end{tabular}
\end{table}
\FloatBarrier

\section{The Body Is Not Your Enemy}

\subsection{The body and impermanence}

If the body is temporary and created by God, why is hatred toward the body considered spirituality?

It would be completely wrong to think that the people who developed the theology of bodily hatred were stupid people or malicious people. They were, in many cases, extraordinarily sensitive and genuinely intelligent human beings who were wrestling with real and difficult problems about the nature of human experience. But they made one fundamental error at the very beginning of their thinking, and that error, being foundational, infected everything built upon it the way a crack in the foundation of a building affects every wall and every ceiling and every floor above it. The building may stand for a very long time. It may even be beautiful in many of its parts. But eventually the crack works its way upward and outward, and nothing built on it is entirely stable.

The foundational error was this: the observation that pleasure fades was taken as evidence that pleasure is unreal. And from this a conclusion was drawn that is not logically necessary but felt psychologically compelling to people living in genuine suffering: if bodily pleasure is temporary and therefore unreal, then the real must lie somewhere else, somewhere outside the body, somewhere beyond the senses. And whatever lies beyond the senses must be protected from the contamination of the senses. Therefore the body is an obstacle to truth. And an obstacle is, by its nature, something to be overcome, to be pushed aside, to be, if possible, eliminated. And the elimination of an obstacle to truth is an act of truth-seeking. And an act of truthseeking is, in the vocabulary of every tradition, a holy act. Therefore: hating the body is holy.

This chain of reasoning, followed step by step, is impeccable. Each step follows from the one before it. The only weak link is the very first one, the equation of temporariness with unreality. And this link is not merely weak; it is entirely unsupported. But because it sits at the beginning of the chain, before anyone has stopped to examine it, before it has been subjected to the kind of scrutiny that later steps in an argument receive, it gets smuggled in as an axiom. It becomes something that is simply assumed rather than something that has been established. And once it is established as an axiom, everything built on it feels as solid as the axiom itself, even though the entire structure is resting on a foundation that was never actually tested.

Let us test it. Let us ask directly: does the fact that something is temporary make it unreal?

Consider a conversation between two people who have known each other for a long time and have not spoken in many years. They meet by accident in a city neither of them lives in. They sit down together and talk for three hours. And in those three hours something happens between them that is entirely real: the recognition of shared history, the specific texture of a friendship that has its own accumulated depth, the particular way that two people who have known each other at formative moments of their lives speak to each other. This conversation lasts three hours. It ends. It will never happen again in exactly this form. The two people will eventually die. Does the temporariness of the conversation make it unreal? Does the fact that it ended mean that it did not happen? Does the fact that both participants will eventually be gone mean that the understanding that passed between them was an illusion?

The argument from impermanence to unreality requires you to say yes to these questions. It requires you to say that the conversation did not matter, that the understanding was illusory, that the beauty of the recognition was a deception, because none of it was permanent. And this is precisely the conclusion that every genuine human experience refuses to confirm. The reality of an experience does not depend on its duration. The significance of a moment does not depend on whether the moment repeats itself forever. The beauty of a thing does not depend on whether it lasts.

Now this is not merely a philosophical observation. It is also a deeply practical one, with immediate consequences for how one lives. Because if you accept the equation of temporariness with unreality, you become unable to be fully present to any actual experience, since all actual experiences are temporary. You are perpetually reaching past what is actually happening toward a permanent something that is never actually here. Your life becomes a constant movement away from where you are toward where you imagine the real to be. And this is precisely the condition that all the teachers of world-rejection have described as the human condition: the condition of a being who is always somewhere else, always in the past or the future, never fully here, never fully present to what is actually in front of it.

But they have made this condition worse by offering, as the cure, an even more radical flight from the present. They say: you are not here because you are attached to the body and its pleasures. Stop being attached to the body. Transcend the body. Move beyond the senses. And then you will be free. But this is asking a man who is already not here to become even less here, in the hope that by achieving some ultimate degree of not-here-ness he will paradoxically arrive at the real, which they have located in a beyond that is by definition the farthest possible point from here.

The genuine insight, the insight that emerges not from philosophical reasoning about permanence and impermanence but from the actual direct experience of a human being who has genuinely arrived at their own existence, is the exact opposite of this. It is that the real is accessed not by moving away from the present and the sensory and the physical but by moving so completely into the present, so fully into the actual texture of what is happening right now in this body in this moment, that the artificial separation between the one who is experiencing and the experience itself dissolves, and what is left is something that has no name in the ordinary vocabulary but that every mystic in every tradition has tried to point at in their own way.

And in this genuine arrival at the present, the body is not an obstacle. The body is the door. The body is the very instrument through which presence becomes possible. Without the body, there is no breath to observe, no heartbeat to feel, no sensation of sunlight on skin or wind across the face or the specific weight of another person's presence in the same room. Without the body, the whole dimension of immediate, sensory, particular existence, which is the dimension in which the present moment actually lives, is simply not available.

Thomas Aquinas said the soul is imprisoned in the body. And I understand what he meant, because he was speaking from within a framework in which the soul was assumed to be the real thing and the body was assumed to be the lesser thing. But consider the possibility that the body is not a prison for the soul but a home for it. That the body is not what keeps the soul from being free but what gives the soul the particular, specific, located experience of being alive that freedom requires as its ground. Because freedom that has no body to move through is not freedom; it is simple absence. And simple absence, the absence of weight and warmth and sensation and particular location, is not the experience that any mystic in the genuine tradition has ever described as liberation.

The mystic does not describe liberation as the experience of not having a body. The mystic describes it as the experience of having a body without being imprisoned by fear of the body. Without being imprisoned by the thought that the body will eventually fail and die and that this failure is a catastrophe. Without being imprisoned by the shame of the body's desires and the guilt that the tradition has layered over those desires. The mystic's liberation is the liberation of a full human being, body and all, from the fear and shame and guilt that have made the body feel like a prison. It is not the liberation of a disembodied consciousness from the inconvenience of having been temporarily housed in flesh.

Now, there is a very important observation to be made here about the history of the traditions that have promoted bodily hatred, and the observation is this: every single tradition that has most vigorously promoted hatred of the body has also, alongside the official doctrine, maintained an underground tradition that said exactly the opposite. And this underground tradition was not a corruption of the official doctrine. It was the truth that the official doctrine was suppressing.

Within Christianity, the official doctrine of the body as a site of sin and the source of temptation coexisted with the mystical tradition of figures like Julian of Norwich, who spoke of the body of Christ as something utterly beautiful and as the very vehicle through which the divine love makes itself known in the world. Or Meister Eckhart, who spoke of the birth of God in the soul as an immediate and embodied experience, not a future post-death reality. Or the Song of Solomon, which sits in the middle of the Hebrew Bible with its extraordinarily frank and beautiful celebration of physical love and desire, which every interpreter has had to work very hard to allegorize into spiritual meaning because the literal meaning is simply too inconveniently alive.

Within Indian philosophy, the official doctrine of bodily denial and the path through austerity to liberation coexisted with the Tantric tradition, with the erotic temple sculptures of Khajuraho, with the entire devotional tradition of bhakti yoga, which used the body's capacity for love and desire as the primary vehicle of the movement toward the divine. The linga and yoni worshipped in every Shiva temple are not despite the official doctrine of transcendence; they are in direct and permanent tension with it, representing a truth that the official doctrine could suppress in its verbal teachings but could never quite remove from the living body of the culture.

This persistent, recurring, universal appearance of the counter-tradition, the tradition that says the body is not the enemy but the vehicle, is not coincidence. It is evidence that the official doctrine of bodily hatred is, at some fundamental level, contradicted by the actual experience of the people it is supposed to be teaching. Because people live in bodies. And in their actual experience of living in bodies, they encounter, alongside the suffering that the tradition correctly identifies, also an encounter with dimensions of experience that feel utterly genuine, utterly real, and utterly sacred, even though they are completely bodily in their nature. The taste of good food eaten with full presence. The warmth of physical human contact. The extraordinary state of consciousness that music can induce when the body is completely given over to the rhythm and the sound. The particular quality of aliveness that comes from sustained physical exertion. All of these are bodily experiences. And none of them feel, in the actual moment of their occurrence, like obstacles to truth. They feel, if anything, like glimpses of it.

And so the honest human being, the human being who trusts their own experience over the doctrine they have been handed, has always found themselves in a private and uncomfortable disagreement with the official theology of the body. They have had to hide from themselves, and from the authorities of their tradition, the simple and persistent fact that their body, rather than being an obstacle to their inner life, has often been the very instrument through which their inner life has been richest and most alive.

The body was not created by God, if we are to use that language at all, in order to be hated. It was created, or it evolved, or it emerged from the same process that produced everything else, in order to be the living interface between the awareness that knows and the world that is known. Hating it is as sensible as hating your eyes for allowing you to see things that are sometimes painful to look at. The solution to the pain of seeing is not to destroy the eyes. It is to develop the wisdom to look with understanding rather than with fear. And wisdom of that kind cannot come from above the body. It can only come through it.

\subsection{Sexuality and moral authority}

If existence itself created sexuality, why has man treated sexual energy like a crime? Is humanity fighting nature in the name of morality?

Yes. And the fighting has been going on for so long, and has been so thorough, and has involved so much sincere effort from so many genuinely intelligent and genuinely well-intentioned people, that it has acquired the appearance of wisdom. It has been going on so long that it no longer looks like a war. It looks like the natural order of things. It looks like the way the world is rather than a very particular and very consequential set of decisions that were made at specific moments in human history for specific reasons that had very little to do with spiritual wisdom and a great deal to do with social organization and the management of populations.

Let us look at this very carefully, because this is not a simple matter and does not yield to simple responses.

Existence produces sexual energy. This is the first fact, and it is a fact of such enormous weight and importance that everything else in this discussion depends on it being held clearly in mind. Existence produces sexual energy. Not humanity. Not society. Not the devil. Existence. The same force or intelligence or accident that produced the universe, that produced the stars and the oceans and the first single-celled organisms trembling at the edge of life and death four billion years ago, that same force produced sexual energy as one of its primary expressions. It is the energy of life reaching toward more life. It is the fundamental creative impulse of the cosmos expressed through biological forms. The tree reaching its roots toward water and its branches toward light, the river finding its way through stone and soil to the sea, the seeds bursting through the cold ground toward whatever warmth spring offers, all of these are expressions of the same fundamental energy that in human beings takes the specific form of sexual longing and sexual joy.

This is not a metaphor. It is a biological and philosophical fact. The biologists tell us that sexual reproduction is one of the most extraordinary developments in the history of life on earth, because it introduced something that asexual reproduction does not have: the possibility of genuine novelty, of offspring that are not merely copies of their parents but genuinely new combinations of possibility. Sexuality, in this understanding, is the universe's mechanism for creating something new, for preventing life from simply repeating itself forever in exactly the same forms. It is the engine of evolutionary creativity. Without it, the biological world would be static. Without it, there would be no diversity, no adaptation, no capacity to meet new challenges with new responses. There would be no you. There would be nothing resembling the human mind or the human heart.

Now, given all of this, given the fundamental and creative and utterly indispensable role that this energy plays in the structure of living existence, the question your question asks is not merely interesting but staggering: why has the culture built around the very beings produced by this energy treated the energy itself as something shameful, something criminal, something to be fought and suppressed and denied at every possible turn?

The standard answer that religious traditions have offered is that sexual energy is tied to the ego, to the attachment to the individual self, to the process of perpetuating the individual self through offspring, and that the spiritual life requires a dissolution of the individual self and therefore requires a dissolution of the energy that perpetuates it. This answer has a certain philosophical logic. I am not going to dismiss it entirely. But it is fatally incomplete, and its incompleteness has produced catastrophic consequences in the actual lives of actual human beings.

The incompleteness lies in the following assumption: that dissolution of the ego is achieved by fighting the forces that are associated with the ego. But this assumption is wrong. Utterly wrong. And we have known it to be wrong, or should have known it, since Freud made his great discovery at the turn of the twentieth century. Freud's central insight, his most important and most enduring contribution to the understanding of the human mind, is this: what you push underground does not go away. It transforms. It reappears in forms that are more powerful and more distorted and more difficult to work with than the original energy in its natural form.

This is not a merely theoretical observation. It is an observation derived from Freud's clinical practice, from decades of sitting with actual human beings and listening to the actual contents of their inner lives. And what he found, over and over again, in case after case, was that the people who had most thoroughly suppressed their sexual energy were the people who were most thoroughly controlled by it. The man who had spent twenty years in rigorous religious practice designed to eliminate all sexual thought was the man whose sexual obsession had simply gone underground, where it was no longer available for conscious examination or conscious direction, but was operating through disguised channels: through excessive religiosity that had an unmistakably erotic character, through displaced desire that attached itself to religious imagery or to relationships with spiritual authority figures, through sudden and violent eruptions of the suppressed material that would break through the crust of control at unexpected moments and in deeply damaging ways.

This pattern, which Freud documented in the context of Victorian European civilization, was not new with Freud. It has been visible in every culture that has practiced aggressive sexual suppression in the name of religion. The history of institutional religion is full of the consequences of this suppression: the sexual abuse that occurs within celibate religious institutions, the obsessive preoccupation with sexual sin that characterizes the most sexually repressive religious cultures, the explosion of sexuality in taboo and commercially organized forms in exactly those societies where the official morality is most rigidly controlling. These are not accidents. They are the predictable consequences of fighting a fundamental life-force rather than understanding and integrating it.

Now let us ask the question that almost never gets asked in these discussions: why did humanity decide to fight this particular energy? What was the actual historical and social logic of this decision?

The answer has several layers, and none of them are spiritual. The first layer is economic. In agrarian societies, which is the context in which most of the world's major religions developed their core teachings, the inheritance of land and property required the certainty of biological paternity. A man needed to know that the sons to whom he was leaving his property were actually his biological offspring. This required control over female sexuality. The control of female sexuality, once established as a social necessity, was then extended, for the sake of internal consistency and social cohesion, to male sexuality as well. The entire edifice of sexual morality that these religions built, with its enormous weight of guilt and shame and sin, began as a property law. This is not a pleasant thing to say about the traditions that have shaped the inner lives of billions of human beings across thousands of years. But it is true.

The second layer is political. Organized religion, in virtually every case in which it has acquired significant social power, has recognized that sexual energy is one of the primary forces that creates strong bonds between individuals that are independent of and potentially competitive with the bonds that the institution wants to create between individuals and itself. A man who is in a deep and complete sexual relationship with another human being has a center of gravity that is not the institution. He has someone he turns to in the night, someone whose opinion matters more than the priest's opinion, someone through whom he experiences something of the same dissolution of boundaries and the same access to a dimension of experience that transcends the ordinary that the institution claims to offer exclusively through its own rituals and sacraments. The institution correctly perceives this as competition. And the way it has dealt with this competition is to make sexual experience shameful, to insist that the pleasure it offers is spiritually contaminating, to require that it be confined to the most tightly controlled possible set of circumstances, ideally circumstances that the institution itself oversees and sanctifies.

The third layer is psychological, and it is perhaps the most subtle. The people who built these systems of sexual suppression were not, in the main, cynical manipulators. They were human beings who had, in many cases, genuinely experienced something of the problematic character of ungoverned sexual desire. They had seen the suffering that comes from desire without wisdom: the devastation of unwanted pregnancy in contexts without adequate support, the suffering caused by jealousy and possessiveness and betrayal, the way that obsessive desire can narrow a person's entire field of consciousness until there is nothing left in it but the object of desire. These are real problems. Desire without wisdom is genuinely painful and genuinely damaging. And the people who built the traditions of sexual suppression were often genuinely trying to address this pain and this damage.

The mistake they made was the mistake of the sledgehammer used where a scalpel was needed. The problem was not the energy but the absence of wisdom in relation to the energy. The solution was not to eliminate the energy but to understand it, to develop wisdom in relation to it, to find ways of engaging with it consciously rather than being driven by it unconsciously. But this requires a much more difficult and much longer and much more personally demanding process of education and inner development. It is far simpler to say: the energy is bad, suppress it. And simple things are what institutions prefer, because simple things can be organized. Complex things are too individual, too dependent on the particular nature and situation of each person, to be managed by a large-scale institutional structure.

And so humanity chose the simple solution. And the simple solution has failed for five thousand years, in every culture that has tried it, with a consistency that ought to long ago have prompted a fundamental re-examination of the premise. But the premise has not been re-examined, because the premise is now so deeply embedded in the culture that to question it is to be seen as endorsing the very excess and suffering it was designed to prevent. The argument against questioning the suppression is always: look at what happens when there is no suppression, look at the chaos and the suffering of uncontrolled desire. And this argument has force, because uncontrolled desire does produce chaos and suffering. But it is the argument of someone who has never considered the third possibility, the possibility that lies between suppression and chaos, the possibility of desire met with full understanding and genuine freedom.

Tantra discovered this possibility and was immediately and systematically persecuted by every established religious authority that encountered it. Because Tantra said something that the institutions could not tolerate: that the sexual energy is not the enemy of the sacred but one of its most powerful expressions, that this energy, when it is met with full awareness rather than with fear, becomes one of the most direct available paths to the experience of the dissolution of the individual ego in something larger. It said that the problem is not the energy but the unconsciousness with which it is ordinarily engaged. And it developed, over centuries, extraordinarily refined methods for engaging with this energy in ways that used it for the expansion of consciousness rather than the perpetuation of bondage.

The institutional religions called this perversion. They called it an excuse for indulgence. They could not see, or would not see, that Tantra was doing something far more demanding than either suppression or indulgence. It was doing what all genuine spiritual work does: it was taking the actual material of human life, exactly as it is, and working with it consciously, transforming it from the inside rather than trying to eliminate it from the outside.

Is humanity fighting nature in the name of morality? Absolutely. And the fighting is visible everywhere, in the distorted forms that the suppressed energy takes when it finally breaks through, in the obsessive preoccupation with sexual sin that characterizes the most sexually repressive cultures, in the profound confusion about the relationship between love and the body that is the central psychological wound of the Western tradition. And the damage is not only personal. It is civilizational. A civilization built on the suppression of one of its most fundamental creative energies is a civilization that does not know itself, that is at war with its own roots, that pours an enormous proportion of its available energy into a self-destructive conflict that it cannot win and cannot afford to stop fighting.

The alternative is not permissiveness. I want to be very clear about this. The alternative is not the elimination of all guidance in relation to this energy and the encouragement of simple indulgence. That produces its own form of suffering, as Freud's successors in the culture of the sexual revolution of the twentieth century eventually had to acknowledge. The alternative is wisdom. The alternative is the development of a genuine understanding of what this energy is and what it is seeking, an understanding that takes it seriously rather than dismissing it either by suppressing it or by reducing it to mere recreation. The alternative is what India knew at its most lucid moments and then forgot, or rather was persuaded by its own institutional religions to pretend to forget, which is that this energy, this fundamental creative energy of existence, is in its deepest nature not different from the divine. That it is, in fact, one of the most direct faces that the divine turns toward human beings. And that to fight it is to fight the divine itself, while imagining that you are serving it.

\par\medskip\begingroup\footnotesize\raggedright
\noindent\textbf{Image credit.} Caspar David Friedrich, The Monk by the Sea. Alte Nationalgalerie, Berlin. Public-domain reproduction via Wikimedia Commons.

Image source: \url{https://commons.wikimedia.org/?curid=2221334}.
\par\endgroup

\end{document}
